\documentclass[11pt,a4paper]{scrartcl}
\usepackage[margin=22mm,headheight=18pt]{geometry}
\usepackage{fontspec}
\usepackage{xeCJK}
\IfFontExistsTF{Songti SC}{\setCJKmainfont{Songti SC}}{\setCJKmainfont{FandolSong-Regular.otf}[BoldFont=FandolSong-Bold.otf]}
\IfFontExistsTF{Heiti SC}{\setCJKsansfont{Heiti SC}}{\setCJKsansfont{FandolHei-Regular.otf}[BoldFont=FandolHei-Bold.otf]}
\usepackage{amsmath,amssymb,mathtools,mathrsfs,bm,cancel}
\usepackage{mathpazo}
\setmainfont{texgyrepagella-regular.otf}[BoldFont=texgyrepagella-bold.otf,ItalicFont=texgyrepagella-italic.otf,BoldItalicFont=texgyrepagella-bolditalic.otf]
\usepackage[normalem]{ulem}
\usepackage[svgnames,table]{xcolor}
\usepackage[most]{tcolorbox}
\usepackage{tikz}
\usetikzlibrary{arrows.meta,positioning,calc,matrix,patterns,decorations.markings,shapes.geometric}
\usepackage{enumitem,booktabs,longtable,array,tabularx,multicol}
\usepackage{microtype,fancyhdr}
\definecolor{ChapterBlue}{RGB}{20,60,120}
\definecolor{NoteRed}{RGB}{179,0,0}
\definecolor{AnnotationBlue}{RGB}{30,65,135}
\setkomafont{section}{\Large\bfseries\color{ChapterBlue}}
\setkomafont{subsection}{\large\bfseries\color{ChapterBlue}}
\setlist{nosep,leftmargin=*,topsep=3pt}
\setlength{\parindent}{0pt}
\setlength{\parskip}{5pt plus 1pt}
\linespread{1.12}
\allowdisplaybreaks
\emergencystretch=3em
\widowpenalty=10000
\clubpenalty=10000
\pagestyle{fancy}
\fancyhf{}
\fancyhead[C]{\small\nouppercase{\leftmark}}
\fancyfoot[C]{\thepage}
\renewcommand{\headrulewidth}{0.3pt}
% Original scan boundaries are internal anchors only; no visible page labels.
\newcommand{\sourcepage}[1]{\par\phantomsection\label{source:#1}}
\newif\ifchapterstart
\newcommand{\noteschapter}[1]{\clearpage\section{#1}\global\chapterstarttrue}
\newcommand{\notessection}[1]{\ifchapterstart\global\chapterstartfalse\else\clearpage\fi\subsection{#1}}
\renewcommand{\sectionmark}[1]{\markboth{}{}}
\renewcommand{\subsectionmark}[1]{\markboth{\thesubsection\quad #1}{}}
\renewcommand{\contentsname}{目录}
\setcounter{tocdepth}{2}
\setcounter{secnumdepth}{2}
\newcommand{\unclear}[2]{\textnormal{\textcolor{NoteRed}{［无法辨认 #1］}}}
\newcommand{\ann}[1]{{\color{AnnotationBlue}#1}}
\newcommand{\redann}[1]{{\color{NoteRed}#1}}
\newcommand{\circled}[1]{\mbox{\textcircled{\scriptsize #1}}}
\newcommand{\R}{\mathbb{R}}
\newcommand{\N}{\mathbb{N}}
\newcommand{\Z}{\mathbb{Z}}
\newcommand{\Q}{\mathbb{Q}}
\newcommand{\C}{\mathbb{C}}
\newcommand{\F}{\mathbb{F}}
\newcommand{\V}{\mathcal{V}}
\newcommand{\E}{\mathcal{E}}
\newcommand{\abs}[1]{\left|#1\right|}
\newcommand{\norm}[1]{\left\|#1\right\|}
\newcommand{\inner}[2]{\langle #1,#2\rangle}
\newcommand{\floor}[1]{\lfloor #1\rfloor}
\newcommand{\ceil}[1]{\lceil #1\rceil}
\DeclareMathOperator{\rank}{rank}
\DeclareMathOperator{\tr}{tr}
\DeclareMathOperator{\Span}{span}
\DeclareMathOperator{\diag}{diag}
\DeclareMathOperator{\im}{Im}
\DeclareMathOperator{\id}{id}
\DeclareMathOperator{\Hom}{Hom}
\DeclareMathOperator{\End}{End}
\DeclareMathOperator{\Ker}{Ker}
\DeclareMathOperator{\sgn}{sgn}
\newtcolorbox{noteblock}[1]{enhanced,breakable,colback=ChapterBlue!2,colframe=ChapterBlue!60,title={#1},fonttitle=\bfseries,boxrule=0.4pt,arc=1mm}
\newenvironment{proof}{\par\textbf{Proof.}\ }{\par}
\tikzset{vertex/.style={circle,draw,inner sep=2pt},edge/.style={thick},dedge/.style={thick,-Stealth}}

% Clickable contents and PDF outline.
\usepackage[unicode,hidelinks,bookmarksnumbered,bookmarksopen,linktoc=all]{hyperref}
\hypersetup{pdftitle={实变函数},pdfsubject={完整转录，含可点击目录},pdfauthor={芬达不带气儿}}

\newcommand{\eps}{\varepsilon}
\DeclareMathOperator{\supp}{supp}
\newcommand{\BV}{\mathrm{B.V.}}

\pdfstringdefDisableCommands{\def\infty{∞}\def\R{ℝ}}

% Coordinate system: x and y are per-mille of original 1415x2000 preview image.
% Approximate each original baseline as (pixel_x/1.415,pixel_y/2).
% Width is also in per-mille of the original image width.
\newlength{\handunit}\setlength{\handunit}{0.194mm}
\newlength{\handyunit}\setlength{\handyunit}{0.265mm}
\newsavebox{\handbox}
\newcommand{\handtext}[6]{%
  \node[anchor=base west,inner sep=0pt,outer sep=0pt] at (#1,#2)
    {{\normalfont\fontsize{#4}{#4}\selectfont\color{#5}%
      \sbox{\handbox}{#6}%
      \ifdim\wd\handbox>#3\handunit\resizebox{#3\handunit}{!}{\usebox{\handbox}}\else\usebox{\handbox}\fi}};}
\newcommand{\handline}[5][black]{\handtext{#2}{#3}{#4}{10.5}{#1}{#5}}
\newcommand{\handtitle}[4]{\handtext{#1}{#2}{#3}{13}{black}{\bfseries #4}}
\newcommand{\handnote}[5][AnnotationBlue]{\handtext{#2}{#3}{#4}{7.2}{#1}{#5}}
\newcommand{\handarrow}[5][AnnotationBlue]{\draw[->,color=#1,line width=.35pt] (#2,#3)--(#4,#5);}
\newcommand{\handdiagram}[5]{\node[anchor=north west,inner sep=0pt] at (#1,#2)
 {\resizebox{#3\handunit}{#4\handyunit}{\begin{tikzpicture}[x=1cm,y=1cm]#5\end{tikzpicture}}};}
\newcommand{\handchapterentry}[1]{\refstepcounter{section}\setcounter{subsection}{0}\addcontentsline{toc}{section}{\protect\numberline{\thesection}#1}}
\newcommand{\handsectionentry}[1]{\refstepcounter{subsection}\addcontentsline{toc}{subsection}{\protect\numberline{\thesubsection}#1}}
\newcommand{\handregister}[1]{\ifcsname handentry@#1\endcsname\csname handentry@#1\endcsname\fi}
\newenvironment{handpage}[1]{\clearpage\thispagestyle{fancy}\markboth{\csname handheader@#1\endcsname}{}\null\handregister{#1}%
 \begin{tikzpicture}[remember picture,overlay]
 \begin{scope}[shift={([xshift=8mm,yshift=-13mm]current page.north west)},x=0.194mm,y=-0.265mm]
}{\end{scope}\end{tikzpicture}\par}

\fancyhead[C]{\sbox{\handbox}{\small\nouppercase{\leftmark}}\ifdim\wd\handbox>\textwidth\resizebox{\textwidth}{!}{\usebox{\handbox}}\else\usebox{\handbox}\fi}
\expandafter\def\csname handheader@1\endcsname{1.1\quad Functions of Bounded Variation}
\expandafter\def\csname handentry@1\endcsname{\setcounter{section}{0}\handchapterentry{Fuctions of Bounded Variation and the Riemann-Stieltjes Integral}\setcounter{section}{1}\setcounter{subsection}{0}\handsectionentry{Functions of Bounded Variation}}
\expandafter\def\csname handheader@2\endcsname{1.1\quad Functions of Bounded Variation}
\expandafter\def\csname handentry@2\endcsname{}
\expandafter\def\csname handheader@3\endcsname{1.1\quad Functions of Bounded Variation}
\expandafter\def\csname handentry@3\endcsname{}
\expandafter\def\csname handheader@4\endcsname{1.1\quad Functions of Bounded Variation}
\expandafter\def\csname handentry@4\endcsname{}
\expandafter\def\csname handheader@5\endcsname{1.2\quad Rectifiable Curves}
\expandafter\def\csname handentry@5\endcsname{\setcounter{section}{1}\setcounter{subsection}{1}\handsectionentry{Rectifiable Curves}}
\expandafter\def\csname handheader@6\endcsname{1.3\quad The Riemann-Stieltjes Integral}
\expandafter\def\csname handentry@6\endcsname{\setcounter{section}{1}\setcounter{subsection}{2}\handsectionentry{The Riemann-Stieltjes Integral}}
\expandafter\def\csname handheader@7\endcsname{1.3\quad The Riemann-Stieltjes Integral}
\expandafter\def\csname handentry@7\endcsname{}
\expandafter\def\csname handheader@8\endcsname{1.3\quad The Riemann-Stieltjes Integral}
\expandafter\def\csname handentry@8\endcsname{}
\expandafter\def\csname handheader@9\endcsname{1.3\quad The Riemann-Stieltjes Integral}
\expandafter\def\csname handentry@9\endcsname{}
\expandafter\def\csname handheader@10\endcsname{2.1\quad Lebesgue Outer Measure and the Cantor Set}
\expandafter\def\csname handentry@10\endcsname{\setcounter{section}{1}\handchapterentry{Lebesgue Measure and Outer Measure}\setcounter{section}{2}\setcounter{subsection}{0}\handsectionentry{Lebesgue Outer Measure and the Cantor Set}}
\expandafter\def\csname handheader@11\endcsname{2.1\quad Lebesgue Outer Measure and the Cantor Set}
\expandafter\def\csname handentry@11\endcsname{}
\expandafter\def\csname handheader@12\endcsname{2.1\quad Lebesgue Outer Measure and the Cantor Set}
\expandafter\def\csname handentry@12\endcsname{}
\expandafter\def\csname handheader@13\endcsname{2.2\quad Lebesgue Measurable Sets}
\expandafter\def\csname handentry@13\endcsname{\setcounter{section}{2}\setcounter{subsection}{1}\handsectionentry{Lebesgue Measurable Sets}}
\expandafter\def\csname handheader@14\endcsname{2.2\quad Lebesgue Measurable Sets}
\expandafter\def\csname handentry@14\endcsname{}
\expandafter\def\csname handheader@15\endcsname{2.2\quad Lebesgue Measurable Sets}
\expandafter\def\csname handentry@15\endcsname{}
\expandafter\def\csname handheader@16\endcsname{2.3\quad Two properties of Lebesgue Measure}
\expandafter\def\csname handentry@16\endcsname{\setcounter{section}{2}\setcounter{subsection}{2}\handsectionentry{Two properties of Lebesgue Measure}}
\expandafter\def\csname handheader@17\endcsname{2.4\quad Characterizations of Measurability}
\expandafter\def\csname handentry@17\endcsname{\setcounter{section}{2}\setcounter{subsection}{3}\handsectionentry{Characterizations of Measurability}}
\expandafter\def\csname handheader@18\endcsname{2.4\quad Characterizations of Measurability}
\expandafter\def\csname handentry@18\endcsname{}
\expandafter\def\csname handheader@19\endcsname{2.4\quad Characterizations of Measurability}
\expandafter\def\csname handentry@19\endcsname{}
\expandafter\def\csname handheader@20\endcsname{2.5\quad Lipschitz Transformation of $\R^n$}
\expandafter\def\csname handentry@20\endcsname{\setcounter{section}{2}\setcounter{subsection}{4}\handsectionentry{Lipschitz Transformation of $\R^n$}}
\expandafter\def\csname handheader@21\endcsname{2.6\quad A Nonmeasurable Set}
\expandafter\def\csname handentry@21\endcsname{\setcounter{section}{2}\setcounter{subsection}{5}\handsectionentry{A Nonmeasurable Set}}
\expandafter\def\csname handheader@22\endcsname{2.6\quad A Nonmeasurable Set}
\expandafter\def\csname handentry@22\endcsname{}
\expandafter\def\csname handheader@23\endcsname{3.1\quad Elementary Properties of Measurable Functions}
\expandafter\def\csname handentry@23\endcsname{\setcounter{section}{2}\handchapterentry{Lebesgue Measurable Functions}\setcounter{section}{3}\setcounter{subsection}{0}\handsectionentry{Elementary Properties of Measurable Functions}}
\expandafter\def\csname handheader@24\endcsname{3.1\quad Elementary Properties of Measurable Functions}
\expandafter\def\csname handentry@24\endcsname{}
\expandafter\def\csname handheader@25\endcsname{3.2\quad Semicontinuous Function}
\expandafter\def\csname handentry@25\endcsname{\setcounter{section}{3}\setcounter{subsection}{1}\handsectionentry{Semicontinuous Function}}
\expandafter\def\csname handheader@26\endcsname{3.3\quad Properties of Measurable functions and Theorems of Egorov and Lusin}
\expandafter\def\csname handentry@26\endcsname{\setcounter{section}{3}\setcounter{subsection}{2}\handsectionentry{Properties of Measurable functions and Theorems of Egorov and Lusin}}
\expandafter\def\csname handheader@27\endcsname{3.3\quad Properties of Measurable functions and Theorems of Egorov and Lusin}
\expandafter\def\csname handentry@27\endcsname{}
\expandafter\def\csname handheader@28\endcsname{3.4\quad Convergence in Measure}
\expandafter\def\csname handentry@28\endcsname{\setcounter{section}{3}\setcounter{subsection}{3}\handsectionentry{Convergence in Measure}}
\expandafter\def\csname handheader@29\endcsname{3.4\quad Convergence in Measure}
\expandafter\def\csname handentry@29\endcsname{}
\expandafter\def\csname handheader@30\endcsname{4.1\quad Definition of the Integral of a Nonnegative Function}
\expandafter\def\csname handentry@30\endcsname{\setcounter{section}{3}\handchapterentry{The Lebesgue Integral}\setcounter{section}{4}\setcounter{subsection}{0}\handsectionentry{Definition of the Integral of a Nonnegative Function}}
\expandafter\def\csname handheader@31\endcsname{4.1\quad Definition of the Integral of a Nonnegative Function}
\expandafter\def\csname handentry@31\endcsname{}
\expandafter\def\csname handheader@32\endcsname{4.2\quad Properties of Integral}
\expandafter\def\csname handentry@32\endcsname{\setcounter{section}{4}\setcounter{subsection}{1}\handsectionentry{Properties of Integral}}
\expandafter\def\csname handheader@33\endcsname{4.2\quad Properties of Integral}
\expandafter\def\csname handentry@33\endcsname{}
\expandafter\def\csname handheader@34\endcsname{4.2\quad Properties of Integral}
\expandafter\def\csname handentry@34\endcsname{}
\expandafter\def\csname handheader@35\endcsname{4.2\quad Properties of Integral}
\expandafter\def\csname handentry@35\endcsname{}
\expandafter\def\csname handheader@36\endcsname{4.3\quad The Integral of Arbitrary Measurable Functions}
\expandafter\def\csname handentry@36\endcsname{\setcounter{section}{4}\setcounter{subsection}{2}\handsectionentry{The Integral of Arbitrary Measurable Functions}}
\expandafter\def\csname handheader@37\endcsname{4.3\quad The Integral of Arbitrary Measurable Functions}
\expandafter\def\csname handentry@37\endcsname{}
\expandafter\def\csname handheader@38\endcsname{4.3\quad The Integral of Arbitrary Measurable Functions}
\expandafter\def\csname handentry@38\endcsname{}
\expandafter\def\csname handheader@39\endcsname{4.3\quad The Integral of Arbitrary Measurable Functions}
\expandafter\def\csname handentry@39\endcsname{}
\expandafter\def\csname handheader@40\endcsname{4.4\quad Relation between Riemann-Stieltjes and Lebesgue Integrals and the $L^p$ Spaces, $0<p<\infty$}
\expandafter\def\csname handentry@40\endcsname{\setcounter{section}{4}\setcounter{subsection}{3}\handsectionentry{Relation between Riemann-Stieltjes and Lebesgue Integrals and the $L^p$ Spaces, $0<p<\infty$}}
\expandafter\def\csname handheader@41\endcsname{4.4\quad Relation between Riemann-Stieltjes and Lebesgue Integrals and the $L^p$ Spaces, $0<p<\infty$}
\expandafter\def\csname handentry@41\endcsname{}
\expandafter\def\csname handheader@42\endcsname{4.4\quad Relation between Riemann-Stieltjes and Lebesgue Integrals and the $L^p$ Spaces, $0<p<\infty$}
\expandafter\def\csname handentry@42\endcsname{}
\expandafter\def\csname handheader@43\endcsname{4.4\quad Relation between Riemann-Stieltjes and Lebesgue Integrals and the $L^p$ Spaces, $0<p<\infty$}
\expandafter\def\csname handentry@43\endcsname{}
\expandafter\def\csname handheader@44\endcsname{4.4\quad Relation between Riemann-Stieltjes and Lebesgue Integrals and the $L^p$ Spaces, $0<p<\infty$}
\expandafter\def\csname handentry@44\endcsname{}
\expandafter\def\csname handheader@45\endcsname{4.5\quad Riemann and Lebesgue Integrals}
\expandafter\def\csname handentry@45\endcsname{\setcounter{section}{4}\setcounter{subsection}{4}\handsectionentry{Riemann and Lebesgue Integrals}}
\expandafter\def\csname handheader@46\endcsname{4.5\quad Riemann and Lebesgue Integrals}
\expandafter\def\csname handentry@46\endcsname{}
\expandafter\def\csname handheader@47\endcsname{5.1\quad Fubini's Theorem}
\expandafter\def\csname handentry@47\endcsname{\setcounter{section}{4}\handchapterentry{Repeated Integration}\setcounter{section}{5}\setcounter{subsection}{0}\handsectionentry{Fubini's Theorem}}
\expandafter\def\csname handheader@48\endcsname{5.1\quad Fubini's Theorem}
\expandafter\def\csname handentry@48\endcsname{}
\expandafter\def\csname handheader@49\endcsname{5.1\quad Fubini's Theorem}
\expandafter\def\csname handentry@49\endcsname{}
\expandafter\def\csname handheader@50\endcsname{5.1\quad Fubini's Theorem}
\expandafter\def\csname handentry@50\endcsname{}
\expandafter\def\csname handheader@51\endcsname{5.1\quad Fubini's Theorem}
\expandafter\def\csname handentry@51\endcsname{}
\expandafter\def\csname handheader@52\endcsname{5.2\quad Tonelli's Theorem}
\expandafter\def\csname handentry@52\endcsname{\setcounter{section}{5}\setcounter{subsection}{1}\handsectionentry{Tonelli's Theorem}}
\expandafter\def\csname handheader@53\endcsname{5.3\quad Applications of Fubini's Theorem}
\expandafter\def\csname handentry@53\endcsname{\setcounter{section}{5}\setcounter{subsection}{2}\handsectionentry{Applications of Fubini's Theorem}}
\expandafter\def\csname handheader@54\endcsname{5.3\quad Applications of Fubini's Theorem}
\expandafter\def\csname handentry@54\endcsname{}
\expandafter\def\csname handheader@55\endcsname{5.3\quad Applications of Fubini's Theorem}
\expandafter\def\csname handentry@55\endcsname{}
\expandafter\def\csname handheader@56\endcsname{6.1\quad The Indefinite Integral}
\expandafter\def\csname handentry@56\endcsname{\setcounter{section}{5}\handchapterentry{Differentiation}\setcounter{section}{6}\setcounter{subsection}{0}\handsectionentry{The Indefinite Integral}}
\expandafter\def\csname handheader@57\endcsname{6.2\quad Lebesgue's Differentiation Theorem}
\expandafter\def\csname handentry@57\endcsname{\setcounter{section}{6}\setcounter{subsection}{1}\handsectionentry{Lebesgue's Differentiation Theorem}}
\expandafter\def\csname handheader@58\endcsname{6.2\quad Lebesgue's Differentiation Theorem}
\expandafter\def\csname handentry@58\endcsname{}
\expandafter\def\csname handheader@59\endcsname{6.2\quad Lebesgue's Differentiation Theorem}
\expandafter\def\csname handentry@59\endcsname{}
\expandafter\def\csname handheader@60\endcsname{6.2\quad Lebesgue's Differentiation Theorem}
\expandafter\def\csname handentry@60\endcsname{}
\expandafter\def\csname handheader@61\endcsname{6.2\quad Lebesgue's Differentiation Theorem}
\expandafter\def\csname handentry@61\endcsname{}
\expandafter\def\csname handheader@62\endcsname{6.2\quad Lebesgue's Differentiation Theorem}
\expandafter\def\csname handentry@62\endcsname{}
\expandafter\def\csname handheader@63\endcsname{6.2\quad Lebesgue's Differentiation Theorem}
\expandafter\def\csname handentry@63\endcsname{}
\expandafter\def\csname handheader@64\endcsname{6.2\quad Lebesgue's Differentiation Theorem}
\expandafter\def\csname handentry@64\endcsname{}
\expandafter\def\csname handheader@65\endcsname{6.2\quad Lebesgue's Differentiation Theorem}
\expandafter\def\csname handentry@65\endcsname{}
\expandafter\def\csname handheader@66\endcsname{6.2\quad Lebesgue's Differentiation Theorem}
\expandafter\def\csname handentry@66\endcsname{}
\expandafter\def\csname handheader@67\endcsname{6.3\quad Vitali Covering Lemma}
\expandafter\def\csname handentry@67\endcsname{\setcounter{section}{6}\setcounter{subsection}{2}\handsectionentry{Vitali Covering Lemma}}
\expandafter\def\csname handheader@68\endcsname{6.4\quad Differentiation of Monotone Functions}
\expandafter\def\csname handentry@68\endcsname{\setcounter{section}{6}\setcounter{subsection}{3}\handsectionentry{Differentiation of Monotone Functions}}
\expandafter\def\csname handheader@69\endcsname{6.4\quad Differentiation of Monotone Functions}
\expandafter\def\csname handentry@69\endcsname{}
\expandafter\def\csname handheader@70\endcsname{6.4\quad Differentiation of Monotone Functions}
\expandafter\def\csname handentry@70\endcsname{}
\expandafter\def\csname handheader@71\endcsname{6.4\quad Differentiation of Monotone Functions}
\expandafter\def\csname handentry@71\endcsname{}
\expandafter\def\csname handheader@72\endcsname{6.5\quad Absolutely continuous and singular functions}
\expandafter\def\csname handentry@72\endcsname{\setcounter{section}{6}\setcounter{subsection}{4}\handsectionentry{Absolutely continuous and singular functions}}
\expandafter\def\csname handheader@73\endcsname{6.5\quad Absolutely continuous and singular functions}
\expandafter\def\csname handentry@73\endcsname{}
\expandafter\def\csname handheader@74\endcsname{6.5\quad Absolutely continuous and singular functions}
\expandafter\def\csname handentry@74\endcsname{}
\expandafter\def\csname handheader@75\endcsname{6.5\quad Absolutely continuous and singular functions}
\expandafter\def\csname handentry@75\endcsname{}
\expandafter\def\csname handheader@76\endcsname{6.6\quad Convex Functions}
\expandafter\def\csname handentry@76\endcsname{\setcounter{section}{6}\setcounter{subsection}{5}\handsectionentry{Convex Functions}}
\expandafter\def\csname handheader@77\endcsname{6.6\quad Convex Functions}
\expandafter\def\csname handentry@77\endcsname{}
\expandafter\def\csname handheader@78\endcsname{6.6\quad Convex Functions}
\expandafter\def\csname handentry@78\endcsname{}
\expandafter\def\csname handheader@79\endcsname{6.6\quad Convex Functions}
\expandafter\def\csname handentry@79\endcsname{}
\begin{document}

\hypersetup{pageanchor=false}
\begin{titlepage}
\thispagestyle{empty}
\null
\begin{tikzpicture}[remember picture,overlay]
\node[anchor=center,inner sep=0pt] at ([yshift=-0.382\paperheight]current page.north)
  {{\sffamily\bfseries\fontsize{40}{50}\selectfont 实变函数}};
\node[anchor=center,inner sep=0pt] at ([yshift=-0.618\paperheight]current page.north)
  {{\rmfamily\fontsize{24}{32}\selectfont By 芬达不带气儿}};
\end{tikzpicture}
\end{titlepage}
\pagenumbering{roman}
\hypersetup{pageanchor=true}
\pagestyle{plain}
\renewcommand{\contentsname}{目录}
{\small\setlength{\parskip}{0pt}\tableofcontents}
\clearpage
\pagenumbering{arabic}
\pagestyle{fancy}

% SOURCE PAGE 001
% CHAPTER: 2 | Fuctions of Bounded Variation and the Riemann-Stieltjes Integral
% SECTION: 2.1 | Functions of Bounded Variation
\begin{handpage}{1}
\handtitle{72}{144}{795}{Chapter 1\quad Fuctions of Bounded Variation and the Riemann--Stieltjes Integral}
\handtitle{72}{173}{730}{\S1.1\quad Functions of Bounded Variation}
\handline{72}{198}{790}{\textbf{Def.} $f$ 在 $[a,b]$ 上有限. $\Gamma=\{a=x_0<x_1<\cdots<x_m=b\}$ 是 $[a,b]$ 的一个分划}
\handline{72}{222}{721}{$S_\Gamma:=S_\Gamma[f;a,b]:=\sum_{i=1}^m|f(x_i)-f(x_{i-1})|$}
\handline{72}{247}{718}{$V:=V[f;a,b]:=\sup_\Gamma S_\Gamma$（无歧义情况下也可记作 $V[a,b]$ 或 $V(f)$）}
\handline{72}{274}{723}{若 $V<+\infty$，则 $f$ 称为 $[a,b]$ 上的有界变差函数\ann{（Bounded Variation）}}
\handline{72}{302}{797}{\textbf{Prop.} \circled{1} 若 $f$ 是 B.V.，则 $f$ 在 $[a,b]$ 上有界}
\handline{72}{329}{733}{\circled{2} 若 $f,g$ 是 B.V.，则 $\alpha f+\beta g$，$f\cdot g$ 是 B.V.}
\handline{72}{356}{709}{若 $\exists\varepsilon>0$，s.t. $|g(x)|\ge\varepsilon$，$\forall x\in[a,b]$，则 $f/g$ 是 B.V.}
\handline{72}{385}{804}{\textbf{Lemma.} 若 $\bar\Gamma$ 是 $\Gamma$ 的一个加细，则 $S_\Gamma\le S_{\bar\Gamma}$\textcolor{Teal}{（用三角不等式）}}
\handline{72}{413}{798}{\textbf{Thm.} \circled{1} 若 $[a',b']\subseteq[a,b]$，则 $V[a',b']\le V[a,b]$}
\handline{72}{443}{733}{\circled{2} 若 $a<c<b$，则 $V[a,b]=V[a,c]+V[c,b]$}
\handline{72}{471}{791}{\textbf{Pf.} \circled{1} 可由 \circled{2} 推出，只证 \circled{2}}
\handline{72}{499}{725}{设 $\Gamma_1,\Gamma_2$ 分别是 $[a,c]$ 和 $[c,b]$ 的分划，则 $\Gamma=\Gamma_1\cup\Gamma_2$ 是 $[a,b]$ 的一个分划.}
\handline{72}{530}{722}{$S_\Gamma[a,b]=S_{\Gamma_1}[a,c]+S_{\Gamma_2}[c,b]$}
\handline{72}{559}{726}{$\therefore\quad S_\Gamma[a,b]\le V[a,c]+V[c,b]\Rightarrow V[a,b]\le V[a,c]+V[c,b]$}
\handline{72}{590}{663}{$S_{\Gamma_1}[a,c]+S_{\Gamma_2}[c,b]\le V[a,b]\Rightarrow V[a,c]+V[c,b]\le V[a,b]$}
\handline{72}{622}{726}{$\therefore V[a,b]=V[a,c]+V[c,b]$.}
\handline{72}{651}{818}{\textbf{Def.} $f$ 在 $[a,b]$ 上有限，$\Gamma=\{x_i\}_{i=0}^m$ 是 $[a,b]$ 的一个分划}
\handline{72}{687}{736}{$P_\Gamma:=P_\Gamma[f;a,b]:=\sum_{i=1}^m(f(x_i)-f(x_{i-1}))^+$，$N_\Gamma:=N_\Gamma[f;a,b]:=\sum_{i=1}^m(f(x_i)-f(x_{i-1}))^-$}
\handline{72}{716}{737}{$P:=\sup_\Gamma P_\Gamma$ 称为正变差. $N:=\sup_\Gamma N_\Gamma$ 称为负变差}
\handline[Teal]{72}{753}{740}{$P_\Gamma+N_\Gamma=S_\Gamma$，$P_\Gamma-N_\Gamma=f(b)-f(a)$}
\end{handpage}

% SOURCE PAGE 002
\begin{handpage}{2}
\handline{72}{143}{817}{\textbf{Thm.} $P,N,V$ 中只要有一个 $<+\infty$，则三个全部 $<+\infty$，且 $P+N=V$，$P-N=f(b)-f(a)$}
\handline[Teal]{72}{168}{755}{$P=\frac12[V+f(b)-f(a)]$，$N=\frac12[V-f(b)+f(a)]$}
\handline{72}{197}{821}{\textbf{Pf.} \circled{1} 若 $V<+\infty$，$\forall\Gamma$，$P_\Gamma+N_\Gamma=S_\Gamma\le V$. $\therefore P_\Gamma,N_\Gamma\le V\Rightarrow P,N\le V<+\infty$}
\handline{72}{223}{747}{若 $P$ 或 $N<+\infty$，由 $P_\Gamma-N_\Gamma=f(b)-f(a)$ 得 $P,N<+\infty$.}
\handline{72}{251}{703}{而 $\forall\Gamma$，$S_\Gamma=P_\Gamma+N_\Gamma\le P+N\Rightarrow V\le P+N<+\infty$ \textcolor{Teal}{$(\star)$}}
\handline{72}{281}{769}{\circled{2} $P,N,V<+\infty$ 时，设 $\{P_{\Gamma_k}\}_{k=1}^\infty$，$\lim_{k\to\infty}P_{\Gamma_k}=P$，则 $N_{\Gamma_k}=P_{\Gamma_k}-f(b)+f(a)\longrightarrow P-f(b)+f(a)$}
\handline{72}{309}{729}{而 $N_{\Gamma_k}\le N$，$\therefore P-f(b)+f(a)\le N$}
\handline{72}{337}{735}{若 $P-f(b)+f(a)<N$，由 $N$ 定义，$\exists\Gamma$，s.t. $N_\Gamma>P-f(b)+f(a)$}
\handline{72}{368}{724}{$\therefore P_\Gamma=N_\Gamma+f(b)-f(a)>P$，矛盾！$\therefore P-f(b)+f(a)=N$}
\handline{72}{401}{731}{$\therefore\lim_{k\to\infty}N_{\Gamma_k}=P-f(b)+f(a)=N$}
\handline{72}{434}{729}{由 $S_{\Gamma_k}=P_{\Gamma_k}+N_{\Gamma_k}$，$S_{\Gamma_k}\le V$，$\therefore\lim_{k\to\infty}S_{\Gamma_k}=P+N\le V$}
\handline{72}{462}{703}{结合\textcolor{Teal}{$(\star)$}可知 $P+N=V$}
\handline{72}{500}{846}{\textbf{Thm.} \ann{（Jordan's Theorem）}}
\handline{72}{531}{799}{$f$ 是 $[a,b]$ 上的 B.V. 函数，当且仅当 $f$ 可以写为两个有界单增函数之差.}
\handline{72}{567}{856}{\textbf{Pf.} “$\Rightarrow$” $V[a,b]<+\infty$，设 $x\in[a,b]$，$[a,x]\subseteq[a,b]\Rightarrow V[a,x]\le V[a,b]<+\infty$}
\handline{72}{599}{732}{于是 $P[a,x],N[a,x]<+\infty$. 设 $f_1(x):=P[a,x]$，$f_2(x):=N[a,x]$}
\handline{72}{629}{733}{则 $f_1,f_2$ 在 $[a,b]$ 上单调增.}
\handline{72}{668}{739}{且 $f_1(x)-f_2(x)=f(x)-f(a)$. $\therefore f(x)=f_1(x)-f_2(x)+f(a)$ 为两个单调函数之差}
\handline{72}{708}{815}{“$\Leftarrow$” $f=f_1-f_2$，其中 $f_1,f_2$ 在 $[a,b]$ 上有界且单增.}
\handline{72}{744}{731}{$\therefore f_1,f_2$ 是 $[a,b]$ 上的 B.V. 函数.}
\handline{72}{778}{753}{由线性性质知 $f$ 是 $[a,b]$ 上的 B.V. 函数.}
\handline[Teal]{72}{814}{855}{此定理可改写为：}
\handline[Teal]{72}{853}{804}{$f$ 是 $[a,b]$ 上 B.V. 函数 $\Longleftrightarrow f$ 可写为两个有界单减函数之和.}
\end{handpage}

% SOURCE PAGE 003
\begin{handpage}{3}
\handline{72}{139}{778}{\textbf{Ex.}}
\handline{72}{165}{775}{\circled{1} $f$ 在 $[a,b]$ 上单调，则 $f$ 是 B.V. 且 $\forall\Gamma$，$S_\Gamma=|f(b)-f(a)|$，$\therefore V=|f(b)-f(a)|$.}
\handline{72}{193}{780}{\circled{2} 设 $[a,b]=\bigcup_{k=1}^n[a_k,a_{k+1}]$，$f$ 在 $[a_k,a_{k+1}]$ 上单调（$\forall k$）}
\handline{72}{222}{735}{则 $V[a,b]=\sum_{k=1}^nV[a_k,a_{k+1}]=\sum_{i=1}^n|f(a_{i+1})-f(a_i)|$}
\handline{72}{252}{272}{\circled{3} $f=\begin{cases}1&x=0\\0&x\in[-1,1]\setminus\{0\}\end{cases}$}
\handline{329}{247}{529}{若 $0\in\Gamma$ 则 $S_\Gamma=2$}
\handline{72}{283.276}{529}{若 $0\notin\Gamma$ 则 $S_\Gamma=0$}
\handline{490}{264}{368}{$\Rightarrow V=2$}
\handline[Teal]{72}{316}{737}{故可以选取 $\{\Gamma_k\}$，使得 $\lim_{k\to\infty}S_{\Gamma_k}\ne V$}
\handline{72}{347}{291}{\circled{4} $f=\begin{cases}1&x\in\Q\cap[0,1]\\0&x\in[0,1]\setminus\Q\end{cases}$}
\handline{383}{347}{488}{$\forall[a,b]\subset[0,1]$，$[a,b]$ 包含无穷有理数和无理数}
\handline{72}{378}{489}{$\therefore V[a,b]=+\infty$}
\handline[NoteRed]{72}{411}{807}{\circled{5} 闭区间上的连续函数不一定是 B.V. 函数.}
\handline{72}{444}{771}{设 $\{a_n\},\{d_n\}$ 是 $[0,1]$ 上单减的数列，满足 $a_1=1$，$\lim_{n\to\infty}a_n=\lim_{n\to\infty}d_n=0$，$\sum_{n=1}^\infty d_n=+\infty$.}
\handline{72}{476}{789}{在 $[a_{k+1},a_k]$ 上，$f$ 是以线段 $\overline{a_{k+1}a_k}$ 为底，高为 $d_k$ 的等腰三角形的腰，$f(a_{k+1})=f(a_k)=0$.}
\handdiagram{603}{490}{325}{160}{
\draw[->](0,0)--(7,0);\draw[->](0,0)--(0,3.5);
\node[below,font=\tiny]at(0,-.15){$0$};\node[below,font=\tiny]at(6,0){$1(a_1)$};
\draw[NoteRed,thick](.3,0)--(.7,.8)--(1.1,0)--(1.5,1.2)--(1.9,0);
\draw[NoteRed,thick](3.4,0)--(3.825,2.1)--(4.25,0)--(4.675,2.7)--(5.1,0)--(5.55,3.2)--(6,0);
\foreach \x/\lab in {.3/{a_{k+2}},1.1/{a_{k+1}},1.9/{a_k},3.4/{a_4},4.25/{a_3},5.1/{a_2}}{\node[below,font=\tiny]at(\x,0){$\lab$};}
\node[below]at(2.65,0){$\cdots$};
\foreach \x/\h/\lab in {.7/.8/{d_{k+1}},1.5/1.2/{d_k},3.825/2.1/{d_3},4.675/2.7/{d_2},5.55/3.2/{d_1}}{\draw[NoteRed,dashed](0,\h)--(\x,\h);\node[NoteRed,left]at(0,\h){$\lab$};}}
\handline{72}{511}{455}{记 $[a_{k+1},a_k]$ 的中点为 $m_k$，$f(m_k)=d_k$.}
\handline{72}{546}{459}{取 $\Gamma_k=\{0\}\cup\{a_j\}_{j=1}^{k+1}\cup\{m_j\}_{j=1}^k$}
\handline{72}{578}{459}{则 $S_{\Gamma_k}=2\sum_{j=1}^kd_j\longrightarrow+\infty$（$k\to\infty$）}
\handline{72}{616}{459}{$\therefore V[f;0,1]=+\infty$，不是 B.V. 函数.}
\handline[NoteRed]{72}{695}{782}{事实上还可以构造在闭区间上连续，但在任意开子区间上不一定是 B.V. 的函数}
\handline[NoteRed]{72}{731}{691}{以及连续的 B.V. 函数在任意开子区间上不单调.}
\handline{72}{771}{829}{\circled{6} $f$：Lipschitz 连续，在 $[a,b]$ 上，$\forall x,y\in[a,b]$，$|f(x)-f(y)|\le L|x-y|$（$L>0$）.}
\handline{72}{809}{776}{设 $\Gamma=\{a_k\}_{k=0}^n$，$S_\Gamma=\sum_{k=1}^n|f(a_k)-f(a_{k-1})|\le\sum_{k=1}^nL|a_k-a_{k-1}|=L(b-a)$}
\handline{72}{845}{735}{$\therefore V[f;a,b]\le L(b-a)<+\infty$，$\therefore f$ 是 $[a,b]$ 上的 B.V. 函数.}
\end{handpage}

% SOURCE PAGE 004
\begin{handpage}{4}
\handline{72}{150}{816}{\textbf{Thm.} B.V. 函数只有第一类间断点，且至多可数个.}
\handline{72}{176}{810}{\textbf{Pf.} 设 $f$ 是 $[a,b]$ 上的 B.V. 函数，不妨设 $f$ 在 $[a,b]$ 上有界且单增 \textcolor{Teal}{(Jordan's Theorem)}}
\handline{72}{202}{744}{于是 $f$ 只有第一类间断点，即跳跃间断点.}
\handline{72}{229}{744}{令 $D$ 表示 $f$ 的所有间断点的集合，于是 $D=\bigcup_{k=1}^\infty\{x\mid f(x+)-f(x-)\ge\frac1k\}$.}
\handline{72}{256}{751}{由于 $f$ 有界，故 $\forall k\in\N$，$\#\{x\mid f(x+)-f(x-)\ge\frac1k\}<+\infty$，$\therefore D$ 是可数集.}
\handline{72}{283}{826}{\textbf{Thm.} $f\in C[a,b]$，则 $V=\lim_{\|\Gamma\|\to0}S_\Gamma$，即 $\forall M<V,\exists\delta>0$，s.t. $\|\Gamma\|<\delta$ 时，$S_\Gamma>M$.}
\handline[AnnotationBlue]{72}{309}{757}{$\Gamma=\{x_0,x_1,\ldots,x_n\}$，$\|\Gamma\|:=\max_{1\le i\le n}\{x_i-x_{i-1}\}$.}
\handline{72}{335}{823}{\textbf{Pf.} $\forall M<V$，取 $\mu>0$，s.t. $M+\mu<V$.}
\handline{72}{362}{747}{取 $\bar\Gamma=\{\bar x_j\}_{j=0}^k$，s.t. $S_{\bar\Gamma}>M+\mu$.}
\handline{72}{389}{749}{$f\in C[a,b]$，故在 $[a,b]$ 上一致连续，$\therefore\exists\eta>0$，s.t. $x_1,x_2\in[a,b],|x_1-x_2|<\eta$ 时}
\handline{72}{417}{731}{$|f(x_1)-f(x_2)|<\frac{\mu}{2(k+1)}$.}
\handline{72}{445}{750}{设 $\Gamma$ 为 $[a,b]$ 的一个分划，$\delta=\min\{\eta,\min_{1\le j\le n}\{\bar x_j-\bar x_{j-1}\}\}$.}
\handline{72}{473}{764}{记 $\Gamma=\{x_i\}_{i=0}^m$，$S_\Gamma=\sum_{i=1}^m|f(x_i)-f(x_{i-1})|=\sum'+\sum''$.}
\handline{72}{504}{763}{$\sum''$ 包含了使得 $(x_{i-1},x_i)$ 中有某个 $\bar x_j$ 的 $|f(x_i)-f(x_{i-1})|$，其余项放在 $\sum'$ 中.}
\handline{72}{534}{757}{于是 $\sum''$ 中至多有 $k+1$ 项.}
\handline{72}{563}{760}{令 $\Gamma'=\Gamma\cup\bar\Gamma$，$S_{\Gamma'}=\sum'+\sum'''$.}
\handline{72}{593}{762}{$\sum'''$ 可由 $\sum''$ 中的项 $|f(x_{i-1})-f(x_i)|$ 替换为 $|f(\bar x_j)-f(x_{i-1})|+|f(x_i)-f(\bar x_j)|$ 得到.}
\handline{72}{625}{746}{当 $\|\Gamma\|<\delta$ 时 $\sum'''<2(k+1)\cdot\frac{\mu}{2(k+1)}=\mu$，$\therefore\sum'=S_{\Gamma'}-\sum'''>S_{\Gamma'}-\mu$.}
\handline{72}{657}{756}{由于 $\Gamma'$ 是 $\bar\Gamma$ 的加细，故 $S_{\Gamma'}\ge S_{\bar\Gamma}$，$\therefore\sum'>S_{\Gamma'}-\mu\ge S_{\bar\Gamma}-\mu$.}
\handline{72}{689}{746}{$\therefore S_\Gamma\ge\sum'>S_{\Gamma'}-\mu\ge S_{\bar\Gamma}-\mu>M$.}
\handline{72}{721}{856}{\textbf{Cor.} 若 $f\in C^1[a,b]$，则 $V=\int_a^b|f'(x)|\,dx$，$P=\int_a^b\{f'(x)\}^+\,dx$，$N=\int_a^b\{f'(x)\}^-\,dx$.}
\handline{72}{751}{856}{\textbf{Pf.} 由中值定理，设 $\Gamma=\{x_i\}_{i=0}^m$.}
\handline{72}{784}{762}{$S_\Gamma=\sum_{i=1}^m|f(x_i)-f(x_{i-1})|=\sum_{i=1}^m|f'(\xi_i)|(x_i-x_{i-1})$，$\xi_i\in(x_{i-1},x_i)$.}
\handline{72}{818}{778}{由定积分定义，$\|\Gamma\|\to0$ 时 $S_\Gamma\to\int_a^b|f'(x)|\,dx$，$\therefore V=\int_a^b|f'(x)|\,dx$.}
\handline{72}{851}{767}{$\therefore P=\frac12(V+f(b)-f(a))=\frac12\bigl(\int_a^b|f'(x)|\,dx+\int_a^bf'(x)\,dx\bigr)=\int_a^b\{f'(x)\}^+\,dx$.}
\handline{72}{885}{729}{$N=\frac12(V-f(b)+f(a))=\frac12\bigl(\int_a^b|f'(x)|\,dx-\int_a^bf'(x)\,dx\bigr)=\int_a^b\{f'(x)\}^-\,dx$.}
\end{handpage}

% SOURCE PAGE 005
% SECTION: 2.2 | Rectifiable Curves
\begin{handpage}{5}
\handtitle{72}{153}{790}{\S1.2 Rectifiable Curves}
\handline{72}{184}{796}{\textbf{Def.} $C:\begin{cases}x=\varphi(t),&a\le t\le b\\y=\psi(t)\end{cases}$，$C$ 的图像为 $\{(x,y)\mid x=\varphi(t),y=\psi(t),a\le t\le b\}$.}
\handline{72}{234}{729}{$\Gamma=\{a=t_0<t_1<\cdots<t_m=b\}$，$P_i=(\varphi(t_i),\psi(t_i))$，将 $P_0,P_1,\ldots,P_m$ 用折线连接.}
\handline{72}{264}{732}{折线段总长为 $l(\Gamma)=\sum_{i=1}^m\bigl([\varphi(t_i)-\varphi(t_{i-1})]^2+[\psi(t_i)-\psi(t_{i-1})]^2\bigr)^{\frac12}$.}
\handline{72}{294}{730}{$C$ 的长度 $L:=l(C)=\sup_\Gamma l(\Gamma)$，若 $L<+\infty$ 则称 $C$ 是可求长曲线.}
\handline{72}{325}{812}{\textbf{Thm.} $C$ 可求长 $\Longleftrightarrow\varphi(t),\psi(t)$ 是 $[a,b]$ 上 B.V. 函数.}
\handline{72}{355}{740}{且 $V(\varphi),V(\psi)\le L\le V(\varphi)+V(\psi)$.}
\handline{72}{386}{808}{\textbf{Pf.} ``$\Rightarrow$'' 若 $C$ 可求长，设 $\Gamma$ 是 $[a,b]$ 的一个分划.}
\handline{72}{415}{766}{$l(\Gamma)=\sum_{i=1}^m\bigl([\varphi(t_i)-\varphi(t_{i-1})]^2+[\psi(t_i)-\psi(t_{i-1})]^2\bigr)^{\frac12}\ge\sum_{i=1}^m|\varphi(t_i)-\varphi(t_{i-1})|,\ \sum_{i=1}^m|\psi(t_i)-\psi(t_{i-1})|$}
\handline{72}{444}{420}{$\therefore V(\varphi),V(\psi)\le L<+\infty$}
\handline{523}{436}{160}{$=S_\Gamma(\varphi)$}
\handline{700}{436}{145}{$=S_\Gamma(\psi)$}
\handline{72}{473}{738}{$\therefore\varphi,\psi$ 是 $[a,b]$ 上的 B.V. 函数.}
\handline{72}{505}{758}{``$\Leftarrow$'' $\sqrt{a^2+b^2}\le a+b$（$a,b>0$），若 $\varphi,\psi$ 是 $[a,b]$ 上 B.V. 函数.}
\handline{72}{537}{738}{$\therefore l(\Gamma)\le S_\Gamma(\varphi)+S_\Gamma(\psi)$，$\therefore L\le V(\varphi)+V(\psi)<+\infty$.}
\handline{72}{571}{735}{$\therefore C$ 可求长.}
\end{handpage}

% SOURCE PAGE 006
% SECTION: 2.3 | The Riemann-Stieltjes Integral
\begin{handpage}{6}
\handtitle{72}{112}{834}{\S1.3 The Riemann-Stieltjes Integral}
\handline{72}{143}{828}{\textbf{Def.} $f,\varphi$ 在 $[a,b]$ 上有限，$\Gamma=\{x_i\}_{i=0}^m$ 是 $[a,b]$ 的一个分划.}
\handline{72}{166}{751}{任选一列介值点 $\{\xi_i\}_{i=1}^m$，$x_{i-1}\le\xi_i\le x_i$.}
\handline{72}{194}{751}{$R_\Gamma:=\sum_{i=1}^mf(\xi_i)[\varphi(x_i)-\varphi(x_{i-1})]$ 称为 $\Gamma$ 的 Riemann-Stieltjes 和.}
\handline{72}{224}{761}{若 $I=\lim_{\|\Gamma\|\to0}R_\Gamma$ 存在且有限，即 $\forall\varepsilon>0,\exists\delta>0$，s.t. $\forall\Gamma,\|\Gamma\|<\delta,\ |I-R_\Gamma|<\varepsilon$.}
\handline{72}{254}{763}{则 $I$ 称为 $f$ 关于 $\varphi$ 在 $[a,b]$ 上的 Riemann-Stieltjes 积分，记 $I=\int_a^bf(x)\,d\varphi(x)=\int_a^bf\,d\varphi$.}
\handline{72}{286}{840}{Cauchy 准则：}
\handline{72}{315}{753}{$\int_a^bf\,d\varphi$ 存在 $\Longleftrightarrow\forall\varepsilon>0,\exists\delta>0$，s.t. $\|\Gamma\|,\|\Gamma'\|<\delta,\ |R_\Gamma-R_{\Gamma'}|<\varepsilon$.}
\handline[Teal]{72}{347}{842}{\circled{1} 若 $\varphi(x)=x$，则 $\int_a^bf\,d\varphi=\int_a^bf(x)\,dx$ 为 Riemann 积分.}
\handline[Teal]{72}{376}{842}{\circled{2} $f\in C[a,b]$，$\varphi\in C^1[a,b]$，则 $\int_a^bf\,d\varphi=\int_a^bf(x)\varphi'(x)\,dx$.}
\handline[AnnotationBlue]{72}{405}{856}{\textbf{Pf.} 由中值定理 $R_\Gamma=\sum_{i=1}^mf(\xi_i)[\varphi(x_i)-\varphi(x_{i-1})]=\sum_{i=1}^mf(\xi_i)\varphi'(\eta_i)(x_i-x_{i-1})$，$x_{i-1}\le\xi_i,\eta_i\le x_i$.}
\handline[AnnotationBlue]{72}{433}{805}{结合 $\varphi'$ 的一致连续性，$\forall\varepsilon>0,\exists\delta>0$，s.t. $|x-y|<\delta$，$|\varphi'(x)-\varphi'(y)|<\frac{\varepsilon}{M(b-a)}$.}
\handline[AnnotationBlue]{72}{463}{795}{$\therefore\bigl|R_\Gamma-\sum_{i=1}^mf(\xi_i)\varphi'(\xi_i)(x_i-x_{i-1})\bigr|=\bigl|\sum_{i=1}^mf(\xi_i)(\varphi'(\xi_i)-\varphi'(\eta_i))(x_i-x_{i-1})\bigr|$.}
\handline[AnnotationBlue]{72}{495}{600}{当 $\|\Gamma\|<\delta$ 时 $|\xi_i-\eta_i|\le\|\Gamma\|<\delta$. 上式 $\le M(b-a)\cdot\frac{\varepsilon}{M(b-a)}=\varepsilon$.}
\handline[NoteRed]{693}{489}{163}{（其中 $M=\|f\|_\infty$）}
\handline[AnnotationBlue]{72}{529}{776}{$\therefore\lim_{\|\Gamma\|\to0}R_\Gamma[f,\varphi]-R_\Gamma[f\varphi',x]=0$. $\therefore\int_a^bf\,d\varphi=\int_a^bf\varphi'\,dx$.}
\handline[Teal]{72}{560}{844}{\circled{3} $\varphi$：$[a,b]$ 上的阶梯函数. $a=\alpha_0<\alpha_1<\cdots<\alpha_m=b$.}
\handline[Teal]{72}{590}{812}{$\varphi$ 在每个小区间 $[\alpha_{i-1},\alpha_i]$ 上取值为常数，$1\le i\le m$.}
\handline[Teal]{72}{622}{812}{令 $\varphi(\alpha_i+)=\lim_{x\to\alpha_i+}\varphi(x)$，$i=0,1,2,\ldots,m-1$，$\varphi(\alpha_i-)=\lim_{x\to\alpha_i-}\varphi(x)$，$i=1,2,\ldots,m$.}
\handline[Teal]{72}{656}{830}{$d_i=\varphi(\alpha_i+)-\varphi(\alpha_i-)$，$i=1,2,\ldots,m-1$，$d_0=\varphi(\alpha_0+)-\varphi(\alpha_0)$，$d_m=\varphi(\alpha_m)-\varphi(\alpha_m-)$.}
\handline[Teal]{72}{689}{814}{于是对于 $f\in C[a,b]$，$\int_a^bf\,d\varphi=\sum_{i=0}^mf(\alpha_i)d_i$.}
\handline[AnnotationBlue]{72}{727}{844}{\textbf{Ex.} Heaviside function $H(x)=\begin{cases}0&x<0\\1&x\ge0\end{cases}$，$f$ 在 $0$ 处连续.}
\handline[AnnotationBlue]{72}{761}{746}{于是 $\int_{-1}^1f\,dH=f(0)$.}
\handline[NoteRed]{72}{791}{788}{$f$ 对一个阶梯函数作 Riemann-Stieltjes 积分}
\handline[NoteRed]{72}{825}{792}{结果为每个区间端点的函数值与该点跳幅乘积之和.}
\end{handpage}

% SOURCE PAGE 007
\begin{handpage}{7}
\handline{72}{95}{792}{\circled{4} 若 $\int_a^bf\,d\varphi$ 存在，则 $f,\varphi$ 没有公共间断点.}
\handline[AnnotationBlue]{72}{125}{796}{\textbf{Pf.} 若 $f,\varphi$ 有公共间断点 $x_0\in(a,b)$. \textcolor{NoteRed}{（$x_0=a$ 或 $x_0=b$ 同理）}}
\handline[AnnotationBlue]{72}{153}{750}{$1^\circ$ $x_0$ 不是 $\varphi$ 的可去间断点.}
\handline[AnnotationBlue]{72}{180}{737}{故 $\exists\varepsilon>0$，s.t. $\forall\delta>0$，$\exists x_1,x_2\in(x_0-\frac\delta2,x_0+\frac\delta2)\setminus\{x_0\}$}
\handline[AnnotationBlue]{72}{207}{709}{s.t. $|\varphi(x_1)-\varphi(x_2)|\ge\varepsilon$（$\delta$ 可以充分小，故 $|x_1-x_2|$ 可以充分小）.}
\handline[AnnotationBlue]{72}{236}{734}{设 $\Gamma$ 是 $[a,b]$ 的分划，且 $\|\Gamma\|<\varepsilon$，并取 $x_{i_0-1}=x_1$，$x_{i_0}=x_2$.}
\handline[AnnotationBlue]{72}{266}{718}{$f$ 在 $x_0$ 不连续，$\therefore\exists\mu>0$，对于 $\eta=\min(\frac{|x_1-x_0|}{2},\frac{|x_0-x_2|}{2})>0$，$\exists\xi_{i_0},\xi'_{i_0}\in B(x_0,\eta)$}
\handline[AnnotationBlue]{72}{292}{709}{s.t. $|f(\xi_{i_0})-f(\xi'_{i_0})|\ge\mu$.}
\handline[AnnotationBlue]{72}{321}{714}{设 $\{\xi_i\}_{i=1}^m$ 与 $\{\xi'_i\}_{i=1}^m$ 为两组介值点，除了 $\xi_{i_0}$ 与 $\xi'_{i_0}$ 之外均相同.}
\handline[AnnotationBlue]{72}{351}{726}{于是 $|R_\Gamma-R_{\Gamma'}|=|f(\xi_{i_0})-f(\xi'_{i_0})|\,|\varphi(x_1)-\varphi(x_2)|\ge\varepsilon\mu$，矛盾!}
\handline[AnnotationBlue]{72}{383}{770}{$2^\circ$ $x_0$ 是 $\varphi$ 的可去间断点.}
\handline[AnnotationBlue]{72}{414}{733}{$\lim_{x\to x_0}\varphi(x)=A\ne\varphi(x_0)$，设 $d:=|\varphi(x_0)-L|>0$.}
\handline[AnnotationBlue]{72}{446}{734}{$\therefore\exists\delta>0$，$x\in(x_0-\delta,x_0+\delta)\cap[a,b]$ 时，$|\varphi(x)-L|<\frac d2$.}
\handline[AnnotationBlue]{72}{477}{725}{$\therefore|\varphi(x_0)-\varphi(x)|\ge|\varphi(x_0)-L|-|\varphi(x)-L|>\frac d2$.}
\handline[AnnotationBlue]{72}{512}{724}{取 $x_1\in(x_0-\delta)\cap[a,b]$，$x_2\in(x_0+\delta)\cap[a,b]$.}
\handline[AnnotationBlue]{72}{545}{724}{而 $|x_0-x_1|,|x_0-x_2|$ 可以充分小.}
\handline[AnnotationBlue]{72}{578}{732}{设 $\Gamma$ 是 $[a,b]$ 的分划，$\forall\varepsilon>0$，当 $\|\Gamma\|\le\varepsilon$ 时.}
\handline[AnnotationBlue]{72}{612}{731}{由于 $|x_0-x_1|,|x_0-x_2|$ 可以充分小，故 $x_1,x_0,x_2$ 可作为分划点.}
\handline[AnnotationBlue]{72}{645}{734}{设 $x_{i_0-1}=x_1$，$x_{i_0}=x_2$.}
\handline[AnnotationBlue]{72}{680}{735}{$f$ 在 $x_0$ 不连续，$\therefore\exists\mu>0$，对于 $\eta=\min(\frac{x_0-x_1}{2},\frac{x_2-x_0}{2})>0$.}
\handline[AnnotationBlue]{72}{715}{680}{$\exists\xi_0\in B(x_0,\eta)$，s.t. $|f(\xi_0)-f(x_0)|\ge\mu$.}
\handline[AnnotationBlue]{72}{751}{736}{取 $\{\xi_i\},\{\xi'_i\}$ 为两组介值点，只有在 $[x_{i_0-1},x_{i_0}]$ 上的点不同.}
\handline[AnnotationBlue]{72}{785}{717}{前者取 $\xi_0$，后者取 $x_0$.}
\handline[AnnotationBlue]{72}{821}{747}{则 $|R_\Gamma-R_{\Gamma'}|=|f(\xi_0)-f(x_0)|\cdot|\varphi(\xi_0)-\varphi(x_0)|$}
\handline[AnnotationBlue]{72}{858}{518}{$>\frac d2\cdot\mu$，矛盾!}
\end{handpage}

% SOURCE PAGE 008
\begin{handpage}{8}
\handline{72}{102}{715}{Thm.\quad 若 $\int_a^b f\,d\phi$ 存在，$a<c<b$，则 $\int_a^c f\,d\phi,\int_c^b f\,d\phi$ 都存在且 $\int_a^b f\,d\phi=\int_a^c f\,d\phi+\int_c^b f\,d\phi$.}
\handline{72}{126}{695}{Pf.\quad 要证 $\int_a^c f\,d\phi$ 存在，即证 $\forall\varepsilon>0,\exists\delta>0$，$\|\Gamma_1\|,\|\Gamma_2\|<\delta$ 时，$|R_{\Gamma_1}[a,c]-R_{\Gamma_2}[a,c]|<\varepsilon$.}
\handline{72}{150}{600}{由 $\int_a^b f\,d\phi$ 存在，知 $\exists\delta_0>0$，$\|\Gamma_1'\|,\|\Gamma_2'\|<\delta_0$ 时 $|R_{\Gamma_1'}[a,b]-R_{\Gamma_2'}[a,b]|<\varepsilon$.}
\handline{72}{175}{390}{设 $\Gamma$ 是 $[c,b]$ 的分划，$\|\Gamma\|<\delta_0$.}
\handline{72}{201}{390}{令 $\Gamma_1'=\Gamma\cup\Gamma_1,\ \Gamma_2'=\Gamma\cup\Gamma_2$，$\delta=\delta_0$.}
\handline{72}{225}{510}{$\Rightarrow|R_{\Gamma_1}[a,c]-R_{\Gamma_2}[a,c]|=|R_{\Gamma_1'}[a,b]-R_{\Gamma_2'}[a,b]|<\varepsilon$}
\handline{72}{254}{410}{故 $\int_a^c f\,d\phi$ 存在，同理 $\int_c^b f\,d\phi$ 存在.}
\handline{72}{275}{300}{Thm.\quad\ann{（分部积分）}}
\handline{72}{303}{680}{若 $\int_a^b f\,d\phi$ 存在，则 $\int_a^b\phi\,df$ 存在.且 $\int_a^b f\,d\phi=[f(b)\phi(b)-f(a)\phi(a)]-\int_a^b\phi\,df$.}
\handline{72}{331}{515}{Pf.\quad 设 $\Gamma=\{a=x_0<x_1<x_2<\cdots<x_m=b\}$，$x_{i-1}\le\xi_i\le x_i$.}
\handline{72}{362}{725}{$ R_\Gamma=\sum_{i=1}^m f(\xi_i)[\phi(x_i)-\phi(x_{i-1})]\overset{\text{Abel变换}}{=}-\sum_{i=1}^{m-1}\phi(x_i)[f(\xi_{i+1})-f(\xi_i)]+f(\xi_m)\phi(b)-f(\xi_1)\phi(a)$}
\handline{72}{388}{716}{$=-\sum_{i=1}^{m-1}\phi(x_i)[f(\xi_{i+1})-f(\xi_i)]+f(\xi_m)\phi(b)-f(b)\phi(b)+f(a)\phi(a)-f(\xi_1)\phi(a)+f(b)\phi(b)-f(a)\phi(a)$}
\handline{72}{418}{710}{$=-\sum_{i=1}^{m-1}\phi(x_i)[f(\xi_{i+1})-f(\xi_i)]+\phi(b)(f(b)-f(\xi_m))+\phi(a)(f(\xi_1)-f(a))+f(b)\phi(b)-f(a)\phi(a)$}
\handline{72}{446}{620}{$=f(b)\phi(b)-f(a)\phi(a)-T_\Gamma$，\quad$T_\Gamma$ 是 $\phi$ 对 $f$ 的 Riemann--Stieltjes 和.}
\handline{72}{475}{370}{$f,\phi$ 的地位是对等的.}
\handline{72}{506}{420}{于是 $\lim_{\|\Gamma\|\to0}R_\Gamma$ 存在 $\iff\lim_{\|\Gamma\|\to0}T_\Gamma$ 存在.}
\handline{72}{535}{410}{即 $\int_a^b f\,d\phi$ 存在 $\iff\int_a^b\phi\,df$ 存在.}
\handline{72}{562}{435}{且 $\int_a^b f\,d\phi=[f(b)\phi(b)-f(a)\phi(a)]-\int_a^b\phi\,df$.}
\handline[Teal]{72}{599}{600}{$\phi$ 和 $f$ 只有有限的限制，并没有可导、连续的限制.}
\handline[Teal]{72}{631}{725}{数分中要求了 $f,\phi\in C^1[a,b]$，是因为数分中无法说清楚 $df,dg$ 是什么.}
\handline{72}{666}{770}{Def.\quad$f$：有界.$\phi$：在 $[a,b]$ 上单增.$\Gamma=\{x_i\}_{i=0}^m$，令 $m_i=\inf_{[x_{i-1},x_i]}f(x)$，$M_i=\sup_{[x_{i-1},x_i]}f(x)$.}
\handline{72}{699}{575}{$ L_\Gamma:=\sum_{i=1}^m m_i(\phi(x_i)-\phi(x_{i-1}))$，$ U_\Gamma:=\sum_{i=1}^m M_i(\phi(x_i)-\phi(x_{i-1}))$.}
\handline{72}{735}{560}{分别称为，对于 $\Gamma$ 的 Riemann--Stieltjes 下和与上和.}
\handline{72}{765}{340}{显然.$L_\Gamma\le R_\Gamma\le U_\Gamma$.}
\end{handpage}

% SOURCE PAGE 009
\begin{handpage}{9}
\handline{72}{143}{450}{Lemma.\quad$f$：有界.$\phi$：在 $[a,b]$ 上单增.}
\handline{72}{170}{490}{\circled{1} 若 $\Gamma'$ 是 $\Gamma$ 的加细，则 $L_{\Gamma'}\ge L_\Gamma$，$U_{\Gamma'}\le U_\Gamma$.}
\handline{72}{195}{430}{\circled{2} 若 $\Gamma_1,\Gamma_2$ 是两个分划，则 $L_{\Gamma_1}\le U_{\Gamma_2}$.}
\handline{72}{221}{695}{Pf.\quad\circled{1} 只需证 $\Gamma'$ 比 $\Gamma$ 多了一个分划点的情形.设 $\Gamma=\{x_i\}_{i=0}^m$，$\Gamma'=\Gamma\cup\{x'\}$.}
\handline{72}{250}{280}{设 $x'\in(x_{i-1},x_i)$.}
\handline{72}{286}{820}{$\sup_{[x_{i-1},x']}f(x)[\phi(x')-\phi(x_{i-1})]+\sup_{[x',x_i]}f(x)[\phi(x_i)-\phi(x')]\le\sup_{[x_{i-1},x_i]}f(x)[\phi(x')-\phi(x_{i-1})]+\sup_{[x_{i-1},x_i]}f(x)[\phi(x_i)-\phi(x')]$}
\handline{72}{311}{290}{$=M_i(\phi(x_i)-\phi(x_{i-1}))$}
\handline{72}{339}{310}{$\therefore U_{\Gamma'}\le U_\Gamma$，同理 $L_{\Gamma'}\ge L_\Gamma$.}
\handline{72}{365}{510}{\circled{2} 令 $\Gamma=\Gamma_1\cup\Gamma_2$，则 $L_{\Gamma_1}\le L_{\Gamma_1\cup\Gamma_2}\le U_{\Gamma_1\cup\Gamma_2}\le U_{\Gamma_2}$.}
\handline{72}{397}{795}{Thm.\quad$f\in C[a,b]$，$\phi\in\BV[a,b]$，则 $\int_a^b f\,d\phi$ 存在.且 $|\int_a^b f\,d\phi|\le(\sup_{[a,b]}|f|)\cdot V[\phi;a,b]$.}
\handline{72}{426}{675}{Pf.\circled{1} 由 Jordan 分解定理及线性性，不妨设 $\phi$ 在 $[a,b]$ 上单增.}
\handline{72}{454}{640}{若 $\phi$ 是常数，trivial.设 $\phi$ 不是常数.令 $\Gamma=\{x_i\}_{i=0}^m$ 为一分划.}
\handline{72}{485}{670}{$f\in C[a,b]$，故一致连续.$\therefore\forall\varepsilon>0,\exists\delta>0$，$\|\Gamma\|<\delta$ 时.$M_i-m_i<\dfrac{\varepsilon}{\phi(b)-\phi(a)}$.}
\handline{72}{516}{470}{$\Rightarrow0\le U_\Gamma-L_\Gamma=\sum_{i=1}^m(M_i-m_i)(\phi(x_i)-\phi(x_{i-1}))<\varepsilon$}
\handline{72}{545}{320}{故 $\lim_{\|\Gamma\|\to0}(U_\Gamma-L_\Gamma)=0$.}
\handline{72}{577}{665}{若 $\lim_{\|\Gamma\|\to0}U_\Gamma$ 不存在，则存在两列分划 $\{\Gamma_k'\},\{\Gamma_k''\}$，$\lim_{k\to\infty}\|\Gamma_k'\|=\lim_{k\to\infty}\|\Gamma_k''\|=0$}
\handline{72}{609}{435}{$\exists K\in\N$.$k>K$ 时.$U_{\Gamma_k'}-U_{\Gamma_k''}>\varepsilon$}
\handline{72}{641}{660}{但当 $k$ 充分大时.$|L_{\Gamma_k'}-U_{\Gamma_k'}|<\dfrac\varepsilon2\Rightarrow L_{\Gamma_k'}>U_{\Gamma_k'}-\dfrac\varepsilon2>U_{\Gamma_k''}+\dfrac\varepsilon2$.}
\handline{72}{675}{670}{$\therefore L_{\Gamma_k'}>U_{\Gamma_k''}$，矛盾！$\therefore\lim_{\|\Gamma\|\to0}U_\Gamma$ 存在，于是 $\lim_{\|\Gamma\|\to0}R_\Gamma,\lim_{\|\Gamma\|\to0}L_\Gamma$ 存在.}
\handline{72}{706}{620}{且 $\lim_{\|\Gamma\|\to0}U_\Gamma=\lim_{\|\Gamma\|\to0}L_\Gamma=\lim_{\|\Gamma\|\to0}R_\Gamma=\int_a^b f\,d\phi$，从而 $\int_a^b f\,d\phi$ 存在.}
\handline{72}{740}{710}{\circled{2} $|R_\Gamma|=\left|\sum_{i=1}^m f(\xi_i)(\phi(x_i)-\phi(x_{i-1}))\right|\le(\sup_{[a,b]}|f|)S_\Gamma(\phi)\le(\sup_{[a,b]}|f|)V[\phi;a,b]$}
\handline{72}{771}{465}{$\therefore\left|\int_a^b f\,d\phi\right|\le(\sup_{[a,b]}|f|)V[\phi;b,a]$.}
\handline{72}{800}{340}{Thm.\quad\ann{（积分中值定理）}}
\handline{72}{830}{755}{$f\in C[a,b]$，$\phi\in\BV[a,b]$ 且单增.则 $\exists\xi\in[a,b]\ \text{s.t.}\ \int_a^b f\,d\phi=f(\xi)[\phi(b)-\phi(a)]$.}
\handline{72}{863}{675}{Pf.\quad$(\min_{[a,b]}f)(\phi(b)-\phi(a))\le R_\Gamma\le(\max_{[a,b]}f)(\phi(b)-\phi(a))\quad(\forall\Gamma)$.}
\handline{72}{896}{645}{$\Rightarrow(\min_{[a,b]}f)(\phi(b)-\phi(a))\le\int_a^b f\,d\phi\le(\max_{[a,b]}f)(\phi(b)-\phi(a))$.}
\handline{72}{929}{755}{由连续函数介值性知 $\exists\xi\in[a,b]$.$\text{s.t.}\ \int_a^b f\,d\phi=f(\xi)[\phi(b)-\phi(a)]$.}
\end{handpage}

% SOURCE PAGE 010
% CHAPTER: 3 | Lebesgue Measure and Outer Measure
% SECTION: 3.1 | Lebesgue Outer Measure and the Cantor Set
\begin{handpage}{10}
\handtitle{72}{119}{750}{Chapter 2. Lebesgue Measure and Outer Measure}
\handtitle{72}{151}{760}{\S2.1 Lebesgue Outer Measure and the Cantor Set}
\handline{72}{177}{655}{$E\subseteq\R^n$.$I=\{x=(x_1,\ldots,x_n)\in\R^n\mid a_i\le x_i\le b_i\}$ 表示闭区间.}
\handline{72}{206}{490}{$I$ 的体积定义为 $ V(I):=\prod_{i=1}^n(b_i-a_i)$}
\handline{72}{239}{680}{Def.\quad$S=\{I_i\}_{i=1}^\infty$，$I_i$：闭区间\quad$\sigma(S):=\sum_{i=1}^\infty V(I_i)$，$ E\subseteq\bigcup_{i=1}^\infty I_i$}
\handline{72}{267}{460}{$m^*(E):=\inf\sigma(S)$，称为 $E$ 的外测度.}
\handline{72}{298}{620}{Lemma.\quad$I$：闭区间，则 $m^*(I)=V(I)$.\quad\ann{$(m^*(\partial I)=0)$}}
\handline{72}{328}{405}{Pf.\quad 由定义 $m^*(I)\le V(I)$.}
\handline{72}{359}{355}{设 $ S=\{I_i\}_{i=1}^\infty$，$I\subseteq\bigcup_{i=1}^\infty I_i$.}
\handline{72}{391}{640}{$\forall\varepsilon>0$，$\forall i\in\N$，$\exists I_i^*$，$\text{s.t.}\ I_i\subseteq\operatorname{int}(I_i^*)$.且 $V(I_i^*)\le(1+\varepsilon)V(I_i)$.}
\handline{72}{424}{660}{$ I\subseteq\bigcup_{i=1}^\infty I_i\subseteq\bigcup_{i=1}^\infty\operatorname{int}(I_i^*)$，由紧致性 $\Rightarrow\exists N$.$\text{s.t.}\ I\subseteq\bigcup_{i=1}^N\operatorname{int}(I_i^*)\subseteq\bigcup_{i=1}^N I_i^*$.}
\handline{72}{456}{560}{$\therefore V(I)\le\sum_{i=1}^N V(I_i^*)\le(1+\varepsilon)\sum_{i=1}^N V(I_i)\le(1+\varepsilon)\sigma(S)$.}
\handline{72}{487}{580}{$\therefore V(I)\le\sigma(S)\Rightarrow V(I)\le m^*(I)\quad\therefore m^*(I)=V(I)$}
\handline{72}{522}{475}{Lemma.\quad$E_1\subseteq E_2$.则 $m^*(E_1)\le m^*(E_2)$}
\handline{72}{555}{770}{Pf.\quad$ E_2\subseteq\bigcup_{i=1}^\infty I_i\quad\therefore E_1\subseteq E_2\subseteq\bigcup_{i=1}^\infty I_i\Rightarrow m^*(E_1)\le\sum_{i=1}^\infty V(I_i)\overset{\text{取inf}}{\Rightarrow}m^*(E_1)\le m^*(E_2)$.}
\handline{72}{591}{565}{Lemma.\quad$E_k\subseteq\R^n$.令 $ E=\bigcup_{k=1}^\infty E_k$.则 $ m^*(E)\le\sum_{k=1}^\infty m^*(E_k)$.}
\handline{72}{623}{450}{Pf.\quad 不妨设 $\forall k\in\N\ m^*(E_k)<\infty$.}
\handline{72}{658}{680}{$\forall\varepsilon>0$，$\exists\{I_i^{(k)}\}_{i=1}^\infty$，$ E_k\subseteq\bigcup_{i=1}^\infty I_i^{(k)}$，$\text{s.t.}\ \sum_{i=1}^\infty V(I_i^{(k)})<m^*(E_k)+\frac\varepsilon{2^k}$}
\handline{72}{692}{715}{$ E\subseteq\bigcup_{k=1}^\infty\bigcup_{i=1}^\infty I_i^{(k)}$，$ m^*(E)\le\sum_{k=1}^\infty\sum_{i=1}^\infty V(I_i^{(k)})\le\sum_{k=1}^\infty m^*(E_k)+\frac\varepsilon{2^k}\le\sum_{k=1}^\infty m^*(E_k)+\varepsilon$}
\handline{72}{729}{460}{令 $\varepsilon\to0+$.$\Rightarrow m^*(E)\le\sum_{k=1}^\infty m^*(E_k)$.}
\handline{72}{761}{540}{Ex.\quad\circled{1} 单点集外测度为 $0$.$E=\{p\}$，$m^*(E)=0$}
\handline{72}{792}{465}{\circled{2} 可数集外测度为 $0$.$E=\{p_k\}_{k=1}^\infty$.}
\handline{72}{827}{530}{$ m^*(E)\le\sum_{k=1}^\infty m^*(\{p_k\})=\sum_{k=1}^\infty0=0.\quad\therefore m^*(E)=0$.}
\handline[AnnotationBlue]{72}{862}{550}{不可数集的测度也可以是 $0$.如 Cantor 集.}
\end{handpage}

% SOURCE PAGE 011
\begin{handpage}{11}
\handline{72}{164}{260}{Cantor Set}
\handdiagram{628}{181}{300}{76}{
\draw (0,2)--(9,2);\draw (0,2)--(0,2.2);\draw (9,2)--(9,2.2);
\draw (0,1)--(3,1);\draw (6,1)--(9,1);\foreach \x in {0,3,6,9}{\draw (\x,1)--(\x,1.2);}
\draw (0,0)--(1,0);\draw (2,0)--(3,0);\draw (6,0)--(7,0);\draw (8,0)--(9,0);\foreach \x in {0,1,2,3,6,7,8,9}{\draw (\x,0)--(\x,.2);}
\node[anchor=west] at (9.2,1.3) {$k=1$};\node[anchor=west] at (9.2,.3) {$k=2$};\node at (4.5,-.7) {$\cdots\quad\cdots$};
}
\handline{403}{175}{207}{$C_k$：$\#\{$区间$\}$\qquad 长度}
\handline{72}{202}{160}{$2^k\qquad\dfrac1{3^k}$}
\handline{72}{231}{265}{$\#\{$移去的开区间$\}$：$2^k-1$}
\handline{72}{297}{440}{$ C:=\bigcap_{k=1}^\infty C_k$，称为 Cantor 集.}
\handline{72}{325}{580}{$m^*(C)\le m^*(C_k)=\left(\dfrac23\right)^k\to0\ (k\to\infty)\quad\therefore m^*(C)=0$}
\handline{72}{352}{480}{\circled{1} $C_k$ 是闭区间 $\therefore C=\bigcap_{k=1}^\infty C_k$ 是闭集.}
\handline{72}{380}{805}{\circled{2} $\forall p\in C$，则 $p\in C_k\ (\forall k\in\N)$.故 $p$ 可以用 $C_k$ 的边界点逼近，从而 $C$ 的每个点都为聚点.}
\handline[AnnotationBlue]{72}{406}{535}{由 \circled{1}\circled{2} $\Rightarrow C$ 是完美集.从而不可数.}
\handline[AnnotationBlue]{72}{437}{740}{由 \circled{2} $\Rightarrow$ 令 $A=\{C_k\text{ 的边界点},\ k\in\N\}$.则 $A$ 是可数集.且在 $C$ 中稠密.}
\handline{72}{503}{690}{Cantor Function.\ann{（或 Cantor--Lebesgue Function）}}
\handdiagram{603}{517}{325}{185}{
\draw[->] (0,0)--(10.1,0);\draw[->] (0,0)--(0,7.6);
\node[left] at (0,0) {$0$};\node[left] at (0,1.75) {$\frac14$};\node[left] at (0,3.5) {$\frac12$};\node[left] at (0,5.25) {$\frac34$};\node[left] at (0,7) {$1$};
\draw[Teal,thick] (0,0)--(3,3.5)--(6,3.5)--(9,7);
\draw[AnnotationBlue,thick] (0,0)--(1,1.75)--(2,1.75)--(3,3.5)--(6,3.5)--(7,5.25)--(8,5.25)--(9,7);
\draw[NoteRed,thick] (0,0)--(.333,.875)--(.667,.875)--(1,1.75)--(2,1.75)--(2.333,2.625)--(2.667,2.625)--(3,3.5)--(6,3.5)--(6.333,4.375)--(6.667,4.375)--(7,5.25)--(8,5.25)--(8.333,6.125)--(8.667,6.125)--(9,7);
\draw[AnnotationBlue,dashed] (1,0)--(1,1.75);\draw[AnnotationBlue,dashed] (2,0)--(2,1.75);\draw[Teal,dashed] (3,0)--(3,3.5);\draw[Teal,dashed] (6,0)--(6,3.5);\draw[AnnotationBlue,dashed] (7,0)--(7,5.25);\draw[AnnotationBlue,dashed] (8,0)--(8,5.25);
\node[below,AnnotationBlue] at (1,0) {$\frac19$};\node[below,AnnotationBlue] at (2,0) {$\frac29$};\node[below,Teal] at (3,0) {$\frac13$};\node[below,Teal] at (6,0) {$\frac23$};\node[below,AnnotationBlue] at (7,0) {$\frac79$};\node[below,AnnotationBlue] at (8,0) {$\frac89$};\node[below] at (9,0) {$1$};
\node[Teal] at (4.6,4.1) {$f_1$};\draw[->,AnnotationBlue] (8.1,5.6)--(9.7,5.4);\node[AnnotationBlue,right] at (9.7,5.4) {$f_2$};\draw[->,NoteRed] (2.5,2.5) to[bend right] (3.6,2.1);\node[NoteRed,right] at (3.6,2.1) {$f_3$};
}
\handline{72}{531}{240}{$f_k:[0,1]\to[0,1]$}
\handline{72}{564}{240}{$|f_k-f_{k+1}|\le\dfrac1{2^k}$}
\handline{72}{596}{485}{$\Rightarrow f_n=f_n-f_{n-1}+f_{n-1}-f_{n-2}+\cdots+f_2-f_1+f_1$}
\handline{72}{630}{470}{$|f_n|\le|f_n-f_{n-1}|+|f_{n-1}-f_{n-2}|+\cdots+|f_2-f_1|+|f_1|$}
\handline{72}{661}{335}{$\le\dfrac1{2^{n-1}}+\cdots+\dfrac12+1$}
\handline{72}{695}{325}{故 $\{f_n\}$ 一致收敛.}
\handline{72}{729}{455}{$f:=\lim_{n\to\infty}f_n$，称为 Cantor 函数.}
\handline{72}{762}{605}{$f_n$ 均连续，由连续性定理 $\Rightarrow f$ 是连续函数.}
\handline{72}{795}{680}{$f$ 在每个移去的开区间上为常数，故可微.且 $f'=0$.}
\handline{72}{826}{465}{而移去的开区间的总长度为：}
\handline{72}{861}{520}{$\dfrac13+\left(\dfrac13\right)^2\cdot2+\cdots+\left(\dfrac13\right)^k\cdot2^{k-1}+\cdots=\dfrac13\cdot\dfrac1{1-\frac23}=1$}
\handline[NoteRed]{72}{895}{535}{（称，$f$ 在 $[0,1]$ 上几乎处处可微）}
\end{handpage}

% SOURCE PAGE 012
\begin{handpage}{12}
\handline{72}{125}{670}{Thm.\quad$E\subseteq\R^n$，$\forall\varepsilon>0$，$\exists$ 开集 $G\subseteq\R^n$.$\text{s.t.}\ E\subseteq G$ 且 $m^*(G)\le m^*(E)+\varepsilon$}
\handline{72}{153}{785}{Pf.\quad 不妨设 $m^*(E)<\infty$，否则选 $G=\R^n$.$\forall\varepsilon>0$，$\exists\{I_k\}_{k=1}^\infty$.$\text{s.t.}\ E\subseteq\bigcup_{k=1}^\infty I_k$，且 $ m^*(E)+\frac\varepsilon2>\sum_{k=1}^\infty V(I_k)$.}
\handline{72}{183}{620}{$\forall k\in\N$.取更大的 $I_k^*$.$\text{s.t.}\ I_k\subseteq\operatorname{int}(I_k^*)$.且 $V(I_k^*)\le V(I_k)+\dfrac\varepsilon{2^{k+1}}$}
\handline{72}{214}{360}{令 $ G=\bigcup_{k=1}^\infty\operatorname{int}(I_k^*)$.}
\handline{72}{242}{630}{则 $ m^*(G)\le\sum_{k=1}^\infty V(I_k^*)\le\sum_{k=1}^\infty V(I_k)+\frac\varepsilon{2^{k+1}}=\sum_{k=1}^\infty V(I_k)+\frac\varepsilon2\le m^*(E)+\varepsilon$.}
\handline{72}{271}{640}{Def.\quad 若 $H$ 形如.$ H=\bigcap_{k=1}^\infty G_k$.其中 $G_k$ 为开集，则称 $H$ 为 $G_\delta$ 集.}
\handline{72}{302}{595}{若 $H$ 形如.$ H=\bigcup_{k=1}^\infty F_k$.其中 $F_k$ 为闭集，则称 $H$ 为 $F_\sigma$ 集.}
\handline{72}{332}{650}{Thm.\quad$E\subseteq\R^n$.$\exists$ 一个 $G_\delta$ 集 $H$.$\text{s.t.}\ E\subseteq H$.且 $m^*(H)=m^*(E)$}
\handline{72}{361}{665}{Pf.\quad$\forall k\in\N$.$\exists$ 开集 $G_k$.$\text{s.t.}\ E\subseteq G_k$.且 $m^*(G_k)\le m^*(E)+\dfrac1k$.}
\handline{72}{391}{490}{令 $ H=\bigcap_{k=1}^\infty G_k$，则 $E\subseteq H$.$\therefore m^*(E)\le m^*(H)$.}
\handline{72}{421}{575}{而 $m^*(H)\le m^*(G_k)\le m^*(E)+\dfrac1k$.\quad$\Rightarrow m^*(H)\le m^*(E)$}
\handline{72}{450}{360}{$\therefore m^*(H)=m^*(E)$.}
\handline{72}{483}{450}{Thm.\quad$E\subseteq\R^n$，则 $m^*(E)=m^*(E)'$.}
\handline{72}{512}{350}{Pf.\quad\circled{1} $I$：闭区间.}
\handline{72}{541}{500}{$\forall\varepsilon>0$，取更大的 $I_1$，$\text{s.t.}\ I\subseteq\operatorname{int}(I_1)$ 且 $V(I_1)\le V(I)+\varepsilon$.}
\handline{72}{571}{390}{令 $\operatorname{int}(I_1)=\bigcup_{k=1}^\infty I_k'$\quad$I_k'$：互不相交}
\handline{72}{604}{470}{$\forall N>0.\ \bigcup_{k=1}^N I_k'\subseteq\operatorname{int}(I_1)\subseteq I_1\quad\therefore\sum_{k=1}^N V(I_k')\le V(I_1)$}
\handline{72}{634}{510}{令 $N\to\infty$.$\Rightarrow\sum_{k=1}^\infty V(I_k')\le V(I_1)\le V(I)+\varepsilon=m^*(I)+\varepsilon$}
\handline{72}{671}{790}{$\therefore m^*(I)'\le m^*(I)+\varepsilon\overset{\varepsilon\to0}{\Rightarrow}m^*(I)'\le m^*(I)$.同理 $m^*(I)\le m^*(I)'$.$\therefore m^*(I)=m^*(I)'$}
\handdiagram{615}{432}{313}{121}{
\draw[->] (2,0)--(6,0) node[right] {$x$};\draw[->] (2,0)--(2,3.5) node[above] {$y$};\node[below left] at (2,0) {$0$};
\draw[->,AnnotationBlue] (2,0)--(5,3.1) node[above right] {$x'$};\draw[->,AnnotationBlue] (2,0)--(-.8,3.2) node[above left] {$y'$};
\draw[Teal] (1.5,2.9) .. controls (2.5,3.1) and (3.4,3.1) .. (3.7,2.2) .. controls (4.2,1.2) and (3.7,.7) .. (2.8,1.1) .. controls (1.8,1.5) and (.7,1.4) .. (1.5,2.9);\node[Teal] at (3,3.4) {$E$};
\draw[Teal] (1.2,1.75)--(1.8,2.95);\draw[Teal] (1.7,1.5)--(2.4,3);\draw[Teal] (2.2,1.35)--(3,2.95);\draw[Teal] (2.7,1.1)--(3.5,2.7);\draw[Teal] (3.2,1)--(3.9,2.2);
\draw (3.3,.6) rectangle (3.75,1.8);\node[right] at (3.75,.8) {$I$};\draw[AnnotationBlue] (1.15,3.35)--(2.55,1.9)--(2.25,1.6)--(.85,3.05)--cycle;\node[above left,AnnotationBlue] at (1.15,3.35) {$I'$};
}
\handline[NoteRed]{640}{574}{230}{不管选哪个坐标系，}
\handline[NoteRed]{604}{603}{260}{\unclear{p012-01}{红字第二行“覆盖”前字词}覆盖集合的矩体.}
\handline[NoteRed]{621}{632}{270}{指的都是边界平行于坐标轴的.}
\handline{72}{707}{230}{\circled{2} $E\subseteq\R^n$.}
\handline{72}{741}{570}{$\forall\varepsilon>0,\{I_k\}_{k=1}^\infty\ \text{s.t.}\ E\subseteq\bigcup_{k=1}^\infty I_k$，$\sum_{k=1}^\infty V(I_k)\le m^*(E)+\frac\varepsilon2$}
\handline{72}{775}{620}{$\forall k\in\N.\ \exists\{I_{kj}'\}_{j=1}^\infty\ \text{s.t.}\ I_k\subseteq\bigcup_{j=1}^\infty I_{kj}'$，$\sum_{j=1}^\infty V(I_{kj}')\le V(I_k)+\frac\varepsilon{2^{k+1}}$}
\handline{72}{806}{410}{$\therefore E\subseteq\bigcup_{k=1}^\infty I_k\subseteq\bigcup_{k,j=1}^\infty I_{kj}'$}
\handline{72}{837}{775}{$ m^*(E)'\le\sum_{k=1}^\infty\sum_{j=1}^\infty V(I_{kj}')\le\sum_{k=1}^\infty V(I_k)+\frac\varepsilon{2^{k+1}}\le m^*(E)+\varepsilon.\overset{\varepsilon\to0}{\Rightarrow}m^*(E)'\le m^*(E)$}
\handline{72}{868}{560}{同理 $m^*(E)\le m^*(E)'$.$\therefore m^*(E)=m^*(E)'$}
\handline[AnnotationBlue]{72}{901}{610}{外测度的值不依赖于直角坐标系的选取.}
\handline[AnnotationBlue]{72}{932}{510}{在刚体变换之下保持不变.}
\end{handpage}

% SOURCE PAGE 013
% SECTION: 3.2 | Lebesgue Measurable Sets
\begin{handpage}{13}
\handtitle{72}{143}{650}{\S2.2 Lebesgue Measurable Sets}
\handline{72}{175}{680}{Def.\quad$E\subseteq\R^n$.若 $\forall\varepsilon>0$，$\exists$ 开集 $G$.$\text{s.t.}\ E\subseteq G$，$m^*(G\setminus E)<\varepsilon$}
\handline{72}{198}{645}{则称 $E$ 是 Lebesgue 可测集，并定义它的测度为 $m(E):=m^*(E)$.}
\handline{72}{230}{710}{Ex.\quad\circled{1} 开集可测：$E$ 开集.取 $G=E$.$m^*(G\setminus E)=0\quad\therefore E$ 可测}
\handline{72}{257}{405}{\circled{2} 若 $m^*(E)=0$，则 $E$ 可测.}
\handline{72}{285}{750}{$\forall\varepsilon>0$，$\exists$ 开集 $G$.$\text{s.t.}\ E\subseteq G$.$m^*(G)\le m^*(E)+\varepsilon=\varepsilon$.$\therefore m^*(G\setminus E)\le m^*(G)<\varepsilon$}
\handline{72}{315}{670}{Thm.\quad$E_k$：可测.则 $ E=\bigcup_{k=1}^\infty E_k$ 也可测.且 $ m(E)\le\sum_{k=1}^\infty m(E_k)$.}
\handline{72}{343}{660}{Pf.\quad$\forall\varepsilon>0$.$\forall k\in\N$.$\exists G_k$ 开集 $\text{s.t.}\ E_k\subseteq G_k$.$m^*(G_k\setminus E_k)<\dfrac\varepsilon{2^k}$}
\handline{72}{374}{645}{令 $ G=\bigcup_{k=1}^\infty G_k$.则 $G$ 是开集.$E\subseteq G$.且 $ G\setminus E\subseteq\bigcup_{k=1}^\infty(G_k\setminus E_k)$}
\handline{72}{403}{570}{$ m^*(G\setminus E)\le\sum_{k=1}^\infty m^*(G_k\setminus E_k)<\sum_{k=1}^\infty\frac\varepsilon{2^k}=\varepsilon$}
\handline{72}{432}{650}{$\therefore E$ 是可测集.且 $ m(E)=m^*(E)\le\sum_{k=1}^\infty m^*(E_k)=\sum_{k=1}^\infty m(E_k)$}
\handline{72}{462}{525}{Thm.\quad$I$：闭区间，则 $I$ 可测，且 $m(I)=V(I)$.}
\handline{72}{492}{810}{Pf.\quad$I=\operatorname{int}(I)\sqcup\partial I$.$\operatorname{int}(I)$ 是开集.$\partial I$ 的外测度为 $0$.$\therefore\operatorname{int}(I),\partial I$ 可测.}
\handline{72}{519}{485}{故 $I$ 可测.$m(I)=m^*(I)=V(I)$.}
\handline[AnnotationBlue]{72}{548}{465}{下面要证明闭集是可测的.}
\handline{72}{585}{850}{Lemma.\quad 若 $\{I_k\}_{k=1}^N$ 是互不相交的闭区间，则 $ E=\bigcup_{k=1}^N I_k$ 可测.且 $ m(E)=\sum_{k=1}^N V(I_k)$.\redann{（作为习题）}}
\handline{72}{619}{856}{Def.\quad$E_1,E_2\subseteq\R^n$.$d(E_1,E_2):=\inf\{|x-y|\mid x\in E_1,y\in E_2\}$.$\operatorname{diam}(E_1):=\sup\{|x-y|\mid x,y\in E_1\}$}
\handline{72}{654}{685}{Thm.\quad 若 $d(E_1,E_2)>0$，则 $m^*(E_1\cup E_2)=m^*(E_1)+m^*(E_2)$.}
\handline{72}{686}{790}{Pf.\quad 显然 $m^*(E_1\cup E_2)\le m^*(E_1)+m^*(E_2)$.下证 $m^*(E_1)+m^*(E_2)\le m^*(E_1\cup E_2)$.}
\handline{72}{719}{710}{$\forall\varepsilon>0,\exists\{I_k\}_{k=1}^\infty\ \text{s.t.}\ E_1\cup E_2\subseteq\bigcup_{k=1}^\infty I_k$.$\sum_{k=1}^\infty V(I_k)\le m^*(E_1\cup E_2)+\varepsilon$}
\handline{72}{752}{805}{不妨设 $\forall k\in\N$.$\operatorname{diam}(I_k)<\dfrac{d(E_1,E_2)}2$，于是 $\exists\{I_k'\},\{I_k''\}$.$\text{s.t.}\ \{I_k\}=\{I_k'\}\cup\{I_k''\}$.}
\handline{72}{787}{455}{且 $ E_1\subseteq\bigcup_{k=1}^\infty I_k'$，$ E_2\subseteq\bigcup_{k=1}^\infty I_k''$.}
\handline{72}{821}{785}{$\therefore m^*(E_1)+m^*(E_2)\le\sum_{k=1}^\infty V(I_k')+\sum_{k=1}^\infty V(I_k'')=\sum_{k=1}^\infty V(I_k)\le m^*(E_1\cup E_2)+\varepsilon$}
\handline{72}{852}{645}{令 $\varepsilon\to0+\Rightarrow m^*(E_1)+m^*(E_2)\le m^*(E_1\cup E_2)$}
\handline{72}{884}{540}{$\therefore m^*(E_1\cup E_2)=m^*(E_1)+m^*(E_2)$}
\handline{72}{920}{856}{Cor.\quad$E_1,E_2$：不相交且均为紧集.则 $d(E_1,E_2)>0$，从而 $m^*(E_1\cup E_2)=m^*(E_1)+m^*(E_2)$}
\end{handpage}

% SOURCE PAGE 014
\begin{handpage}{14}
\handline{72}{123}{395}{Thm.\quad 每个闭集都是可测集.}
\handline{72}{149}{365}{Pf.\quad 设 $E\subseteq\R^n$，$E$ 是闭集.}
\handline{72}{176}{355}{$1^\circ$ 若 $E$ 有界，则 $E$ 是紧集.}
\handline{72}{207}{500}{$\forall\varepsilon>0$，$\exists$ 开集 $G$.$\text{s.t.}\ E\subseteq G$，$m(G)\le m^*(E)+\varepsilon$}
\handline{72}{237}{520}{$G\setminus E$ 是开集 $\therefore\exists\{I_k\}$ 互不相交.$\text{s.t.}\ G\setminus E=\bigcup_{k=1}^\infty I_k$}
\handline{72}{268}{455}{对 $\forall N\in\N$，$\bigcup_{k=1}^N I_k$，$E$ 都是紧集，且不交}
\handline{72}{299}{600}{$\therefore m\left(\bigcup_{k=1}^N I_k\right)+m^*(E)=m^*\left(E\cup\left(\bigcup_{k=1}^N I_k\right)\right)\le m(G)\le m^*(E)+\varepsilon$}
\handline{72}{330}{705}{$\therefore m\left(\bigcup_{k=1}^N I_k\right)\le\varepsilon$，令 $N\to\infty\Rightarrow m^*(G\setminus E)=m\left(\bigcup_{k=1}^\infty I_k\right)\le\varepsilon\quad\therefore E$ 可测}
\handline{72}{361}{665}{$2^\circ$ 若 $E$ 无界，令 $\overline B_k=\{|x|\le k\}$ 则 $E_k=E\cap\overline B_k$ 有界闭.$\therefore E_k$ 可测}
\handline{72}{391}{345}{$\therefore E=\bigcup_{k=1}^\infty E_k$ 可测.}
\handline{72}{423}{400}{Thm.\quad$E$ 是可测集，则 $E^c$ 也可测}
\handline{72}{454}{590}{Pf.\quad$\forall k\in\N$.$\exists$ 开集 $G_k$.$\text{s.t.}\ E\subseteq G_k$.$m^*(G_k\setminus E)<\dfrac1k$.}
\handline{72}{486}{520}{令 $F_k=G_k^c\subseteq E^c$.$ F=\bigcup_{k=1}^\infty F_k\subseteq E^c$，$F$ 是闭集}
\handline{72}{518}{560}{令 $ z=E^c\setminus F=E^c\setminus\left(\bigcup_{k=1}^\infty F_k\right)\subseteq E^c\setminus F_k=G_k\setminus E$}
\handline{72}{551}{640}{$\therefore m^*(z)\le m^*(G_k\setminus E)<\dfrac1k$，$k\to\infty\Rightarrow m^*(z)=0\quad\therefore z$ 可测且零测}
\handline{72}{582}{370}{$\therefore E^c=F\cup z$ 是可测集.}
\handline{72}{615}{295}{Cor.\quad$E_k$ 可测，则}
\handline{72}{643}{285}{\circled{1} $E_1\setminus E_2$ 可测}
\handline{72}{675}{465}{\circled{2} 若 $E_1\subseteq E_2$，则 $m(E_2)\le m(E_1)+m(E_2\setminus E_1)$}
\handline{72}{708}{470}{当 $m(E_1)<\infty$ 时，$m(E_2)-m(E_1)\le m(E_2\setminus E_1)$}
\handline{72}{747}{620}{\circled{3} $ E=\bigcap_{k=1}^\infty E_k$ 可测\quad\ann{$\left(E^c=\bigcup_{k=1}^\infty E_k^c\text{ 可测}\Rightarrow E\text{ 可测}\right)$}}
\handline[AnnotationBlue]{72}{780}{650}{从可测集之间的运算，可抽象出 $\sigma$ 代数的概念.}
\handline{72}{813}{560}{Def.\quad$X$：集合.$\Sigma\subseteq\mathcal P(X)$ 且 $\Sigma\ne\varnothing$}
\handline{72}{845}{610}{若满足 \circled{1} $E\in\Sigma\Rightarrow E^c\in\Sigma$ \circled{2} $E_k\in\Sigma\Rightarrow\bigcup_{k=1}^\infty E_k\in\Sigma$}
\handline{72}{875}{460}{则称 $\Sigma$ 是 $X$ 的一个 $\sigma$-代数}
\handline[NoteRed]{72}{910}{725}{特别地，$\mathcal M=\{\text{全体 Lebesgue 可测集}\}$ 构成 $\R^n$ 的一个 $\sigma$-代数.}
\end{handpage}

% SOURCE PAGE 015
\begin{handpage}{15}
\handline{72}{139}{775}{\underline{Def.} $\{E_k\}_{k=1}^{\infty},\quad\limsup E_k:=\bigcap_{j=1}^{\infty}\bigcup_{k\ge j}^{\infty}E_k,\quad\liminf E_k:=\bigcup_{j=1}^{\infty}\bigcap_{k\ge j}^{\infty}E_k$}
\handline[AnnotationBlue]{72}{164}{660}{显然, 若 $E_k\in\mathcal M\ (\forall k\in\N)$, 则 $\limsup E_k,\ \liminf E_k\in\mathcal M$}
\handline{72}{195}{775}{\underline{Ex.} $\Lambda$ 是一个指标集. 令 $\mathcal F=\{\Sigma_\alpha\mid\Sigma_\alpha\text{ 是 }X\text{ 的 }\sigma\text{-代数},\ \alpha\in\Lambda\},\ \mathcal E=\bigcap_{\alpha\in\Lambda}\Sigma_\alpha$}
\handline{72}{223}{620}{则 $\mathcal F,\ \mathcal E$ 都是 $\sigma$-代数.}
\handline{72}{254}{750}{\underline{Def.} \circled{1} 给定集合 $S$. 考虑包含 $S$ 的最小 $\sigma$-代数 $\sigma(S)$}
\handline{72}{281}{225}{$\sigma(S):=\bigcap_{S\in\Sigma}\Sigma$}
\handline[AnnotationBlue]{269}{274}{430}{或称为 $S$ 生成的 $\sigma$-代数}
\handline{72}{309}{750}{\circled{2} 当 $S=\{\R^n\text{ 所有开集}\}$ 时, $\sigma(S)$ 称为 Borel $\sigma$-代数, 记为 $\mathcal B$.}
\handline{72}{338}{550}{$\mathcal B$ 中的元素称为 Borel 集.}
\handline[Teal]{72}{367}{620}{显然, 所有开集, 闭集都是 Borel 集.}
\handline[Teal]{72}{397}{630}{$G_\delta,\ F_\sigma,\ G_{\delta\sigma},\ F_{\sigma\delta},\ F_{\sigma\delta\sigma}\ \cdots$ 都是 Borel 集.}
\handline[AnnotationBlue]{72}{425}{770}{Borel 集就是从开集出发, 通过可数并和补集运算生成的集合}
\handline[AnnotationBlue]{72}{457}{715}{$G_\delta, F_\sigma\cdots$ 都是 Borel 集, 但并不是所有 Borel 集都长这样}
\handline[AnnotationBlue]{72}{487}{705}{Borel $\sigma$-代数足够大, 且元素足够“规范”}
\handline[AnnotationBlue]{72}{516}{650}{包含了实际分析中的大多数集合}
\handline[AnnotationBlue]{72}{544}{750}{并且排除了需要选择公理才能构造的病态: 不可测集}
\end{handpage}

% SOURCE PAGE 016
% SECTION: 3.3 | Two properties of Lebesgue Measure
\begin{handpage}{16}
\handtitle{72}{111}{735}{$\S2.3\quad$ Two properties of Lebesgue Measure}
\handline{72}{139}{765}{\underline{Lemma} 若 $E\in\mathcal M$, 则 $\forall\varepsilon>0$, $\exists$ 闭集 $F$, s.t. $F\subseteq E$, 且 $m(E\setminus F)<\varepsilon$}
\handline{72}{166}{745}{Pf. $E^c\in\mathcal M$. $\therefore\forall\varepsilon>0$, $\exists$ 开集 $G$, s.t. $E^c\subseteq G$. 且 $m(G\setminus E^c)<\varepsilon$}
\handline{72}{192}{690}{令 $F=G^c$. 则 $F\subseteq E$, 而 $E\setminus F=E\cap F^c=E\cap G=G\setminus E^c$}
\handline{72}{216}{610}{$\therefore m(E\setminus F)=m(G\setminus E^c)<\varepsilon$}
\handline{72}{244}{700}{\underline{Thm.} \ann{（可数可加性）}}
\handline{72}{270}{730}{$\{E_k\}_{k=1}^{\infty}$ 两两不交且可测, 则 $\bigcup_{k=1}^{\infty}E_k\in\mathcal M$ 且 $m\left(\bigcup_{k=1}^{\infty}E_k\right)=\sum_{k=1}^{\infty}m(E_k)$.}
\handline{72}{297}{660}{Pf. 不妨设 $\forall k\in\N$, $E_k$ 有界}
\handline{72}{325}{710}{$\forall\varepsilon>0$, $\exists$ 闭集 $F_k$, s.t. $F_k\subseteq E_k$, $m(E_k\setminus F_k)<\dfrac\varepsilon{2^k}$}
\handline{72}{350}{650}{$F_k$ 两两不交且是紧集}
\handline{72}{380}{665}{$\therefore\sum_{k=1}^{N}m(F_k)=m\left(\bigcup_{k=1}^{N}F_k\right)\le m\left(\bigcup_{k=1}^{\infty}F_k\right)$}
\handline{72}{409}{680}{令 $N\to\infty\Rightarrow\sum_{k=1}^{\infty}m(F_k)\le m\left(\bigcup_{k=1}^{\infty}E_k\right)\quad\ann{\circled{1}}$}
\handline{72}{441}{745}{另一方面 $E_k=(E_k\setminus F_k)\cup F_k\quad\therefore m(E_k)\le m(E_k\setminus F_k)+m(F_k)<\dfrac\varepsilon{2^k}+m(F_k)$}
\handline{72}{471}{695}{$\Rightarrow\sum_{k=1}^{\infty}m(E_k)<\sum_{k=1}^{\infty}\left(\dfrac\varepsilon{2^k}+m(F_k)\right)=\sum_{k=1}^{\infty}m(F_k)+\varepsilon$}
\handline{72}{501}{660}{$\therefore\sum_{k=1}^{\infty}m(E_k)\le\sum_{k=1}^{\infty}m(F_k).\quad\ann{\circled{2}}$}
\handline{72}{531}{690}{由 \ann{\circled{1}\circled{2}} $\Rightarrow\sum_{k=1}^{\infty}m(E_k)=m\left(\bigcup_{k=1}^{\infty}E_k\right)$}
\handline{72}{561}{785}{\underline{Cor.} $\{I_k\}_{k=1}^{\infty}$, 两两不交, 则 $m\left(\bigcup_{k=1}^{\infty}I_k\right)=\sum_{k=1}^{\infty}m(I_k)$.}
\handline{72}{592}{790}{\underline{Cor.} $E_1\subseteq E_2$, $E_1,E_2\in\mathcal M$. $m(E_1)<\infty$, 则 $m(E_2\setminus E_1)=m(E_2)-m(E_1)$}
\handline{72}{623}{750}{\underline{Thm.} $E_k\in\mathcal M$, $k\in\N$ \ann{（测度的上下连续性）}}
\handline{72}{652}{745}{\circled{1} $E_k\uparrow E$ 则 $E\in\mathcal M$, $m(E)=\lim_{k\to\infty}m(E_k)$}
\handline{72}{683}{750}{\circled{2} $E_k\downarrow E$, $m(E_1)<\infty$ 则 $E\in\mathcal M$, $m(E)=\lim_{k\to\infty}m(E_k)$}
\handline[NoteRed]{72}{715}{750}{\circled{2} 中若没有 $m(E_1)<\infty$, 则有反例: $E_k=[k,+\infty)\downarrow\varnothing$, 但 $m(E_k)=\infty$, $m(\varnothing)=0$.}
\handline{72}{752}{800}{\underline{Thm.} $E_k\uparrow E$, 则 $m^*(E)=\lim_{k\to\infty}m^*(E_k)$ \ann{（外测度的上连续性）}}
\handline{72}{781}{805}{Pf. $\forall k\in\N$, $\exists G_\delta$ 集 $H_k$, s.t. $E_k\subseteq H_k$, $m^*(E_k)=m^*(H_k)$}
\handline{72}{813}{770}{令 $V_m=\bigcap_{j=m}^{\infty}H_j$, 则 $V_m\uparrow V=\bigcup_{m=1}^{\infty}V_m$ 且 $V_m,V\in\mathcal M$. $m(V)=\lim_{m\to\infty}m(V_m)$}
\handline{72}{843}{793}{$\forall j\ge m$, $E_m\subseteq E_j\subseteq H_j\quad\therefore E_m\subseteq V_m\subseteq H_m\quad\therefore m^*(E_m)\le m(V_m)\le m(H_m)=m^*(E_m)$}
\handline{72}{876}{735}{$\therefore m^*(E_m)=m(V_m)\Rightarrow m(V)=\lim_{m\to\infty}m^*(E_m)$}
\handline{72}{899}{750}{而 $E=\bigcup_{k=1}^{\infty}E_k\subseteq\bigcup_{k=1}^{\infty}V_k=V\quad\therefore m^*(E)\le m(V)=\lim_{m\to\infty}m^*(E_m)$}
\handline{72}{925}{640}{而 $E_m\subseteq E\Rightarrow m^*(E_m)\le m^*(E)\Rightarrow\lim_{m\to\infty}m^*(E_m)\le m^*(E)$}
\handline{670}{890}{220}{$\left.\rule{0pt}{7mm}\right\}\Rightarrow m^*(E)=\lim_{m\to\infty}m^*(E_m)$}
\end{handpage}

% SOURCE PAGE 017
% SECTION: 3.4 | Characterizations of Measurability
\begin{handpage}{17}
\handtitle{72}{136}{790}{$\S2.4\quad$ Characterizations of Measurability}
\handline{72}{169}{770}{\underline{Thm.} (i) $E\in\mathcal M\Longleftrightarrow\exists G_\delta$ 集 $H$, s.t. $E\subseteq H$, 且 $m(H\setminus E)=0$}
\handline{72}{194}{730}{(ii) $E\in\mathcal M\Longleftrightarrow\exists F_\sigma$ 集 $H$, s.t. $H\subseteq E$, 且 $m(E\setminus H)=0$}
\handline{72}{225}{560}{\underline{Thm.} 设 $m^*(E)<\infty$}
\handline{72}{247}{785}{则 $E\in\mathcal M\Longleftrightarrow\forall\varepsilon>0,\ \exists S=\bigcup_{k=1}^{N}I_k$, $I_k$ 互不相交, s.t. $E=(S\cup N_1)\setminus N_2$, $m^*(N_1),m^*(N_2)<\varepsilon$.}
\handline{72}{277}{600}{Pf. “$\Rightarrow$” 若 $E\in\mathcal M$.}
\handline{72}{305}{700}{$\forall\varepsilon>0$, $\exists$ 开集 $G$, s.t. $E\subseteq G$, $m(G\setminus E)<\varepsilon$.}
\handline{72}{332}{670}{设 $G=\bigcup_{k=1}^{\infty}I_k$. 而 $m(G)\le m(G\setminus E)+m(E)<\infty$}
\handline{72}{357}{650}{$\therefore\quad m(G)=\sum_{k=1}^{\infty}V(I_k)$ 收敛.}
\handline{72}{387}{710}{$\therefore\exists N\in\N$, s.t. $n>N$ 时 $\sum_{k=N+1}^{\infty}V(I_k)<\varepsilon$.}
\handline{72}{416}{680}{令 $S=\bigcup_{k=1}^{N}I_k.\quad N_1=E\setminus S,\quad N_2=S\setminus E$.}
\handline{72}{446}{640}{则 $E=(S\cup N_1)\setminus N_2$.}
\handline{72}{475}{680}{且 $m^*(N_1)=m^*(E\setminus S)\le m^*(G\setminus S)=\sum_{k=N+1}^{\infty}V(I_k)<\varepsilon$}
\handline{72}{505}{660}{$m^*(N_2)=m^*(S\setminus E)\le m(G\setminus E)<\varepsilon$}
\handline{72}{534}{755}{“$\Leftarrow$” $\forall\varepsilon>0$, $\exists S,N_1,N_2.\quad E=(S\cup N_1)\setminus N_2.\quad m^*(N_1),m^*(N_2)<\dfrac\varepsilon3$}
\handline{72}{567}{680}{由于 $S$ 为有限闭区间之并 $\therefore\exists$ 开集 $V$, s.t. $S\subseteq V.\ m(V\setminus S)<\dfrac\varepsilon3$.}
\handline{72}{599}{650}{由外测度定义, $\exists$ 开集 $W$, s.t. $N_1\subseteq W$, $m(W)<\dfrac\varepsilon3$.}
\handline{72}{630}{670}{令 $U=V\cup W$. 则 $E\subseteq U$. 且 $U$ 为开集}
\handline{72}{663}{710}{$U\setminus E=(V\cup W)\setminus E=(V\setminus E)\cup(W\setminus E)\subseteq(V\setminus E)\cup W\subseteq(V\setminus S)\cup(S\setminus E)\cup W$}
\handline{72}{695}{725}{$\therefore m^*(U\setminus E)\le m^*(V\setminus S)+m^*(N_2)+m^*(W)<\varepsilon.\qquad\subseteq(V\setminus S)\cup N_2\cup W$.}
\handline{72}{735}{695}{取 $\varepsilon_n=\dfrac1n>0$. $\exists U_n$ 开集 s.t. $E\subseteq U_n.\ m^*(U_n\setminus E)<\dfrac1n$}
\handline{72}{767}{650}{令 $G=\bigcap_{n=1}^{\infty}U_n$ 为 $G_\delta$ 集. 且 $E\subseteq G$.}
\handline{72}{798}{650}{$m^*(G\setminus E)\le m^*(U_n\setminus E)<\dfrac1n\to0\ (n\to\infty)$}
\handline{72}{834}{560}{$\therefore m(G\setminus E)=0$}
\handline{72}{860}{460}{$\therefore E\in\mathcal M$.}
\end{handpage}

% SOURCE PAGE 018
\begin{handpage}{18}
\handline{72}{120}{680}{\underline{Thm.} \ann{(Carathéodory)}}
\handline{72}{150}{750}{$E\subseteq\R^n$, $E\in\mathcal M\Longleftrightarrow\forall A\subseteq\R^n$, $m^*(A)=m^*(A\cap E)+m^*(A\setminus E)$}
\handline{72}{177}{745}{Pf. “$\Rightarrow$” $\forall A\subseteq\R^n$. $\exists G_\delta$ 集 $H$, s.t. $A\subseteq H$, $m^*(A)=m(H)$.}
\handline{72}{204}{650}{$H=(H\cap E)\cup(H\setminus E)\Rightarrow m(H)=m(H\cap E)+m(H\setminus E)$}
\handline{72}{232}{675}{$\therefore m^*(A)=m(H)=m(H\cap E)+m(H\setminus E)\ge m^*(A\cap E)+m^*(A\setminus E)$}
\handline{72}{260}{660}{而 $A=(A\cap E)\cup(A\setminus E)\Rightarrow m^*(A)\le m^*(A\cap E)+m^*(A\setminus E)$}
\handline{72}{286}{630}{$\therefore m^*(A)=m^*(A\cap E)+m^*(A\setminus E)$}
\handline{72}{317}{640}{“$\Leftarrow$” 若 $m^*(E)<\infty$.}
\handline{72}{350}{630}{$\exists G_\delta$ 集 $H$. s.t. $E\subseteq H$. $m(H)=m^*(E)$}
\handline{72}{379}{670}{$\therefore m(H)=m^*(H\cap E)+m^*(H\setminus E)=m^*(E)+m^*(H\setminus E)$}
\handline{72}{408}{630}{$\therefore m^*(H\setminus E)=0.\quad\therefore H\setminus E\in\mathcal M$.}
\handline{72}{438}{600}{$\therefore E=H\setminus(H\setminus E)\in\mathcal M$.}
\handline{72}{469}{670}{一般地. 令 $\overline B_k=\{|x|\le k\}$. $E_k=E\cap\overline B_k$. 则 $m^*(E_k)<\infty$, $E=\bigcup_{k=1}^{\infty}E_k$}
\handline{72}{500}{680}{$\forall k\in\N$, $\exists G_\delta$ 集 $H_k$, s.t. $E_k\subseteq H_k$. $m(H_k)=m^*(E_k)$}
\handline{72}{531}{670}{$m(H_k)=m^*(H_k\cap E)+m^*(H_k\setminus E)\ge m^*(E_k)+m^*(H_k\setminus E)$}
\handline{72}{563}{640}{$\therefore m^*(H_k\setminus E)=0.\quad\therefore H_k\setminus E\in\mathcal M$.}
\handline{72}{596}{680}{$\therefore H\setminus E=\bigcup_{k=1}^{\infty}(H_k\setminus E)\in\mathcal M\quad\therefore E=H\setminus(H\setminus E)\in\mathcal M$.}
\handline{72}{632}{765}{\underline{Cor.} $A\subseteq\R^n$, $E\subseteq A$, $E\in\mathcal M$. 则 $m^*(A)=m(E)+m^*(A\setminus E)$}
\handline{72}{669}{690}{若 $m(E)<\infty$, 则 $m^*(A\setminus E)=m^*(A)-m(E)$.}
\end{handpage}

% SOURCE PAGE 019
\begin{handpage}{19}
\handline{72}{128}{785}{\underline{Thm.} $E\subseteq\R^n$, $\exists G_\delta$ 集 $H$, s.t. $E\subseteq H$, 且 $\forall M\in\mathcal M$, $m^*(E\cap M)=m(H\cap M)$}
\handline{72}{158}{770}{Pf. 若 $m^*(E)<\infty$, 则 $\exists G_\delta$ 集 $H$, s.t. $E\subseteq H$. $m(H)=m^*(E)$.}
\handline{72}{184}{720}{$\forall M\in\mathcal M.\quad m^*(E)=m^*(E\cap M)+m^*(E\setminus M)$}
\handline{72}{212}{620}{$m(H)=m(H\cap M)+m(H\setminus M)$}
\handnote[black]{81}{199}{45}{$\parallel$}
\handline{72}{239}{680}{但 $m^*(E\cap M)\le m(H\cap M)$, $m^*(E\setminus M)\le m(H\setminus M)$}
\handline{72}{266}{640}{故 $m^*(E\cap M)=m(H\cap M)$}
\handline{72}{295}{730}{一般地, 令 $\overline B_k=\{|x|\le k\}$, $E_k=E\cap\overline B_k$, 则 $E=\bigcup_{k=1}^{\infty}E_k$.}
\handline{72}{324}{770}{$\forall k\in\N$. $\exists G_\delta$ 集 $U_k$, s.t. $E_k\subseteq U_k$ 且 $m^*(E_k\cap M)=m(U_k\cap M)\ \ann{(*)}$}
\handline{72}{354}{765}{令 $H_m=\bigcap_{k\ge m}U_k$, 则 $H_m\uparrow H=\bigcup_{m=1}^{\infty}H_m$, $H$ 是一个 $G_{\delta\sigma}$ 集 $\therefore H\in\mathcal M$.}
\handline{72}{384}{705}{由于 $E_k$ 是升列. $\therefore\forall k\ge m$, $E_m\subseteq E_k\subseteq U_k$}
\handline{72}{411}{650}{$\therefore E_m\subseteq H_m\subseteq U_m$}
\handline{72}{441}{755}{$\therefore m^*(E_m\cap M)\le m(H_m\cap M)\le m(U_m\cap M)\overset{\ann{(*)}}=m^*(E_m\cap M)$}
\handline{72}{470}{710}{$\therefore m^*(E_m\cap M)=m(H_m\cap M)$}
\handline{72}{501}{690}{$E_m\cap M\uparrow E\cap M,\quad H_m\cap M\uparrow H\cap M$}
\handline{72}{530}{735}{由外测度的上连续性 $\Rightarrow m^*(E\cap M)=m(H\cap M)$}
\handline[NoteRed]{72}{562}{730}{最后由于 $H$ 不是 $G_\delta$ 集, 是 $G_{\delta\sigma}$ 集, 还需要一步调整}
\handline{72}{592}{760}{$H\in\mathcal M$, $\therefore\exists G_\delta$ 集 $H_1$, s.t. $H\subseteq H_1$, $m(H)=m(H_1)$, $m(H_1\setminus H)=0$.}
\handline{72}{624}{800}{$\therefore m(H_1\cap M)=m(H\cap M)+m((H_1\setminus H)\cap M)=m(H+M)+0=m^*(E\cap M)$}
% 末行后一个 m(H+M) 的运算符原稿明确写成加号，照录.
\end{handpage}

% SOURCE PAGE 020
% SECTION: 3.5 | Lipschitz Transformation of R^n
\begin{handpage}{20}
\handtitle{72}{107}{760}{$\S2.5\quad$ Lipschitz Transformation of $\R^n$}
\handline{72}{138}{800}{\underline{Def.} $f:\R^n\to\R^n$, 若 $\exists L>0$, $\forall x,y\in\R^n$, $|f(x)-f(y)|\le L|x-y|$}
\handline{72}{166}{720}{则称 $f$ 是 Lipschitz 的, $L$ 称为 Lipschitz 常数}
\handline{72}{197}{720}{\underline{Thm.} $f:\R^n\to\R^n:$ Lipschitz.}
\handline{72}{222}{700}{若 $M\in\mathcal M$, 则 $f(M)\in\mathcal M$.}
\handline{72}{249}{790}{Pf. $M\in\mathcal M\quad\therefore\exists F_\sigma$ 集 $F$, s.t. $F\subseteq M$, $M=F\cup z$ 且 $m(z)=0$.}
\handline{72}{278}{730}{由 $f$ Lipschitz $\Rightarrow f$ 连续 $\Rightarrow f(F)$ 是 $F_\sigma$ 集, $\therefore f(F)\in\mathcal M$.}
\handline{72}{307}{730}{$m(z)=0\Rightarrow\forall\varepsilon>0$, $\exists\{I_k\}_{k=1}^{\infty}$, s.t. $z\subseteq\bigcup_{k=1}^{\infty}I_k$ 且 $\sum_{k=1}^{\infty}V(I_k)<\varepsilon$.}
\handline{72}{336}{720}{$\forall k\in\N$, 由 $f$ Lipschitz $\Rightarrow m^*(f(I))\le L'\,m(I)$.}
\draw[Teal,dashed,line width=.5pt] (189,343) rectangle (951,542);
\handline{72}{363}{680}{设 $I=\{x=(x_1,x_2,\cdots,x_n)\mid a_i\le x_i\le b_i\}$}
\handline{72}{391}{700}{$f(x)=(f_1(x),\cdots,f_n(x))$. 则 $f_i:\R^n\to\R$ 且连续.}
\handline{72}{421}{715}{故 $f_i(I)$ 有最大值, 最小值 $c_i\le f_i(x)\le d_i$. 设 $f_i(x_{i1})=c_i$, $f_i(x_{i2})=d_i$.}
\handline{72}{450}{715}{则 $d_i-c_i=|f_i(x_{i1})-f_i(x_{i2})|\le\alpha_i|f(x_{i1})-f(x_{i2})|\le\alpha_i L|x_{i1}-x_{i2}|$}
\handline{72}{474}{400}{$\le\alpha_i\cdot L\cdot\beta_i(b_i-a_i)$.}
\handline{72}{504}{730}{$\therefore m^*(f(I))\le\prod_{i=1}^n(d_i-c_i)\le L^n\prod_{i=1}^n\alpha_i\beta_i(b_i-a_i)\quad$令 $L'=L^n\prod_{i=1}^n\alpha_i\beta_i$ 即可}
\handline[AnnotationBlue]{72}{534}{660}{$\alpha_i,\beta_i$ 是由范数等价性得到的常数}
\handline{72}{591}{730}{$\therefore m^*(f(z))\le\sum_{k=1}^{\infty}m^*(f(I_k))\le\sum_{k=1}^{\infty}L'V(I_k)<L'\varepsilon$}
\handline{72}{627}{660}{$\therefore m^*(f(z))=0.\quad f(z)\in\mathcal M$.}
\handline{72}{658}{660}{$\therefore f(M)=f(F)\cup f(z)\in\mathcal M$.}
\handline{72}{689}{795}{\underline{Thm.} $T$ 是 $\R^n$ 上的线性变换, 表示阵为 $A=(a_{ij})_{n\times n}$ \ann{（$T$ 显然 Lipschitz）}}
\handline{72}{720}{700}{则 $\forall E\in\mathcal M$. $m(T(E))=|A|m(E)$.}
\handline{72}{751}{790}{Pf. $\forall\varepsilon>0$, $\exists\{I_k\}_{k=1}^{\infty}$ s.t. $E\subseteq\bigcup_{k=1}^{\infty}I_k$, $\sum_{k=1}^{\infty}V(I_k)\le m(E)+\varepsilon$.}
\handline{72}{783}{790}{$\therefore T(E)\subseteq\bigcup_{k=1}^{\infty}T(I_k)\quad\therefore m(T(E))\le\sum_{k=1}^{\infty}m(T(I_k))=|A|\sum_{k=1}^{\infty}V(I_k)\le|A|(m(E)+\varepsilon)$}
\handline{72}{813}{660}{$\therefore m(T(E))\le|A|m(E)$.}
\handline{72}{848}{690}{若 $|A|=0$ 则 $m(T(E))=0=|A|m(E)$}
\handline{72}{884}{740}{若 $|A|\ne0$ 则 $T$ 可逆, 同理 $\Rightarrow m(E)\le\dfrac1{|A|}m(T(E))$}
\handline{72}{918}{420}{$\therefore m(T(E))=|A|m(E)$.}
\end{handpage}

% SOURCE PAGE 021
% SECTION: 3.6 | A Nonmeasurable Set
\begin{handpage}{21}
\handtitle{72}{132}{750}{$\S2.6\quad$ A Nonmeasurable Set}
\handline{72}{161}{720}{Zermelo's Axiom \ann{（选择公设）}}
\handline{72}{182}{770}{设 $\mathcal F$ 是两两不交的非空集 $E_\alpha$ 构成的集合族, $\alpha\in\Lambda$}
\handline{72}{209}{730}{则存在一个集合 $C$. s.t. $\forall\alpha\in\Lambda.\quad\#(C\cap E_\alpha)=1$}
\handline[AnnotationBlue]{72}{237}{810}{（等价地, $\exists$ 一个函数 $f:\mathcal F\to\bigcup\mathcal F$. s.t. $\forall\alpha\in\Lambda.\ f(E_\alpha)\in E_\alpha$.}
\handline{72}{265}{810}{\ann{$f$ 称为选择函数）} $\bigcup\mathcal F=\{x\mid\exists E_\alpha\in\mathcal F,\ x\in E_\alpha\}$, 即把 $\mathcal F$ 中所有集合中的元素并起来}
\handline{72}{295}{740}{\underline{Thm.} $E\subseteq\R$, $E\in\mathcal M$, $m(E)>0$}
\handline{72}{320}{750}{则差分集 $\{x-y\mid x,y\in E\}$ 包含了一个以 $0$ 为中心的区间}
\handline{72}{350}{775}{Pf. $\forall\varepsilon>0$, $\exists$ 开集 $G$, s.t. $E\subseteq G$, $m(G)\le(1+\varepsilon)m(E)$}
\handline{72}{377}{700}{令 $G=\bigcup_{k=1}^{\infty}I_k$, $I_k$ 互不重叠}
\handline{72}{406}{745}{令 $E_k=E\cap I_k$, 则 $E=\bigcup_{k=1}^{\infty}E_k$. $E_k$ 互不重叠.}
\handline{72}{437}{735}{$\therefore\sum_{k=1}^{\infty}V(I_k)\le(1+\varepsilon)m(E)=(1+\varepsilon)\sum_{k=1}^{\infty}m(E_k)$}
\handline{72}{466}{735}{$\Rightarrow\exists k_0$, s.t. $V(I_{k_0})\le(1+\varepsilon)m(E_{k_0})\ \redann{(*)}$}
\handline{72}{499}{745}{取 $\varepsilon=\dfrac13$, 则 $V(I_{k_0})\le\dfrac43m(E_{k_0})$, 考虑 $|d|<\dfrac12V(I_{k_0})\quad d\in\R$}
\handline{72}{531.961}{750}{断言: $\left(-\dfrac12V(I_{k_0}),\dfrac12V(I_{k_0})\right)\subseteq\{x-y\mid x,y\in E\}$.}
\handline{72}{567.332}{740}{$\forall|d|<\dfrac12V(I_{k_0})$. 令 $E_d=\{x+d\mid x\in E_{k_0}\}$ 则 $m(E_{k_0})=m(E_d)$}
\handline{72}{588}{760}{若 $E_{k_0}\cap E_d=\varnothing$, 则 $m(E_{k_0}\cup E_d)=m(E_{k_0})+m(E_d)=2m(E_{k_0})$}
\handline{72}{623}{745}{但 $E_{k_0}\cup E_d\subseteq I_{k_0}\cup I_d$, 其中 $I_d=\{x+d\mid x\in I_{k_0}\}$.}
\handdiagram{780}{631}{148}{62}{
  \draw[line width=.7pt] (0,0)--(4,0);
  \draw[Teal] (0,-.12)--(0,.55) (3.35,-.12)--(3.35,.55) (3.85,-.12)--(3.85,.55);
  \draw[NoteRed] (.5,-.18)--(.5,.15) (3.85,-.18)--(3.85,.15);
  \draw[AnnotationBlue] (0,.58)--(.15,.66)--(.3,.58)--(1.65,.58)--(1.75,.76)--(1.9,.58)--(3.2,.58)--(3.35,.45);
  \node[AnnotationBlue,above,font=\small] at (.25,.65) {$I_{k_0}$};
  \draw[NoteRed] (.5,-.3)..controls(.7,-.15)and(.9,-.7)..(2.0,-.75)..controls(2.8,-.75)and(3.5,-.3)..(3.85,-.3);
  \node[NoteRed,below,font=\small] at (2,-.7) {$I_d$};
  \draw[Teal] (0,-.3)--(.2,-.46)--(.5,-.3) (3.35,.7)--(3.6,.85)--(3.85,.7);
  \node[Teal,below,font=\small] at (.2,-.46) {$d$};
  \node[Teal,above,font=\small] at (3.6,.85) {$d$};
}
\handline{72}{663}{570}{$\therefore m(E_{k_0}\cup E_d)\le m(I_{k_0}\cup I_d)\le V(I_{k_0})+|d|$}
\handline{72}{694}{580}{$\therefore2m(E_{k_0})\le V(I_{k_0})+|d|<\dfrac32V(I_{k_0})$}
\handline{72}{725}{720}{$\therefore V(I_{k_0})>\dfrac43m(E_{k_0})$, 与 \redann{$(*)$} 矛盾.}
\handline{72}{759}{750}{$\therefore\exists x\in E_{k_0}$, $y\in E_{k_0}$, s.t. $x=y+d\quad\therefore d=x-y\in\{x-y\mid x,y\in E\}$.}
\handline{72}{790}{740}{$\therefore\left(-\dfrac12V(I_{k_0}),\dfrac12V(I_{k_0})\right)\subseteq\{x-y\mid x,y\in E\}$.}
\end{handpage}

% SOURCE PAGE 022
\begin{handpage}{22}
\handline{72}{125}{170}{\textbf{Thm.}\quad\ann{(Vitali)}}
\handline{72}{154}{270}{存在 Lebesgue 不可测集.}
\handline{72}{184}{490}{\textbf{Pf.}\quad 在 $\R$ 上定义关系 $\sim$：$x\sim y\iff x-y\in\Q$}
\handline{72}{210}{295}{易验证这是一个等价关系}
\handline{72}{240}{455}{$\therefore\R=\bigcup_{r\in\R\setminus\Q}\Q_r$，其中 $\Q_r=\Q+r$，$r\in\R\setminus\Q$}
\handline{72}{269}{755}{由 Zermelo's Axiom，$\exists$ 集合 $E=\{\Q_r\text{ 中的一个元素}\}$\redann{（从每个等价类中选一个代表元）}}
\handline{72}{301}{410}{令 $E_p:=E+p$，$p\in\Q$，则 $\R=\bigcup_{p\in\Q}E_p$}
\handline{72}{329}{505}{若 $E\in\mathcal M$，则 $E_p=E+p\in\mathcal M$，且 $m(E)=m(E_p)$}
\handline{72}{360}{465}{$1^\circ\ m(E)=0$，则 $m(\R)=\sum_{p\in\Q}m(E_p)=0$，矛盾!}
\handline{72}{391}{590}{$2^\circ\ m(E)>0$，则 $E_p$ 的差分集包含了一个以 $0$ 为中心的区间.}
\handline{72}{422}{600}{但是断言，$\forall d\in\Q\setminus\{0\}$，$d\notin E_p$ 的差分集，从而矛盾!}
\handline{72}{456}{505}{$\forall x\ne y\in E_p$，$\exists x_0,y_0\in E$ s.t. $x=x_0+p,y=y_0+p$}
\handline{72}{489}{405}{$\therefore x_0\ne y_0$，$\therefore x_0,y_0$ 来自不同等价类.}
\handline{72}{520}{285}{$\therefore x-y=x_0-y_0\notin\Q$}
\handline[NoteRed]{553}{485}{160}{$\Rightarrow\forall x\ne y\in E_p$}
\handline[NoteRed]{603}{517}{115}{$x-y\notin\Q$}
\handline{72}{552}{405}{$\therefore E_p$ 的差分集中没有非零有理数.}
\handline{72}{586}{500}{\textbf{Thm.}\quad$F\subseteq\R$，$m^*(F)>0$，则 $F$ 有一个不可测子集.}
\handline{72}{620}{515}{\textbf{Pf.}\quad 令 $E$ 为 \ann{Vitali} 中构造的集合，则 $\R=\bigcup_{p\in\Q}E_p$}
\handline{72}{658}{240}{$\therefore F=\bigcup_{p\in\Q}(F\cap E_p)$}
\handline{72}{692}{485}{若 $\forall p\in\Q$，$F\cap E_p\in\mathcal M$.下证 $m(F\cap E_p)=0$.}
\handline{72}{730}{630}{若 $m(F\cap E_p)>0$，则 $F\cap E_p$ 的差分集包含一个以 $0$ 为中心的区间.}
\handline{72}{764}{635}{而 $F\cap E_p\subseteq E_p=E+p$，由 $E$ 的定义 $\Rightarrow\forall x\ne y\in E$，$x-y\notin\Q$.}
\handline{72}{798}{500}{故 $F\cap E_p$ 不可能包含区间，$\therefore m(F\cap E_p)=0$.}
\handline{72}{834}{590}{$\therefore m^*(F)\leq\sum_{p\in\Q}m(F\cap E_p)=0$.与 $m^*(F)>0$ 矛盾.}
\handline[Teal]{72}{868}{650}{若 $m^*(F)=0$，则 $F$ 任一子集 $G\subseteq F$，$m^*(G)=0$，$\therefore G$ 可测.}
\handline[Teal]{72}{906}{465}{$\therefore m^*(F)=0$ 时，$F$ 没有不可测子集.}
\end{handpage}

% SOURCE PAGE 023
% CHAPTER: 4 | Lebesgue Measurable Functions
% SECTION: 4.1 | Elementary Properties of Measurable Functions
\begin{handpage}{23}
\handtitle{72}{149}{490}{Chapter 3 Lebesgue Measurable Functions}
\handline[AnnotationBlue]{72}{125}{310}{（本章中的 $E$ 默认为可测集）}
\handtitle{72}{177}{590}{§3.1 Elementary Properties of Measurable Functions}
\handline{72}{203}{445}{\textbf{Thm.}\quad$f:E\to\R\cup\{\pm\infty\}$，$E\in\mathcal M$，TFAE：}
\handline{72}{229}{505}{$(i)\ \forall a\in\R$，$\{f>a\}\in\mathcal M$\quad$(ii)\ \forall a\in\R$，$\{f\leq a\}\in\mathcal M$}
\handline{72}{256}{510}{$(iii)\ \forall a\in\R$，$\{f<a\}\in\mathcal M$\quad$(iv)\ \forall a\in\R$，$\{f\geq a\}\in\mathcal M$}
\handline{72}{286}{780}{\textbf{Cor.}\quad 若 $f$ 可测，则 $\{f>-\infty\},\{f<+\infty\},\{f=+\infty\},\{a\leq f\leq b\},\{f=a\}\in\mathcal M$}
\handline{72}{313}{415}{$f$ 可测 $\iff\forall a\in\R$，$\{a<f<+\infty\}\in\mathcal M$.}
\handline{72}{344}{705}{\textbf{Thm.}\quad$f:\R^n\to\R\cup\{\pm\infty\}$，则 $f$ 可测 $\iff\forall$ 开集 $G\subseteq\R$，$f^{-1}(G)\in\mathcal M$.}
\handline[Teal]{72}{374}{470}{若 $f^{-1}(G)$ 是 Borel 集，则称 $f$ 是 Borel 可测的.}
\handline{72}{408}{650}{\textbf{Pf.}\quad“$\Rightarrow$” $G\subseteq\R$，$ G=\bigcup_{i=1}^{\infty}(a_i,b_i)$，$f^{-1}(a_i,b_i)=\{a_i<f<b_i\}$}
\handline{72}{435}{440}{$f$ 可测 $\Rightarrow f^{-1}(a_i,b_i)\in\mathcal M\Rightarrow f^{-1}(G)\in\mathcal M$.}
\handline{72}{463}{645}{“$\Leftarrow$”令 $G=(a,+\infty)$ 则 $f^{-1}(G)=\{a<f<+\infty\}\in\mathcal M$，$\therefore f$ 可测.}
\handline{72}{498}{615}{\textbf{Thm.}\quad 设 $D\subseteq\R$ 稠密，若 $\forall a\in D$，$\{a<f\}\in\mathcal M$，则 $f$ 可测.}
\handline{72}{531}{540}{\textbf{Pf.}\quad$\forall\lambda\in\R$，取 $\{d_k\}\subseteq D$，s.t. $d_k\to\lambda$ 且 $d_k\geq\lambda$.}
\handline{72}{562}{475}{$\therefore\{f>\lambda\}=\bigcup_{k=1}^{\infty}\{f>d_k\}\in\mathcal M$，$\therefore f$ 可测!}
\handline{72}{597}{780}{\textbf{Thm.}\quad$f,g$ 在 $\R^n$ 上可测，$f=g$ a.e.，则 $m(\{f>a\})=m(\{g>a\})$，$\forall a\in\R$.}
\handline{72}{630}{675}{\textbf{Pf.}\quad 令 $Z=\{f\ne g\}$，则 $m(Z)=0$，$\forall a\in\R$，$\{f>a\}\cup Z=\{g>a\}\cup Z$}
\handline{72}{663}{615}{$\therefore m(\{f>a\})=m(\{f>a\}\cup Z)=m(\{g>a\}\cup Z)=m(\{g>a\})$.}
\handline[NoteRed]{72}{696}{710}{可测函数复合不一定可测，但可测函数与连续函数复合是可测函数.}
\handline{72}{731}{650}{\textbf{Thm.}\quad$f$ 可测，几乎处处有限，$g:\R\to\R$ 连续，则 $g\circ f$ 可测.}
\handline{72}{762}{470}{\textbf{Pf.}\quad$f$ 几乎处处有限 $\Rightarrow g\circ f$ 几乎处处有定义.}
\handline{72}{794}{650}{令 $G\subseteq\R$ 为开集，则 $g^{-1}(G)$ 为开集，$\therefore f^{-1}(g^{-1}(G))\in\mathcal M$}
\handline{72}{826}{605}{$\therefore(g\circ f)^{-1}(G)=f^{-1}(g^{-1}(G))\in\mathcal M$，$\therefore g\circ f$ 可测}
\handline[Teal]{72}{860}{856}{特别地，若 $f$ 可测，则 $|f|$，$|f|^p\ (p>0)$，$e^{kf}\ (k\in\R)$，$f^+=\max(0,f)$，$f^-=\max(0,-f)$}
\handline[Teal]{72}{891}{220}{都是可测函数.}
\end{handpage}

% SOURCE PAGE 024
\begin{handpage}{24}
\handline{72}{118}{540}{\textbf{Thm.}\quad$f,g$ 可测，则 $\{f>g\}\in\mathcal M$.}
\handline{72}{142}{570}{\textbf{Pf.}\quad 设 $\{r_k\}$ 是所有有理数.}
\handline{72}{169}{650}{$\Rightarrow\{f>g\}=\bigcup_{k=1}^{\infty}\{f>r_k>g\}=\bigcup_{k=1}^{\infty}(\{f>r_k\}\cap\{g<r_k\})\in\mathcal M$.}
\handline{72}{196}{785}{\textbf{Thm.}\quad$f,g$ 可测，$\lambda\in\R$，则 $f+g$，$f+\lambda$，$\lambda f$，$f\cdot g$，$f/g\ (g\ne0,\mathrm{a.e.})$ 可测.}
\handline{72}{222}{650}{\textbf{Thm.}\quad$f_k$：可测，则 $\sup_{k\in\N}f_k$，$\inf_{k\in\N}f_k$ 可测}
\handline{72}{247}{610}{进而 $\limsup_k f_k$，$\liminf_k f_k$ 可测.}
\handline{72}{273}{715}{\textbf{Pf.}\quad$\forall a\in\R,\quad\{\sup f_k\leq a\}=\bigcap_{k=1}^{\infty}\{f_k\leq a\}\in\R\quad\therefore\sup_{k\in\N}f_k$ 可测}
\handline{72}{301}{580}{$\therefore\inf_{k\in\N}f_k=-\sup_{k\in\N}(-f_k)$ 也可测}
\handline{72}{329}{700}{$\therefore\limsup_k f_k=\inf_{n\in\N}(\sup_{k\geq n}f_k),\quad\liminf_k f_k=\sup_{n\in\N}(\inf_{k\geq n}f_k)$ 也都可测.}
\handline{72}{385}{730}{\textbf{Def.}\quad$E\subseteq\R^n$，$\chi_E(x):=\begin{cases}1&x\in E\\0&x\notin E\end{cases}$，$\chi_E$ 称为 $E$ 的特征函数.}
\handline{72}{416.026}{620}{设 $E_k\subseteq\R^n$ 两两不交，$a_k\in\R$，两两不同}
\handline{72}{441}{600}{$ S(x):=\sum_{k=1}^N a_k\chi_{E_k}(x)$ 称为简单函数.}
\handline{72}{472}{790}{\textbf{Prop.}\quad$ S(x)=\sum_{k=1}^N a_k\chi_{E_k}(x)$ 可测 $\iff E_k\in\mathcal M$，$k=1,2,\ldots,N$.}
\handline[AnnotationBlue]{72}{499}{740}{\textbf{Thm.}\quad 任意 $\R^n$ 上的实值函数 $f$，可写为简单函数列 $\{f_k\}$ 的极限. \redann{（逐点）}}
\handline{72}{529}{720}{\textbf{Pf.}\quad 不妨设 $f\geq0$，否则考虑 $f=f^+-f^-$.}
\handline{72}{563}{770}{令 $ f_k(x)=\sum_{j=1}^{k\cdot2^k}\left(\frac{j-1}{2^k}\right)\chi_{\{\frac{j-1}{2^k}\leq f\leq\frac{j}{2^k}\}}(x)+k\chi_{\{f\geq k\}}(x)$，是简单函数.}
\handline{72}{614}{600}{显然，$f_k(x)\leq f_{k+1}(x)$.}
\handline{72}{654}{815}{若 $f(x)<+\infty$，则 $\exists$ 充分大的 $k_0\in\N$，s.t. $k_0>f(x)$，$\Rightarrow f(x)-f_k(x)\leq\dfrac1{2^k}\ (k\geq k_0)$.}
\handline{72}{687}{490}{$\therefore f_k(x)\nearrow f(x)$.}
\handline{72}{721}{760}{若 $f(x)=+\infty$，则 $f_k(x)=k$，$\therefore\lim_{k\to\infty}f_k(x)=\lim_{k\to\infty}k=+\infty=f(x)$}
\handline{72}{752}{660}{$\therefore f_k(x)$ 逐点收敛到 $f(x)$.}
\handline[NoteRed]{72}{788}{710}{特别地，若 $f$ 有界，$|f|\leq M$，$\exists k_0\in\N$，s.t. $k_0\geq M\geq|f|$.}
\handline[NoteRed]{72}{822}{540}{$\Rightarrow |f-f_k|\leq\dfrac1{2^k}\ (k\geq k_0)\Rightarrow f_k\rightrightarrows f$.}
\handline[NoteRed]{72}{854}{520}{即此时 $f_k$ 一致收敛到 $f$.}
\end{handpage}

% SOURCE PAGE 025
% SECTION: 4.2 | Semicontinuous Function
\begin{handpage}{25}
\handtitle{72}{130}{670}{§3.2 Semicontinuous Function}
\handline{72}{157}{650}{\textbf{Def.}\quad\circled{1} $E\subseteq\R^n$，$E\in\mathcal M$，$x_0\in E$.}
\handline{72}{184}{720}{若 $\limsup_{x\to x_0;\,x\in E}f(x)\leq f(x_0)$ 则称 $f$ 在 $x_0$ 处上半连续. \ann{（u.s.c.）}}
\handline{72}{212}{680}{$\liminf_{x\to x_0;\,x\in E}f(x)\geq f(x_0)$ 则称 $f$ 在 $x_0$ 处下半连续. \ann{（l.s.c.）}}
\handline{72}{243}{790}{\circled{2} $f$ 在 $E$ 上上半连续（下半连续）$\overset{\mathrm{def}}\iff\forall x\in E$，$f$ 在 $x$ 处上半连续（下半连续）}
\handline{72}{271}{680}{\textbf{Ex.}\quad\circled{1} $f(x_0)=+\infty$，则 $f$ 在 $x_0$ 处 u.s.c.}
\handline{72}{297}{630}{$f(x_0)=-\infty$，则 $f$ 在 $x_0$ 处 l.s.c.}
\handline{72}{332}{700}{\circled{2} 若 $f(x_0)\in\R$，则 $f$ 在 $x_0$ 处 u.s.c.}
\handline{72}{359}{735}{$\iff\forall M>f(x_0),\ \exists\delta>0,\ \mathrm{s.t.}\ \forall x\in(x_0-\delta,x_0+\delta)\cap E,\ f(x)<M$.}
\handline[NoteRed]{72}{382}{675}{$\downarrow$ 从上往下看，看不到洞，故上半连续，而从下往上看，可以看到一个洞}
\handline{72}{395}{80}{\circled{3}}
\handdiagram{686}{393}{242}{96}{
\draw[AnnotationBlue] (0,.9)--(3,.9);\fill[white] (1.2,.9)circle(.055);\draw[AnnotationBlue](1.2,.9)circle(.055);\fill[AnnotationBlue](1.2,1.45)circle(.035);\draw[AnnotationBlue,dashed](1.2,.9)--(1.2,.6);\node[AnnotationBlue,left]at(0,.9){$0$};\node[AnnotationBlue,left]at(1.03,1.45){$1$};\node[AnnotationBlue,right]at(1.2,.6){$x_0$};
\draw[AnnotationBlue](0,0)--(3,0);\fill[white](1.2,0)circle(.055);\draw[AnnotationBlue](1.2,0)circle(.055);\fill[AnnotationBlue](1.2,-.55)circle(.035);\node[AnnotationBlue,left]at(0,0){$0$};\node[AnnotationBlue,left]at(1.05,-.55){$-1$};
}
\handline[AnnotationBlue]{72}{418}{510}{$f(x_0)=1$，$f(x)=0$，$x\ne x_0$，则 $f$ 在 $x_0$ 处 u.s.c.}
\handline[AnnotationBlue]{72}{447}{510}{$f(x_0)=-1$，$f(x)=0$，$x\ne x_0$，则 $f$ 在 $x_0$ 处 l.s.c.}
\handline[NoteRed]{72}{478}{596}{$\uparrow$ 从下往上看，看不到洞，故下半连续，而从上往下看，可以看到一个洞.}
\handline{72}{517}{730}{\textbf{Thm.}\quad$E\subseteq\R^n$，$f:E\to\R\cup\{\pm\infty\}$.}
\handline{72}{547}{740}{则 $f$ 在 $E$ 上 u.s.c. \ann{（l.s.c.）}$\iff\{f<a\}$ \ann{（$\{f>a\}$）}是开集，$\forall a\in\R$.}
\handline{72}{579}{650}{\textbf{Pf.}\quad“$\Rightarrow$” 设 $p\in\{f<a\}$，则 $f(p)<a$}
\handline{72}{607}{640}{$\therefore\exists\delta>0$，s.t. $x\in(p-\delta,p+\delta)\cap E$ 时，$f(x)<a$}
\handline{72}{638}{650}{$\therefore B(p,\delta)\cap E\subseteq\{f<a\}\quad\therefore\{f<a\}$ 是开集}
\handline{72}{674}{680}{“$\Leftarrow$” $x_0\in E$，不妨设 $f(x_0)<\infty$.}
\handline{72}{706}{600}{设 $m>f(x_0)$，则 $\{f<m\}$ 是开集}
\handline{72}{739}{635}{$x_0\in\{f<m\}\quad\therefore\exists\delta>0$，s.t. $B(x_0,\delta)\cap E\subseteq\{f<m\}$}
\handline{72}{774}{620}{$\therefore\forall x\in B(x_0,\delta)\cap E,\quad f(x)<m$}
\handline{72}{805}{500}{$\therefore f$ 在 $x_0$ 处 u.s.c.}
\handline{72}{842}{820}{\textbf{Cor.}\quad$E\subseteq\R^n$，可测，若 $f$ 在 $E$ 上 u.s.c.（或 l.s.c.）则 $f$ 可测}
\handline[AnnotationBlue]{72}{873}{720}{更进一步，由于 $\forall a\in\R$，$\{f<a\}$（$\{f>a\}$）是开集}
\handline[AnnotationBlue]{72}{904}{680}{故 $f$ 不仅可测，还是 Borel 可测函数.}
\end{handpage}

% SOURCE PAGE 026
% SECTION: 4.3 | Properties of Measurable functions and Theorems of Egorov and Lusin
\begin{handpage}{26}
\handtitle{72}{126}{790}{§3.3 Properties of Measurable functions and Theorems of Egorov and Lusin}
\handline{72}{150}{740}{\textbf{Ex.}\quad\circled{1} $E=[0,1]$，$f_n(x)=n$，$\lim_{n\to\infty}f_n(x)=+\infty$}
\handline{72}{176}{685}{\circled{2} $E_n=[0,n]$，$f_n(x)=\chi_{E_n}(x)$，$\lim_{n\to\infty}f_n(x)=\chi_{[0,+\infty)}$}
\handline[AnnotationBlue]{72}{200}{680}{但 \circled{1}\circled{2} 都是逐点收敛，不是一致收敛.}
\handline{72}{227}{690}{\textbf{Thm.}\quad\ann{（Egorov's theorem）}}
\handline{72}{266}{850}{$E\subseteq\R^n$，$E\in\mathcal M$，$m(E)<\infty$，$f_n$ 是 $E$ 上可测函数，逐点收敛到一个几乎处处有限的函数 $f$.}
\handline{72}{293}{845}{则 $\forall\varepsilon>0$，$\exists E$ 的闭子集 $F$，s.t. $m(E\setminus F)<\varepsilon$，且 $f_n$ 在 $F$ 上一致收敛到 $f$.}
\handline[NoteRed]{72}{320}{745}{若 $\R^n$ 换成其他测度空间，要改成“可测子集”.}
\handline{72}{353}{856}{\textbf{Pf.}\quad(I) $\forall\varepsilon,\eta>0$，令 $E_m=\{x\in E\mid|f_k(x)-f(x)|<\varepsilon,\ k>m\}=\bigcap_{k>m}\{x\in E\mid|f_k(x)-f(x)|<\varepsilon\}\in\mathcal M$}
\handline{72}{382}{755}{$E_m\subseteq E_{m+1}$；由 $f_n\to f$ a.e. $\Rightarrow E_m\nearrow E\setminus Z$，其中 $m(Z)=0$.}
\handline{72}{415}{770}{$\therefore\lim_{m\to\infty}m(E_m)=m(E)<\infty\quad\therefore\exists m_0$，s.t. $E_{m_0}\subseteq E$ 且 $m(E\setminus E_{m_0})<\dfrac\eta2$}
\handline{72}{446}{750}{$E_{m_0}$ 可测，$\therefore\exists$ 闭子集 $F$，s.t. $m(E_{m_0}\setminus F)<\dfrac\eta2$}
\handline{72}{478}{770}{$\therefore F\subseteq E$，且 $m(E\setminus F)\leq m(E\setminus E_{m_0})+m(E_{m_0}\setminus F)<\eta$}
\handline{72}{508}{735}{由 $F\subseteq E_{m_0}\Rightarrow\forall x\in F$，$k>m_0$，$|f_k(x)-f(x)|<\varepsilon$.}
\handline{72}{541}{760}{(II) 特别地将 (I) 中 $\varepsilon$ 取为 $\dfrac1m\ (m\in\N)$，$\eta$ 取为 $\dfrac\varepsilon{2^m}$.}
\handline{72}{574}{780}{故 $\forall\varepsilon>0$，$m\in\N$，$\exists$ 闭子集 $F_m$ s.t. $m(E\setminus F_m)<\dfrac\varepsilon{2^m}$.}
\handline{72}{605}{745}{$\exists k(m,\varepsilon)\in\N$，s.t. $\forall x\in F_m$，$k>k(m,\varepsilon)$ 时，$|f_k(x)-f(x)|<\dfrac1m$}
\handline{72}{638}{790}{令 $\widetilde F=\bigcap_{m=1}^{\infty}F_m$，则 $\widetilde F$ 是 $E$ 的闭子集，且 $ m(E\setminus\widetilde F)\leq\sum_{m=1}^{\infty}m(E\setminus F_m)<\sum_{m=1}^{\infty}\frac\varepsilon{2^m}=\varepsilon$}
\handline{72}{668}{685}{由于 $\forall m\in\N$，$\widetilde F\subseteq F_m$.}
\handline{72}{701}{750}{固定 $\varepsilon=\varepsilon_0>0$，$\forall\delta>0$，$\exists n_0=n_0(\delta)\in\N$ s.t. $n>n_0$ 时 $\dfrac1n<\delta$.}
\handline{72}{736}{690}{$\exists k_0=k_0(n_0,\varepsilon_0)=k_0(n_0(\delta),\varepsilon_0)=g_0(\delta)$，s.t. $k>k_0$ 时.}
\handline{72}{769}{610}{$\forall x\in\widetilde F\subseteq F_n$，$|f_k(x)-f(x)|<\dfrac1n<\delta$.}
\handline{72}{804}{550}{$\therefore f_n\rightrightarrows f$ 在 $\widetilde F$ 上.}
\handline[NoteRed]{72}{838}{660}{$\left(\forall\delta>0,\ \exists k_0=g_0(\delta)\in\N.\ \mathrm{s.t.}\ k>k_0.\right.$}
\handline[NoteRed]{72}{868}{630}{$\left.\forall x\in\widetilde F,\quad |f_k(x)-f(x)|<\delta\right)$}
\end{handpage}

% SOURCE PAGE 027
\begin{handpage}{27}
\handline{72}{119}{550}{\textbf{Def.}\quad$E\subseteq\R^n$，$E\in\mathcal M$，$f$ 定义在 $E$ 上，若 $\forall\varepsilon>0$，$\exists$ 闭子集 $F\subseteq E$ s.t.}
\handline{72}{147}{660}{\circled{1} $m(E\setminus F)<\varepsilon$\quad\circled{2} $f|_F$ 是连续的，则称 $f$ 有性质 C.}
\handline[AnnotationBlue]{646}{112}{223}{（完全是为叙述 Lusin}
\handline[AnnotationBlue]{72}{133}{215}{的方便性而定义的）}
\handline{72}{175}{650}{\textbf{Thm.}\quad 简单函数都具有性质 C.}
\handline{72}{202}{785}{\textbf{Pf.}\quad 设 $ f(x)=\sum_{k=1}^N a_k\chi_{E_k}(x)$，$a_k\in\R$ 互不相同，$E_k\in\mathcal M\Rightarrow f$ 可测}
\handline{72}{227}{730}{$\forall\varepsilon>0$，$\exists$ 闭集 $F_k\subseteq E_k$，s.t. $m(E_k\setminus F_k)<\dfrac\varepsilon N$，}
\handline{72}{254}{760}{令 $ F=\bigcup_{k=1}^N F_k$，$ E=\bigcup_{k=1}^N E_k$，则 $F$ 是 $E$ 的闭子集，$ m(E\setminus F)\leq\sum_{k=1}^N m(E_k\setminus F_k)<N\cdot\frac\varepsilon N=\varepsilon$}
\handline{72}{282}{690}{$\forall k\in\N$，相对于 $F$ 上的子拓扑，$F_k$ 是 $F$ 上的开集}
\handline{72}{309}{670}{而 $f|_{F_k}\equiv a_k\quad\therefore f|_{F_k}$ 连续 $\quad\therefore f|_F$ 连续.}
\handline{72}{336}{640}{\textbf{Thm.}\quad\ann{（Lusin's theorem）}}
\handline{72}{364}{810}{$E\subseteq\R^n$，$E\in\mathcal M$，$f$ 定义在 $E$ 上，则 $f$ 可测且几乎处处有限 $\iff f$ 具有性质 C.}
\handline{72}{392}{740}{\textbf{Pf.}\quad“$\Rightarrow$” 存在简单函数 s.t. $S_k(x)\to f(x)$.}
\handline{72}{419}{630}{$1^\circ\quad m(E)<\infty$}
\handline{72}{446}{765}{$\forall\varepsilon>0$，$\exists$ 闭子集 $F_k\subseteq E$，s.t. $m(E\setminus F_k)<\dfrac\varepsilon{2^{k+1}}$，$S_k|_{F_k}$ 连续，}
\handline{72}{474}{785}{用 Egorov 定理，$\exists$ 闭集 $F_0\subseteq E$ s.t. $m(E\setminus F_0)<\dfrac\varepsilon2$，$S_k\rightrightarrows f$ 在 $F_0$ 上.}
\handline{72}{502}{805}{令 $ F=F_0\cap\left(\bigcap_{k=1}^{\infty}F_k\right)\subseteq E$ 是闭集，则 $ m(E\setminus F)<\sum_{k=0}^{\infty}m(E\setminus F_k)<\frac\varepsilon2+\frac\varepsilon2=\varepsilon$}
\handline{72}{530}{790}{$F\subseteq F_k\ (\forall k\in\N)\quad\therefore S_k|_F$ 连续，又 $F\subseteq F_0\quad\therefore S_k|_F\rightrightarrows f|_F$}
\handline{72}{558}{680}{由连续性定理 $\Rightarrow f|_F$ 连续.}
\handline{72}{589}{850}{$2^\circ$ 一般地，令 $E_k=E\cap R_k$，$R_k=\{x\in\R^n\mid k-1\leq|x|<k\}$，$\therefore E=\bigcup_{k=1}^{\infty}E_k$，$E_k$ 两两不交.}
\handline{72}{619}{795}{$m(E_k)<\infty$，由 $1^\circ$：$\forall k\in\N$，$\exists$ 闭子集 $F_k\subseteq E_k$，s.t. $m(E_k\setminus F_k)<\dfrac\varepsilon{2^k}$，且 $f|_{F_k}$ 连续}
\handline{72}{650}{780}{$F_k$ 有界闭 $\Rightarrow F_k$ 紧致，\quad$F_k\subseteq E_k\Rightarrow F_k$ 两两不交.\quad$\therefore d(F_k,F_j)>0\ (k\ne j)$}
\handline{72}{679}{730}{令 $ F=\bigcup_{k=1}^{\infty}F_k$，则 $F$ 是闭集}
\handline{72}{709}{760}{且 $ m(E\setminus F)\leq\sum_{k=1}^{\infty}m(E\setminus F_k)<\sum_{k=1}^{\infty}\frac\varepsilon{2^k}=\varepsilon$，$f|_F$ 连续.}
\handline{72}{745}{795}{“$\Leftarrow$” $\forall k\in\N$，$\exists$ 闭集 $F_k\subseteq E$ s.t. $m(E\setminus F_k)<\dfrac1k$ 且 $f|_{F_k}$ 连续.}
\handline{72}{776}{805}{令 $ F=\bigcup_{k=1}^{\infty}F_k$ 是 $F_\sigma$ 集，故可测，且 $m(E\setminus F)\leq m(E\setminus F_k)<\dfrac1k\quad\therefore m(E\setminus F)=0$.}
\handline{72}{808}{805}{令 $E=F\cup Z$，$Z$ 是零测集.\quad$f|_F$ 连续，$\therefore f$ 在 $F$ 上有限，$\therefore f$ 几乎处处有限}
\handline{72}{841}{790}{$\forall a\in\R,\quad\{x\in E\mid f(x)>a\}=\left(\bigcup_{k=1}^{\infty}\{x\in F_k\mid f(x)>a\}\right)\cup\{x\in Z\mid f(x)>a\}$.}
\handline{72}{878}{795}{而 $f|_F$ 连续，$\therefore\{x\in F\mid f(x)>a\}$ 是相对于 $F$ 的子拓扑的开集.}
\handline{72}{908}{775}{$\therefore\{x\in F_k\mid f(x)>a\}=F_k\cap\{x\in F\mid f(x)>a\}$ 是 Borel 集，故可测.}
\handline{72}{933}{520}{而 $\{x\in Z\mid f(x)>a\}\subseteq Z\quad\therefore$ 是零测集}
\handline{72}{950}{360}{$\therefore\{x\in E\mid f(x)>a\}$ 可测.\quad$\therefore f$ 可测.}
\end{handpage}

% SOURCE PAGE 028
% SECTION: 4.4 | Convergence in Measure
\begin{handpage}{28}
\handtitle{72}{57}{690}{§3.4 Convergence in Measure}
\handline{72}{83}{775}{\textbf{Def.}\quad$E\subseteq\R^n$，$E\in\mathcal M$，$f,f_k$ 在 $E$ 上可测}
\handline{72}{111}{745}{若 $\forall\varepsilon>0,\quad\lim_{k\to\infty}m(\{x\in E\mid|f(x)-f_k(x)|>\varepsilon\})=0$}
\handline{72}{138}{690}{则称 $f_k$ 依测度收敛到 $f$，记为 $f_k\xrightarrow{m}f$}
\handline{72}{168}{790}{\textbf{Thm.}\quad$E\subseteq\R^n$，$m(E)<+\infty$，$f,f_k$ 可测，且在 $E$ 上几乎处处有限}
\handline{72}{194}{720}{若 $f_k\to f$ a.e. 则 $f_k\xrightarrow{m}f$.}
\handline{72}{223}{805}{\textbf{Pf.}\quad 由 Egorov 定理，$\forall\varepsilon>0,\ \eta>0$，$\exists$ 闭子集 $F\subseteq E$，$m(E\setminus F)<\eta$，$m_0\in\N$}
\handline{72}{251}{695}{s.t. $k>m_0$ 时，$\{x\in E\mid|f_k(x)-f(x)|>\varepsilon\}\subseteq E\setminus F$}
\handline{72}{279}{675}{$\therefore m(\{x\in E\mid|f_k(x)-f(x)|>\varepsilon\})<m(E\setminus F)<\eta$}
\handline{72}{305}{705}{令 $\eta\to0+\Rightarrow\lim_{k\to\infty}m(\{x\in E\mid|f(x)-f_k(x)|>\varepsilon\})=0$}
\handline{72}{332}{570}{$\therefore f_k\xrightarrow{m}f$.}
\handline{72}{369}{760}{\textbf{Ex.}\quad\circled{1} 依测度收敛 \redann{$\nRightarrow$} 几乎处处收敛}
\handline{72}{400}{635}{令 $f_1=\chi_{[0,1]}$}
\handline{72}{429}{650}{$f_2=\chi_{[0,\frac12]}$，$f_3=\chi_{[\frac12,1]}$}
\handline{72}{457}{700}{$f_4=\chi_{[0,\frac14]}$，$f_5=\chi_{[\frac14,\frac12]}$，$f_6=\chi_{[\frac12,\frac34]}$，$f_7=\chi_{[\frac34,1]}$}
\handline{72}{485}{500}{$\ldots\ldots$}
\handline{72}{516}{700}{$n\in\N$，$2^k\leq n<2^{k+1}$，$ f_n=\chi_{[\frac{n-2^k}{2^k},\frac{n-2^k+1}{2^k}]}$}
\handline{72}{549}{735}{则令 $E=[0,1]$，$\forall\varepsilon>0$，$m(\{x\in E\mid|f_n(x)|>\varepsilon\})\leq\dfrac1{2^k}\to0$.}
\handline{72}{580}{590}{$\therefore f_n\xrightarrow{m}0$}
\handline{72}{612}{735}{但 $\forall x\in(0,1)$，$\#\{n\mid f_n(x)=1\}=\infty\quad\therefore f_n(x)\not\to f(x)$.}
\handline{72}{646}{740}{\circled{2} 去掉 $m(E)<\infty$：几乎处处收敛 \redann{$\nRightarrow$} 依测度收敛}
\handline{72}{680}{700}{$E=[0,+\infty)$，$f=\chi_{[0,+\infty)}$}
\handline{72}{715}{745}{令 $f_k=\chi_{[k,+\infty)}$，则 $\forall x\in(0,+\infty)$，当 $k>x$ 时：$f_k(x)=\chi_{[0,k]}(x)=1$}
\handline{72}{747}{530}{$\therefore f_k(x)\to f(x)=1$}
\handline{72}{782}{730}{但 $\forall k\in\N$，$m(\{x\in E\mid|f_k(x)-f(x)|>\dfrac12\})=+\infty$}
\handline{72}{816}{600}{$\therefore f_k\not\xrightarrow{m}f\equiv1$.}
\end{handpage}

% SOURCE PAGE 029
\begin{handpage}{29}
\handline{72}{130}{817}{\textbf{Thm.} $E\subseteq\R^n$，$E\in\mathcal M$，$f_k,f$ 在 $E$ 上可测，$f_k\xrightarrow m f$}
\handline{72}{156}{763}{则存在子列 $\{f_{N_j}\}$ s.t. $f_{N_j}\to f$ a.e.（$j\to\infty$）.}
\handline{72}{183}{820}{\textbf{Pf.} $\forall j\in\N$，$\lim_{k\to\infty}m(\{x\in E\mid|f_k(x)-f(x)|>\frac1j\})=0$.}
\handline{72}{211}{756}{$\therefore\exists N_j\in\N$，s.t. $k\ge N_j$ 时 $m(\{x\in E\mid|f_k(x)-f(x)|>\frac1j\})<\frac1{2^j}$.}
\handline{72}{238}{755}{不妨设 $\forall j\in\N$，$N_{j+1}>N_j$.}
\handline{72}{267}{758}{令 $E_j=\{x\in E\mid|f_{N_j}(x)-f(x)|>\frac1j\}$，则 $m(E_j)<\frac1{2^j}$.}
\handline{72}{297}{759}{对 $m\in\N$，令 $H_m=\bigcup_{j\ge m}E_j$，则 $m(H_m)\le\sum_{j=m}^\infty\frac1{2^j}=\frac1{2^{m-1}}\to0$（$m\to\infty$）.}
\handline{72}{326}{761}{$\therefore H_m\downarrow Z$，$m(Z)=0$，$Z=\bigcap_{m=1}^\infty H_m$.}
\handline{72}{355}{758}{$x\notin H_m\Longleftrightarrow\forall j\ge m,x\notin E_j\Longleftrightarrow\forall j\ge m,\ |f_{N_j}(x)-f(x)|\le\frac1j$，$\therefore f_{N_j}$ 在 $x$ 处收敛到 $f$.}
\handline{72}{386}{763}{而 $x\notin Z\Longleftrightarrow\exists m_0\in\N$，s.t. $x\notin H_{m_0}\Rightarrow f_{N_j}$ 在 $x$ 处收敛到 $f$.}
\handline{72}{414}{761}{由于 $m(Z)=0$，$\therefore f_{N_j}\to f$ a.e.（$j\to\infty$）.}
\handline{72}{444}{838}{\textbf{Thm.} $E\subseteq\R^n$，$E\in\mathcal M$，$f_k,f$ 在 $E$ 上可测. \textcolor{AnnotationBlue}{（Cauchy 准则）}}
\handline{72}{474}{790}{则 $f_k\xrightarrow m f\Longleftrightarrow\forall\varepsilon>0,\lim_{j,k\to\infty}m(\{x\in E\mid|f_j(x)-f_k(x)|>\varepsilon\})=0$.}
\handline{72}{501}{835}{\textbf{Pf.} ``$\Rightarrow$'' trivial. 注意到 $\{|f_j-f_k|>\varepsilon\}\subseteq\{|f_j-f|>\frac\varepsilon2\}\cup\{|f_k-f|>\frac\varepsilon2\}$.}
\handline{72}{531}{792}{``$\Leftarrow$'' $\forall j\in\N$，令 $E_j=\{x\in E\mid|f_l(x)-f_k(x)|>\frac1{2^j}\}$.}
\handline{72}{559}{728}{$\lim_{l,k\to\infty}m(E_j)=0$，$\therefore\exists N_j\in\N$，s.t. $l,k>N_j$ 时，$m(E_j)<\frac1{2^j}$.}
\handline{72}{589}{599}{不妨设 $N_{j+1}>N_j$（$\forall j\in\N$）. 从而得到一个函数子列 $\{f_{N_j}\}$.}
\handline{72}{619}{599}{$\forall x\notin H_m\Rightarrow\forall j\ge m,x\notin E_j$. $\therefore|f_{N_{j+1}}(x)-f_{N_j}(x)|\le\frac1{2^j}$.}
\handline{72}{651}{605}{$\therefore\bigl|\sum_{j=m}^\infty f_{N_{j+1}}(x)-f_{N_j}(x)\bigr|\le\sum_{j=m}^\infty\frac1{2^j}=\frac1{2^{m-1}}$.}
\handline{72}{680}{434}{$\therefore f_{N_j}$ 在 $H_m^c$ 上一致收敛.}
\handline{72}{710}{430}{设收敛到 $f$，再令 $m\to\infty$，$H_m^c\uparrow Z^c$.}
\handline{72}{740}{721}{易验证 $f_{N_j}$ 在 $Z^c$ 上逐点收敛，$m(Z)=0$.}
\handline{72}{771}{719}{$\therefore f_{N_j}\to f$ a.e. $\therefore f_{N_j}\xrightarrow m f$.}
\handline{72}{803}{733}{$\therefore\forall\varepsilon>0$，$m(\{|f_{N_j}-f|>\frac\varepsilon2\})\to0$.}
\handline{72}{833}{713}{而 $\{|f_k-f|>\varepsilon\}\subseteq\{|f_k-f_{N_j}|>\frac\varepsilon2\}\cup\{|f_{N_j}-f|>\frac\varepsilon2\}$.}
\handline{72}{868}{735}{$\therefore m(\{|f_k-f|>\varepsilon\})\le m(\{|f_k-f_{N_j}|>\frac\varepsilon2\})+m(\{|f_{N_j}-f|>\frac\varepsilon2\})$.}
\handline{72}{899}{713}{$k,j\to\infty$ 时 RHS $\to0\Rightarrow m(\{|f_k-f|>\varepsilon\})\to0$（$k\to\infty$）.}
\handline{72}{933}{705}{$\therefore f_k\xrightarrow m f$.}
\end{handpage}

% SOURCE PAGE 030
% CHAPTER: 5 | The Lebesgue Integral
% SECTION: 5.1 | Definition of the Integral of a Nonnegative Function
\begin{handpage}{30}
\handtitle{72}{119}{790}{Chapter 4 The Lebesgue Integral}
\handtitle{72}{148}{789}{\S4.1 Definition of the Integral of a Nonnegative Function}
\handline{72}{178}{790}{\textbf{Def.} $E\subseteq\R^n$，$E\in\mathcal M$，$f:E\to\R\cup\{\pm\infty\}$，$f\ge0$.}
\handline{72}{205}{724}{令 $T(f;E):=\{(x,f(x))\in\R^{n+1}\mid x\in E,\ f(x)<+\infty\}$.}
\handline{72}{233}{701}{$R(f;E):=\{(x,y)\in\R^{n+1}\mid0\le y\le f(x)\ (f(x)<+\infty),\ 0\le y<+\infty\ (f(x)=+\infty)\}$.}
\handline{72}{262}{732}{若 $R(f;E)$ 是 $\R^{n+1}$ 中的可测集，则称 $f$ 在 $E$ 上的 Lebesgue 积分存在.}
\handline{72}{292}{717}{且记 $\int_Ef(x)\,dx=m(R(f;E))$. \textcolor{NoteRed}{（取值不一定有限，可以为 $+\infty$）}}
\handline{72}{324}{807}{\textbf{Thm.} $E\subseteq\R^n$，$E\in\mathcal M$，$f\ge0$，则 $f$ 在 $E$ 上可测 $\Longleftrightarrow\int_Ef(x)\,dx$ 存在.}
\handline[NoteRed]{72}{353}{734}{只证 ``$\Rightarrow$''，``$\Leftarrow$'' 需要 Fubini.}
\handline{72}{386}{815}{\textbf{Lemma.} $E_a:=\{(x,y)\mid x\in E,\ 0\le y\le a\}$，若 $a=+\infty$，$E_a:=\{(x,y)\mid x\in E,\ 0\le y<+\infty\}$.}
\handline{72}{416}{699}{$E\in\mathcal M$，则 $E_a$ 是 $\R^{n+1}$ 中可测集，且 $m(E_a)=a\,m(E)$.}
\handline{72}{445}{810}{\textbf{Pf.} (i) $a<+\infty$}
\handline{72}{479}{740}{\circled{1} $m(E)=0$，或 $E$ 是区间（闭、开、半开半闭） \textcolor{AnnotationBlue}{trivial，可作为习题}}
\handline{72}{509}{742}{\circled{2} $E$ 是开集，令 $E=\bigcup_{k=1}^\infty I_k$，$I_k$ 两两不交，半开半闭.}
\handline{72}{538}{718}{$\therefore E_a=\bigcup_{k=1}^\infty I_{k,a}$，$\therefore m(E_a)=\sum_{k=1}^\infty m(I_{k,a})=\sum_{k=1}^\infty a\,m(I_k)=a\,m(E)$.}
\handline{72}{568}{744}{\circled{3} $E$ 是 $G_\delta$ 集，且 $m(E)<+\infty$，$E=\bigcap_{k=1}^\infty G_k$，不妨设 $G_k\downarrow$（否则令 $G'_1=G_1,G'_2=G_1\cap G_2,\ldots$）}
\handline{72}{600}{717}{不妨设 $m(G_1)<+\infty$，$m(G_{k,a})=a\,m(G_k)$，且 $G_{k,a}\downarrow E_a$.}
\handline{72}{632}{733}{$\therefore m(E_a)=\lim_{k\to\infty}m(G_{k,a})=\lim_{k\to\infty}a\,m(G_k)=a\,m(E)$.}
\handline{72}{666}{746}{\circled{4} $E$ 可测且 $m(E)<+\infty$，$\exists G_\delta$ 集 $H$ s.t. $E\subseteq H$，$m(H\setminus E)=0$，即 $m(H)=m(E)<+\infty$.}
\handline{72}{699}{749}{$\therefore m(E_a)=m(H_a\setminus(H_a\setminus E_a))=m(H_a)-m(H_a\setminus E_a)=a\,m(H)-a\,m(H\setminus E)=a\,m(E)$.}
\handline{72}{733}{749}{\circled{5} $E$ 可测且 $m(E)=+\infty$，令 $E=\bigcup_{k=1}^\infty E_k$，$E_k=E\cap R_k$，则 $E_k$ 两两不交且 $m(E_k)<+\infty$.}
\handline{72}{765}{740}{$\therefore m(E_a)=\sum_{k=1}^\infty m(E_{k,a})=\sum_{k=1}^\infty a\,m(E_k)=a\,m(E)$.}
\handline{72}{800}{771}{(ii) $a=+\infty$ 取一个正数列 $a_k$ s.t. $\lim_{k\to\infty}a_k=+\infty$ 且单调增.}
\handline{72}{834}{789}{则 $m(E_{a_k})=a_km(E)$}
\handline{72}{866}{752}{$m(E_a)=\lim_{k\to\infty}m(E_{a_k})=\lim_{k\to\infty}a_km(E)=+\infty=a\,m(E)$（$m(E)\ne0$）.}
\handline{72}{899}{768}{若 $m(E)=0$，$\forall k$，$m(E_{a_k})=a_km(E)=0\Rightarrow m(E_a)=\lim_{k\to\infty}m(E_{a_k})=0$.}
\handline{72}{933}{741}{\textcolor{NoteRed}{作约定：$+\infty\cdot0=0$}，于是 $a\,m(E)=0$，$\therefore m(E_a)=a\,m(E)$.}
\end{handpage}

% SOURCE PAGE 031
\begin{handpage}{31}
\handline{72}{133}{846}{\textbf{Lemma.} $f\ge0$，且为 $E$ 上可测函数，$0\le m(E)\le+\infty$.}
\handline{72}{162}{747}{则 $T(f;E)$ 是 $\R^{n+1}$ 中的零测集.}
\handline{72}{193}{834}{\textbf{Pf.} $\forall\varepsilon>0$，令 $E_k=\{x\in E\mid k\varepsilon\le f(x)<(k+1)\varepsilon\}$，$E_\infty=\{x\in E\mid f(x)=+\infty\}$.}
\handline{72}{222}{782}{$\therefore E=E_\infty\cup(\bigcup_{k=1}^\infty E_k)$，$\sum_{k=1}^\infty m(E_k)\le m(E)$.}
\handline{72}{251}{781}{而 $T(f;E)=\bigcup_{k=1}^\infty T(f;E_k)$，$m^*(T(f;E_k))\le\varepsilon m(E_k)$.}
\handline{72}{281}{773}{$\therefore m^*(T(f;E))\le\sum_{k=1}^\infty m^*(T(f;E_k))\le\sum_{k=1}^\infty\varepsilon m(E_k)\le\varepsilon m(E)$.}
\handline{72}{313}{782}{\circled{1} $m(E)<+\infty\Rightarrow$ 令 $\varepsilon\to0+$，$m^*(T(f;E))=0$，$\therefore T(f;E)$ 为零测集.}
\handline{72}{344}{785}{\circled{2} $m(E)=+\infty$ 令 $E_k^{(1)}=E\cap R_k$，则 $m(E_k^{(1)})<+\infty$，$E_k$ 两两不交.}
\handline{72}{375}{767}{则 $E=\bigcup_{k=1}^\infty E_k^{(1)}$，$T(f;E)=\bigcup_{k=1}^\infty T(f;E_k^{(1)})$，且由\circled{1}知 $m(T(f;E_k^{(1)}))=0$.}
\handline{72}{408}{782}{$\therefore m^*(T(f;E))\le\sum_{k=1}^\infty T(f;E_k^{(1)})=0$.}
\handline{72}{437}{783}{$\therefore T(f;E)$ 是零测集.}
\handline[NoteRed]{72}{467}{856}{最后借助这两个引理来证明定理的 ``$\Rightarrow$'' 部分.}
\handline{72}{499}{856}{\textbf{Pf.} $f(x)$ 是 $E$ 上可测函数，故可取一列简单函数 $S_k(x)\to f(x)$.}
\handline{72}{531}{792}{于是 $R(S_k;E)\cup T(f;E)\uparrow R(f;E)$.}
\handline{72}{562}{794}{$\forall k\in\N$，设 $S_k=\sum_{j=1}^Na_j\chi_{E_j}$.}
\handline{72}{597}{785}{$\therefore R(S_k;E)=\bigcup_{j=1}^NE_{j,a_j}\in\mathcal M_{n+1}$.}
\handline{72}{629}{781}{而 $T(f;E)$ 是零测集，故 $T(f;E)\in\mathcal M_{n+1}$.}
\handline{72}{665}{796}{于是 $R(f;E)=\lim_{k\to\infty}R(S_k;E)\cup T(f;E)$ 也是 $\R^{n+1}$ 中的可测集.}
\handline{72}{701}{782}{即 $\int_Ef(x)\,dx$ 存在.}
\handline{72}{738}{856}{\textbf{Cor.} $S_k(x)=\sum_{j=1}^Na_j\chi_{E_j}(x)$，则 $\int_ES_k(x)\,dx=m(R(S_k;E))=\sum_{j=1}^Na_jm(E_j)$.}
\end{handpage}

% SOURCE PAGE 032
% SECTION: 5.2 | Properties of Integral
\begin{handpage}{32}
\handtitle{72}{120}{780}{\S4.2 Properties of Integral}
\handline{72}{149}{790}{\textbf{Prop.} \circled{1} $E\subseteq\R^n$，$E\in\mathcal M$，若 $f,g$ 在 $E$ 上可测，$0\le g\le f$}
\handline{72}{179}{690}{则 $\int_Eg\le\int_Ef$，特别地 $\int_E(\inf_E f)\le\int_Ef$.}
\handline{72}{208}{722}{\circled{2} $E_1\subseteq E_2\subseteq\R^n$，$E_1,E_2\in\mathcal M$，$f$ 在 $E_2$ 上可测，$f\ge0$.}
\handline{72}{238}{690}{则 $\int_{E_1}f\le\int_{E_2}f$.}
\handline{72}{268}{730}{\circled{3} $E\subseteq\R^n$，$E\in\mathcal M$，$f$ 在 $E$ 上可测，$f\ge0$.}
\handline{72}{300}{697}{若 $\int_Ef<+\infty$，则 $f<+\infty$ a.e.}
\handline{72}{332}{802}{\textbf{Pf.} 只证\circled{3}. $\forall\alpha>0$，令 $E_\alpha=\{x\in E\mid f(x)>\alpha\}$，$E_\infty=\{x\in E\mid f(x)=+\infty\}$.}
\handline{72}{363}{662}{则 $E_\infty\subseteq E_\alpha$.}
\handline{72}{395}{653}{$\alpha\,m(E_\infty)\le\alpha\,m(E_\alpha)\le\int_{E_\alpha}f\le\int_Ef<+\infty$.}
\handline{72}{427}{668}{$\Rightarrow m(E_\infty)=0$，$\therefore f<+\infty$ a.e.}
\handline{72}{461}{806}{\textbf{Thm.} \textcolor{AnnotationBlue}{(Levi's Lemma)}}
\handline{72}{493}{774}{$f_k,f\ge0$ 在 $E$ 上可测，$f_k\uparrow f$ a.e.}
\handline{72}{526}{777}{则 $\int_Ef_k\uparrow\int_Ef$，即 $\lim_{k\to\infty}\int_Ef_k=\int_E\lim_{k\to\infty}f_k=\int_Ef$.}
\handline{72}{560}{785}{\textbf{Pf.} $R(f_k;E)\cup T(f;E)\uparrow R(f;E)$}
\handline{72}{595}{712}{$\therefore\lim_{k\to\infty}m(R(f_k;E))+m(T(f;E))=m(R(f;E))$.}
\handline{72}{630}{662}{即 $\lim_{k\to\infty}\int_Ef_k=\int_Ef$.}
\handline{72}{665}{814}{\textbf{Prop.} $f\ge0$，若 $E=\bigcup_{k=1}^\infty E_k$，$E_k$ 可测，两两不交.}
\handline{72}{699}{731}{则 $\int_Ef=\sum_{k=1}^\infty\int_{E_k}f$.}
\handline{72}{734}{805}{\textbf{Pf.} $R(f;E)=\bigcup_{k=1}^\infty R(f;E_k)$}
\handline{72}{769}{699}{$\therefore m(R(f;E))=\sum_{k=1}^\infty m(R(f;E_k))$，$\therefore\int_Ef=\sum_{k=1}^\infty\int_{E_k}f$.}
\end{handpage}

% SOURCE PAGE 033
\begin{handpage}{33}
\handline{72}{130}{856}{\textbf{Thm.} $f\ge0$，在 $E$ 上可测，则 $\int_Ef=\sup\sum_{j=1}^m(\inf_{E_j}f)\,m(E_j)$.}
\handline[NoteRed]{72}{158}{810}{$\sup$ 取遍形如 $E=\bigcup_{j=1}^mE_j$（$E_j$ 两两不交，可测）的有限分解.}
\handline{72}{187}{856}{\textbf{Pf.} 若 $E=\bigcup_{j=1}^mE_j$，$E_j$ 可测，两两不交.}
\handline{72}{217}{792}{令 $S(x)=\sum_{j=1}^m(\inf_{E_j}f)\chi_{E_j}(x)$，则 $S(x)\le f(x)$.}
\handline{72}{247}{764}{$\therefore\sum_{j=1}^m(\inf_{E_j}f)\,m(E_j)=\int_ES(x)\,dx\le\int_Ef(x)\,dx$.}
\handline{72}{280}{762}{$\Rightarrow\sup\sum_{j=1}^m(\inf_{E_j}f)\,m(E_j)\le\int_Ef$.}
\handline{72}{336}{785}{另一方面，令 $S_k(x)=\sum_{j=1}^{k\cdot2^k}(\inf_{E_j}f)\chi_{E_j}(x)+k\chi_{E_0}(x)$.}
\handline{72}{370}{783}{其中 $E_0=\{x\in E\mid f(x)\ge k\}$，$E_j:=\{x\in E\mid\frac{j-1}{2^k}\le f(x)<\frac j{2^k}\}$.}
\handline{72}{402}{786}{于是 $S_k(x)\uparrow f(x)$.}
\handline{72}{433}{790}{由 Levi 引理 $\Rightarrow\int_ES_k\uparrow\int_Ef$.}
\handline{72}{466}{792}{即 $\sum_{j=1}^{k\cdot2^k}(\inf_{E_j}f)\,m(E_j)+(\inf_{E_0}f)\,m(E_0)=\int_ES_k\uparrow\int_Ef$.}
\handline{72}{500}{774}{$\therefore\sup\sum_{j=1}^m(\inf_{E_j}f)\,m(E_j)\ge\int_Ef$.}
\handline{72}{557}{779}{$\therefore\int_Ef=\sup\sum_{j=1}^m(\inf_{E_j}f)\,m(E_j)$.}
\handline[NoteRed]{72}{589}{801}{若 $\sup$ 与 $\inf$ 位置互换则结论不成立.}
\handline{72}{625}{856}{\textbf{Cor.} $f\ge0$，在 $E$ 上可测，若 $m(E)=0$，则 $\int_Ef=0$.}
\handline{72}{657}{856}{\textbf{Pf.} 两种证法.}
\handline{72}{690}{799}{\circled{1} $\int_Ef=\sup\sum_{j=1}^m(\inf_{E_j}f)\,m(E_j)$，但 $\forall E_j\subseteq E$，$m(E_j)=0$，$\therefore\int_Ef=0$.}
\handline{72}{724}{802}{\circled{2} 设 $S_k(x)\uparrow f(x)$，$S_k(x)=\sum_{j=1}^{N_k}a_{k,j}\chi_{E_{k,j}}(x)$.}
\handline{72}{757}{769}{$\forall k\in\N$，$1\le j\le N_k$，$E_{k,j}\subseteq E\Rightarrow m(E_{k,j})=0$.}
\handline{72}{792}{772}{$\therefore\forall k\in\N$，$\int_ES_k=\sum_{j=1}^{N_k}a_{k,j}m(E_{k,j})=0$. $\therefore\int_Ef=\lim_{k\to\infty}\int_ES_k=0$.}
\handline{72}{826}{856}{\textbf{Cor.} $f,g\ge0$，在 $E$ 上可测，若 $f=g$ a.e.，则 $\int_Ef=\int_Eg$.}
\handline[AnnotationBlue]{72}{858}{823}{改变零测集上的函数值，不改变积分结果.}
\handline{72}{891}{856}{\textbf{Ex.} Dirichlet 函数，$D(x)=\begin{cases}1&x\in[0,1]\setminus\Q\\0&x\in[0,1]\cap\Q\end{cases}$.}
\handline{72}{927}{814}{$D(x)$ 不是 Riemann 可积的，但 $\int_{[0,1]}D(x)\,dx=\int_{[0,1]\setminus\Q}D(x)\,dx=1\cdot m([0,1]\setminus\Q)=1$.}
\end{handpage}

% SOURCE PAGE 034
\begin{handpage}{34}
\handline{72}{103}{813}{\textbf{Prop.} $f\ge0$，在 $E$ 上可测，则 $\int_Ef=0\Longleftrightarrow f=0$ a.e.}
\handline{72}{131}{806}{\textbf{Pf.} ``$\Leftarrow$'' trivial.}
\handline{72}{160}{739}{``$\Rightarrow$'' $\forall\alpha>0$，令 $E_\alpha=\{x\in E\mid f(x)>\alpha\}$.}
\handline{72}{190}{691}{则 $0\le\alpha\,m(E_\alpha)=\int_{E_\alpha}\alpha\le\int_{E_\alpha}f\le\int_Ef=0$.}
\handline{72}{220}{686}{$\therefore m(E_\alpha)=0$.}
\handline{72}{249}{711}{$\therefore m(\{f>0\})=m(\bigcup_{n=1}^\infty\{f>\frac1n\})\le\sum_{n=1}^\infty m(E_{1/n})=0$.}
\handline{72}{278}{717}{$\therefore f=0$ a.e.}
\handline{72}{309}{822}{\textbf{Thm.} \textcolor{AnnotationBlue}{(Tchebyshev's Inequality).}}
\handline{72}{339}{791}{$f\ge0$，在 $E$ 上可测，则 $\forall\varepsilon>0$，$m(\{x\in E\mid f(x)>\varepsilon\})\le\frac1\varepsilon\int_Ef$.}
\handline{72}{369}{812}{\textbf{Pf.} $\forall\varepsilon>0$，令 $E_\varepsilon=\{x\in E\mid f(x)>\varepsilon\}$.}
\handline{72}{402}{701}{则 $\varepsilon\,m(E_\varepsilon)=\int_{E_\varepsilon}\varepsilon\le\int_{E_\varepsilon}f\le\int_Ef$.}
\handline{72}{433}{677}{$\therefore m(E_\varepsilon)\le\frac1\varepsilon\int_Ef$.}
\handline{72}{465}{831}{\textbf{Prop.} $f,g\ge0$，在 $E$ 上可测，$c\ge0$，则 \circled{1} $\int_Ecf=c\int_Ef$ \circled{2} $\int_E(f+g)=\int_Ef+\int_Eg$.}
\handline{72}{495}{822}{\textbf{Pf.} $1^\circ$ $f,g$ 是简单函数.}
\handline{72}{526}{510}{$f(x)=\sum_{j=1}^ma_j\chi_{E_j}(x)$，$g(x)=\sum_{k=1}^nb_k\chi_{F_k}(x)$.}
\handline{526}{525}{262}{$a_j,b_k\ge0$，$E_j$ 两两不交.}
\handline{699}{554}{132}{$F_k$ 两两不交.}
\handline{72}{558}{612}{$\int_Ecf=\sum_{j=1}^mca_jm(E_j)=c\sum_{j=1}^ma_jm(E_j)=c\int_Ef$.}
\handline{72}{589}{722}{$f(x)+g(x)=\sum_{j,k}(a_j+b_k)\chi_{E_j\cap F_k}(x)$.}
\handline{72}{622}{757}{$\therefore\int_E(f+g)=\sum_{j,k}(a_j+b_k)m(E_j\cap F_k)=\sum_{j,k}a_jm(E_j\cap F_k)+\sum_{j,k}b_km(E_j\cap F_k)$}
\handline{72}{655}{626}{$=\sum_ja_jm(E_j)+\sum_kb_km(F_k)=\int_Ef+\int_Eg$.}
\handline{72}{718}{765}{$2^\circ$. 一般地，设 $S_k(x)\uparrow f(x)$，$T_k(x)\uparrow g(x)$，则 $(S_k(x)+T_k(x))\uparrow(f(x)+g(x))$.}
\handline{72}{750}{244}{$cS_k(x)\uparrow cf(x)$.}
\handline{72}{781}{740}{$\therefore c\int_Ef=c\lim_{k\to\infty}\int_ES_k=\lim_{k\to\infty}\int_EcS_k=\int_Ecf$.}
\handline{72}{817}{705}{$\int_E(f+g)=\lim_{k\to\infty}\int_E(S_k+T_k)=\lim_{k\to\infty}\int_ES_k+\lim_{k\to\infty}\int_ET_k=\int_Ef+\int_Eg$.}
\handline{72}{887}{856}{\textbf{Cor.} $\varphi,f$ 在 $E$ 上可测，$0\le f\le\varphi$，$\int_Ef<+\infty$.}
\handline{72}{925}{785}{则 $\int_E(\varphi-f)=\int_E\varphi-\int_Ef$.}
\end{handpage}

% SOURCE PAGE 035
\begin{handpage}{35}
\handline{72}{165}{832}{\textbf{Thm.} $f_k\ge0$，在 $E$ 上可测，则 $\int_E(\sum_{k=1}^\infty f_k)=\sum_{k=1}^\infty\int_Ef_k$.}
\handline{72}{195}{834}{\textbf{Pf.} 令 $S_N=\sum_{k=1}^Nf_k$，则 $S_N\uparrow\sum_{k=1}^\infty f_k$.}
\handline{72}{224}{732}{$\therefore\int_ES_N\uparrow\int_E(\sum_{k=1}^\infty f_k)$.}
\handline{72}{253}{739}{而 $\int_ES_N=\int_E\sum_{k=1}^Nf_k=\sum_{k=1}^N\int_Ef_k\uparrow\sum_{k=1}^\infty\int_Ef_k$.}
\handline{72}{282}{732}{$\therefore\int_E(\sum_{k=1}^\infty f_k)=\sum_{k=1}^\infty\int_Ef_k$.}
\handline{72}{316}{844}{\textbf{Thm.} \textcolor{AnnotationBlue}{(Fatou's Lemma)}}
\handline{72}{344}{774}{$f_k\ge0$，在 $E$ 上可测，则 $\int_E\liminf_k f_k\le\liminf_k\int_Ef_k$.}
\handline{72}{375}{844}{\textbf{Pf.} 令 $g_k=\inf_{j\ge k}f_j$，则 $g_k\uparrow\liminf_j f_j$.}
\handline{72}{405}{721}{$\therefore$ 由 Levi 引理}
\handline{72}{434}{689}{$\int_Eg_k\uparrow\int_E\liminf_j f_j$.}
\handline{72}{466}{774}{即 $\lim_{k\to\infty}\int_Eg_k=\liminf_k\int_Eg_k\le\liminf_k\int_Ef_k$.}
\handline{72}{526}{768}{$\therefore\int_E\liminf_k f_k\le\liminf_k\int_Ef_k$.}
\handline{72}{558}{856}{\textbf{Thm.} \textcolor{AnnotationBlue}{(L.D.C.T., Lebesgue Dominated Convergence Theorem for Nonnegative Function)}}
\handline{72}{590}{830}{$f_k,f\ge0$，在 $E$ 上可测，若 \circled{1} $f_k\to f$ a.e. \circled{2} $\exists\varphi$，$\int_E\varphi<+\infty$，s.t. $0\le f_k\le\varphi$ a.e.}
\handline{72}{620}{790}{则 $\int_Ef_k\to\int_Ef$.}
\handline{72}{650}{856}{\textbf{Pf.} 由于 $f_k\ge0$，$\therefore$ 由 Fatou 引理 $\int_Ef=\int_E\lim_{k\to\infty}f_k=\int_E\liminf_kf_k\le\liminf_k\int_Ef_k$.}
\handline{72}{681}{799}{由于 $\varphi-f_k\ge0$，$\therefore$ 由 Fatou 引理}
\handline{72}{715}{791}{$\int_E\varphi-\int_Ef=\int_E(\varphi-f)=\int_E\lim_k(\varphi-f_k)=\int_E\liminf_k(\varphi-f_k)$}
\handline{72}{746}{707}{$\le\liminf_k\int_E(\varphi-f_k)=\int_E\varphi-\limsup_k\int_Ef_k$.}
\handline{72}{779}{746}{$\Rightarrow\limsup_k\int_Ef_k\le\int_Ef\le\liminf_k\int_Ef_k$.}
\handline{72}{811}{725}{$\therefore\liminf_k\int_Ef_k=\limsup_k\int_Ef_k=\int_Ef$.}
\handline{72}{843}{729}{$\therefore\int_Ef=\lim_{k\to\infty}\int_Ef_k$.}
\handline{72}{874}{856}{\textbf{Cor.} $f_k\ge0$，$f_k\to f$ a.e.，若 $\forall k\in\N$，$\int_Ef_k\le M$，则 $\int_Ef\le M$.}
\handline{72}{903}{856}{\textbf{Pf.} 由 Fatou 引理.}
\handline{72}{937}{762}{$\int_Ef=\int_E\liminf_kf_k\le\liminf_k\int_Ef_k\le\liminf_k M=M$.}
\end{handpage}

% SOURCE PAGE 036
% SECTION: 5.3 | The Integral of Arbitrary Measurable Functions
\begin{handpage}{36}
\handtitle{72}{124}{780}{\S4.3 The Integral of Arbitrary Measurable Functions}
\handline{72}{151}{690}{Def.\quad$f$ 是 $E$ 上可测函数.$f^+=\max(0,f)$，$f^-\equiv\max(0,-f)$.}
\handline{72}{177}{610}{则 $f^+,f^-\ge0$.$f=f^+-f^-$，$\int_E f^+,\int_E f^-$ 一定存在}
\handline{72}{202}{620}{若 $\int_E f^+,\int_E f^-$ 至少有一个有限，则定义 $\int_E f=\int_E f^+-\int_E f^-$}
\handline{72}{231}{710}{进一步，若 $\int_E f$ 存在且有限，则称 $f$ 在 $E$ 上 Lebesgue 可积，并记 $f\in L(E)$.}
\handline[Teal]{72}{257}{490}{$L(E)=\{f\mid\int_E f\text{ 存在且有限}\}$.}
\handline[Teal]{72}{284}{615}{若 $\int_E f=\int_E f^+-\int_E f^-$ 存在且有限 $\Rightarrow\int_E f^+,\int_E f^-$ 均有限.}
\handline[Teal]{72}{313}{530}{$\therefore\int_E|f|=\int_E f^++\int_E f^-$ 存在且有限 $\Rightarrow|f|\in L(E)$.}
\handline[Teal]{72}{341}{630}{反之若 $|f|\in L(E)$ 则 $f\in L(E)$.即“可积”$\iff$“绝对可积”}
\handline{72}{373}{685}{Def.\quad$0<p<\infty$.$L^p(E):=\{f\mid(\int_E|f|^p)^{\frac1p}<\infty\}$.称为 $L^p$-空间.}
\handline{72}{402}{610}{Prop.\quad 若 $f\in L^1(E)$，则 $f$ 在 $E$ 上几乎处处有限.}
\handline{72}{430}{715}{Pf.\quad$f\in L^1(E)\iff|f|\in L^1(E)\iff\int_E|f|<\infty\Rightarrow|f|$ 在 $E$ 上几乎处处有限}
\handline{72}{458}{400}{$\therefore f$ 在 $E$ 上几乎处处有限.}
\handline{72}{496}{560}{Prop.\quad 若 $\int_E f,\int_E g$ 存在，$f\le g$ a.e.则 $\int_E f\le\int_E g$.}
\handline{72}{528}{525}{Pf.\quad$f\le g$ a.e.$\Rightarrow f^+\le g^+$，$g^-\le f^-$ a.e.}
\handline{72}{558}{425}{$\therefore\int_E f^+\le\int_E g^+$，$\int_E g^-\le\int_E f^-$.}
\handline{72}{588}{260}{$\therefore\int_E f\le\int_E g$.}
\handline[NoteRed]{602}{517}{280}{由于此处未证明对一般函数}
\handline[NoteRed]{514}{547}{333}{的积分线性性，故不可令 $\phi=g-f$.}
\handline[NoteRed]{501}{586}{285}{$\int_E\phi\ge0$ 来导出 $\int_E g\ge\int_E f$.}
\handline{72}{623}{585}{Cor.\quad 若 $\int_E f,\int_E g$ 存在.$f=g$ a.e.则 $\int_E f=\int_E g$.}
\handline{72}{655}{805}{Pf.\quad$f=g$ a.e.$\Rightarrow f\le g,g\le f$ a.e.$\Rightarrow\int_E f\le\int_E g$，$\int_E g\le\int_E f\Rightarrow\int_E f=\int_E g$}
\handline{72}{689}{575}{Prop.\quad$E_1\subseteq E_2$.若 $\int_{E_2}f$ 存在则 $\int_{E_1}f$ 存在.}
\handline{72}{720}{585}{Pf.\quad$\int_{E_2}f$ 存在 $\Rightarrow\int_{E_2}f^+$ 或 $\int_{E_2}f^-$ 至少有一个有限}
\handline{72}{753}{600}{$\Rightarrow\int_{E_1}f^+$ 或 $\int_{E_1}f^-$ 至少有一个有限 $\Rightarrow\int_{E_1}f$ 存在.}
\handline{72}{791}{750}{Prop.\quad 若 $\int_E f$ 存在.$ E=\bigcup_{k=1}^\infty E_k$，$E_k$ 两两不交且可测，则 $\int_E f=\sum_{k=1}^\infty\int_{E_k}f$.}
\handline{72}{823}{585}{Pf.\quad$\int_E f$ 存在 $\Rightarrow\int_E f^+,\int_E f^-$ 至少有一个有限.}
\handline{72}{858}{680}{而 $\int_E f^+=\sum_{k=1}^\infty\int_{E_k}f^+$，$\int_E f^-=\sum_{k=1}^\infty\int_{E_k}f^-\Rightarrow$ 至少有一个级数收敛}
\handline{72}{895}{620}{$\Rightarrow\sum_{k=1}^\infty\int_{E_k}f^+-\sum_{k=1}^\infty\int_{E_k}f^-=\sum_{k=1}^\infty\int_{E_k}f^+-f^-=\sum_{k=1}^\infty\int_{E_k}f$}
\handline{72}{932}{350}{$\therefore\int_E f=\sum_{k=1}^\infty\int_{E_k}f$.}
\end{handpage}

% SOURCE PAGE 037
\begin{handpage}{37}
\handline{72}{114}{490}{Prop.\quad 若 $m(E)=0$，或 $f=0$ a.e.则 $\int_E f=0$.}
\handline{72}{145}{655}{Pf.\quad 若 $m(E)=0$ 或 $f=0$ a.e.$\Rightarrow\int_E f^+=\int_E f^-=0\Rightarrow\int_E f=0$}
\handline{72}{177}{600}{Prop.\quad\circled{1} 若 $\int_E f$ 存在，则 $\int_E(-f)$ 存在且 $\int_E(-f)=-\int_E f$}
\handline{72}{202}{575}{\circled{2} 若 $\int_E f$ 存在，$c\in\R$，则 $\int_E cf$ 存在，且 $\int_E cf=c\int_E f$.}
\handline{72}{232}{610}{\circled{3} 若 $f,g\in L^1(E)$，则 $f+g\in L^1(E)$.且 $\int_E f+\int_E g=\int_E f+g$.}
\handline{72}{261}{510}{Pf.\quad\circled{1} $\int_E f$ 存在 $\Rightarrow\int_E f^+$ 或 $\int_E f^-$ 有限.}
\handline{72}{291}{715}{而 $(-f)^+=f^-$，$(-f)^-=f^+$.$\therefore\int_E(-f)^+$ 或 $\int_E(-f)^-$ 有限.$\therefore\int_E(-f)$ 存在.}
\handline{72}{322}{620}{且 $\int_E(-f)=\int_E(-f)^+-\int_E(-f)^-=\int_E f^--\int_E f^+=-\int_E f$.}
\handline{72}{352}{420}{\circled{2}\quad$c=0$ trivial.$c=-1$ 已证.}
\handline{72}{383}{695}{$c>0$：$(cf)^+=cf^+$.$(cf)^-=cf^-\Rightarrow\int_E(cf)^+$ 或 $\int_E(cf)^-$ 有限}
\handline{72}{415}{600}{$\therefore\int_E(cf)$ 存在.且 $\int_E(cf)=\int_E(cf)^+-\int_E(cf)^-$}
\handline{72}{444}{365}{$=c\int_E f^+-c\int_E f^-=c\int_E f$}
\handline{72}{478}{805}{$c<0$：将 $c$ 写作 $-|c|$.再利用 $c=-1$ 和 $c>0$ 的结论 $\Rightarrow\int_E cf=c\int_E f$.}
\handline{72}{513}{520}{\circled{3}\quad$f,g\in L^1(E)\Rightarrow\int_E|f|,\int_E|g|<+\infty$}
\handline{72}{544}{730}{$|f+g|\le|f|+|g|\Rightarrow\int_E|f+g|\le\int_E|f|+|g|=\int_E|f|+\int_E|g|<+\infty$.}
\handline{72}{580}{280}{$\therefore f+g\in L^1(E)$}
\handline{72}{613}{335}{$\int_E f+g=\int_E f+\int_E g$}
\handline{72}{646}{420}{需分六种情况讨论，此处略}
\handline[NoteRed]{549}{576}{305}{这里用的是对正可测函数}
\handline[NoteRed]{544}{605}{340}{积分的可加性，在 \S4.2 中已证.}
\handarrow[NoteRed]{516.0}{551}{520.0}{568}
\handarrow[NoteRed]{520.0}{568}{544.0}{568}
\handline{72}{680}{605}{Cor.\quad$f,\phi$ 在 $E$ 上可测.$f\ge\phi$ a.e.$\phi\in L^1(E)$}
\handline{72}{711}{385}{则 $\int_E(f-\phi)=\int_E f-\int_E\phi$.}
\handline{72}{744}{830}{Pf.\quad$f\ge\phi\Rightarrow f^-\le\phi^-$，$\phi\in L^1(E)\Rightarrow\int_E\phi^-<+\infty\quad\therefore\int_E f^-<+\infty$.}
\handline{72}{776}{450}{(i) $f\in L^1(E)$.则由线性性显然.}
\handline{72}{810}{585}{(ii) $f\notin L^1(E)$ 由 $\int_E f^-<+\infty\Rightarrow\int_E f^+=+\infty$.}
\handline{72}{844}{575}{$\therefore\int_E f-\int_E\phi=\int_E f^+-\int_E f^--\int_E\phi=+\infty$}
\handline{72}{877}{610}{而由于 $f-\phi\notin L^1(E)$\quad\ann{（否则 $f=f-\phi+\phi\in L^1(E)$）}}
\handline{72}{911}{500}{且 $f-\phi\ge0$ a.e.$\therefore\int_E f-\phi=+\infty$}
\handline{72}{943}{380}{$\therefore\int_E f-\phi=\int_E f-\int_E\phi$}
\end{handpage}

% SOURCE PAGE 038
\begin{handpage}{38}
\handline{72}{123}{800}{Thm.\quad$f\in L^1(E)$，$f\ge0$ a.e.$g$ 在 $E$ 上可测 $\exists M>0$，$\text{s.t.}\ |g|\le M$ a.e.则 $fg\in L^1(E)$}
\handline{72}{152}{750}{Pf.\quad$|fg|\le M|f|$.$\therefore\int_E|fg|\le\int_E M|f|=M\int_E|f|<+\infty\quad\therefore fg\in L^1(E)$.}
\handline{72}{181}{500}{Cor.\quad$f\in L^1(E)$，$f\ge0$ a.e.$g$ 在 $E$ 上可测}
\handline{72}{208}{615}{$\exists\alpha,\beta\in\R$.$\text{s.t.}\ \alpha\le g\le\beta$ a.e.则 $\alpha\int_E f\le\int_E fg\le\beta\int_E f$.}
\handline{72}{238}{735}{Pf.\quad$f\ge0$ a.e.$\Rightarrow\alpha f\le fg\le\beta f$ a.e.$\Rightarrow\alpha\int_E f\le\int_E fg\le\beta\int_E f$.}
\handline[AnnotationBlue]{72}{265}{675}{下面将 Levi，Fatou，Lebesgue 的几个收敛定理推广到一般可测函数}
\handline{72}{296}{335}{Thm.\quad\ann{（Levi's Lemma）}}
\handline{72}{326}{310}{$f_k$ 在 $E$ 上可测}
\handline{72}{355}{765}{(i) 若 $f_k\uparrow f$ a.e.，$\exists\phi\in L^1(E)$，$\text{s.t.}\ f_k\ge\phi$ a.e.$(\forall k\in\N)$，则 $\int_E f_k\uparrow\int_E f$}
\handline{72}{384}{775}{(ii) 若 $f_k\downarrow f$ a.e.，$\exists\phi\in L^1(E)$，$\text{s.t.}\ f_k\le\phi$ a.e.$(\forall k\in\N)$，则 $\int_E f_k\downarrow\int_E f$.}
\handline{72}{420}{540}{Pf.\quad(i) $f_k-\phi\ge0$ a.e.且 $f_k-\phi\uparrow f-\phi$}
\handline{72}{450}{460}{$\therefore\int_E f_k-\phi\uparrow\int_E f-\phi$，由 $\int_E\phi$ 有限}
\handline{72}{480}{310}{$\Rightarrow\int_E f_k\uparrow\int_E f$.}
\handline{72}{509}{315}{(ii) 同理可证.}
\handline[NoteRed]{72}{544}{480}{若没有 $\phi$ 的控制，则有反例：}
\handline[Teal]{72}{576}{745}{$E=[0,+\infty)$.$f_k=\chi_{[k,+\infty)}$.$f_k\downarrow f\equiv0$.但 $\int_E f_k=+\infty$，$\int_E f=0$.}
\handline{72}{612}{370}{Thm.\quad\ann{（一致收敛定理）}}
\handline{72}{644}{720}{$f_k\in L^1(E)$，$f_k\rightrightarrows f$，$m(E)<\infty$，则 $f\in L^1(E)$.且 $\int_E f_k\to\int_E f$.}
\handline{72}{678}{680}{Pf.\quad$f_k\rightrightarrows f$.$\therefore\exists N\in\N$，$\text{s.t.}\ k>N$ 时.$|f-f_k|\le1\quad\forall x\in E$.}
\handline{72}{706}{475}{$\therefore|f|\le|f-f_k|+|f_k|\le1+|f_k|$.}
\handline{72}{739}{650}{$\therefore\int_E|f|\le\int_E1+|f_k|=m(E)+\int_E|f_k|<+\infty\quad\therefore f\in L^1(E)$.}
\handline{72}{774}{650}{$\forall\varepsilon>0$.$\exists N_\varepsilon\in\N$.$\text{s.t.}\ k>N_\varepsilon$ 时.$|f-f_k|\le\varepsilon\quad\forall x\in E$.}
\handline{72}{809}{670}{$\therefore\left|\int_E f_k-\int_E f\right|\le\int_E|f-f_k|\le\varepsilon m(E)\to0\ (\varepsilon\to0+)$.}
\handline{72}{841}{350}{$\therefore\int_E f_k\to\int_E f$.}
\handline[NoteRed]{72}{874}{485}{若去掉 $m(E)<+\infty$ 则有反例：}
\handline[Teal]{72}{909}{775}{$E=[0,+\infty)$\quad$f_n(x)=\dfrac1n\chi_{[0,n]}$.$f\equiv0$.$\|f_n-f\|_\infty=\dfrac1n\to0\quad\therefore f_n\rightrightarrows f$.}
\handline{72}{947}{505}{但 $\int_E f_n=\dfrac1n\int_0^n dx=1$，$\int_E f=0$.}
\end{handpage}

% SOURCE PAGE 039
\begin{handpage}{39}
\handline{72}{126}{365}{Thm.\quad\ann{（Fatou's Lemma）}}
\handline{72}{153}{856}{$f_k$ 在 $E$ 上可测.若 $\exists\phi\in L^1(E)$.$\text{s.t.}\ f_k\ge\phi$ a.e.$(\forall k\in\N)$ 则 $\int_E\liminf_k f_k\le\liminf_k\int_E f_k$.}
\handline{72}{182}{350}{Pf.\quad$f_k-\phi\ge0$ a.e.}
\handline{72}{211}{530}{$\therefore\int_E\liminf_k(f_k-\phi)\le\liminf_k\int_E(f_k-\phi)$.}
\handline{72}{244}{500}{$\overset{\prime\prime}{\int_E\liminf_k f_k-\int_E\phi}\qquad\qquad\overset{\prime\prime}{\liminf_k\int_E f_k-\int_E\phi}$}
\handline{72}{273}{430}{$\therefore\int_E\liminf_k f_k\le\liminf_k\int_E f_k$.}
\handline{72}{309}{856}{Cor.\quad$f_k$ 在 $E$ 上可测.若 $\exists\phi\in L^1(E)\ \text{s.t.}\ f_k\le\phi$ a.e.则 $\int_E\limsup_k f_k\ge\limsup_k\int_E f_k$.}
\handline{72}{346}{280}{Thm.\quad\ann{（L.D.C.T）}}
\handline{72}{379}{800}{$f_k$ 在 $E$ 上可测.$f_k\to f$ a.e.若 $\exists\phi\in L^1(E)$，$\text{s.t.}\ \phi\ge0$，$|f_k|\le\phi$ a.e.}
\handline{72}{410}{475}{则 $\lim_{k\to\infty}\int_E f_k=\int_E\lim_{k\to\infty}f_k=\int_E f$.}
\handline{72}{444}{570}{Pf.\quad$|f_k|\le\phi$ a.e.$\Rightarrow0\le f_k+\phi\le2\phi$ a.e.}
\handline{72}{477}{565}{由正可测版本的 LDCT $\Rightarrow\int_E f_k+\phi\to\int_E f+\phi$}
\handline{72}{509}{435}{又 $\phi\in L^1(E)\quad\therefore\int_E f_k\to\int_E f$.}
\handline[NoteRed]{72}{544}{630}{若没有 $\phi$ 的控制.则会有与 Levi 引理一样的反例.}
\handline{72}{582}{340}{Cor.\quad\ann{（有界收敛定理）}}
\handline{72}{616}{730}{$f_k$ 在 $E$ 上可测.$f_k\to f$ a.e.$m(E)<\infty$.\quad\redann{$\Rightarrow M\in L^1(E)$.}}
\handline{72}{649}{650}{若 $\exists M>0$，$\text{s.t.}\ |f_k|<M$，$\forall k\in\N$.则 $\int_E f_k\to\int_E f$}
\handline{72}{688}{740}{Pf.\quad 令 $\phi\equiv M$ 即可.$\int_E\phi=Mm(E)<+\infty\quad\therefore\phi\in L^1(E)$.}
\handline[NoteRed]{72}{723}{555}{L.D.C.T 中若失去 $\phi$ 的控制，还有反例：}
\handline[Teal]{72}{757}{650}{$E=[0,+\infty)$.$f_n(x)=n\chi_{(0,\frac1n]}(x)$.$f(x)\equiv0$.}
\handline{72}{797}{690}{则 $f_n\to f$ a.e.但 $\int_E f_n=\int_0^{\frac1n}n\,dx=1$.$\int_E f=0$.}
\end{handpage}

% SOURCE PAGE 040
% SECTION: 5.4 | Relation between Riemann-Stieltjes and Lebesgue Integrals and the L^p Spaces, 0<p<infinity
\begin{handpage}{40}
\handtitle{72}{98}{825}{\S4.4 Relation between Riemann--Stieltjes and Lebesgue Integrals}
\handtitle{72}{132}{570}{and the $L^p$ Spaces，$0<p<\infty$.}
\handline{72}{163}{830}{Def.\quad$f$ 在 $E$ 上可测，$\alpha\in\R$，令 $w(\alpha)=w_f(\alpha)=w_{f,E}(\alpha):=m(\{x\in E\mid f(x)>\alpha\})$.}
\handline{72}{190}{445}{称为 $f$ 在 $E$ 上的分布函数.}
\handline[Teal]{72}{220}{795}{注：\circled{1} $\forall\alpha>\beta$.$\{x\in E\mid f(x)>\alpha\}\subseteq\{x\in E\mid f(x)>\beta\}$.$\therefore w(\alpha)\le w(\beta)$}
\handline[Teal]{72}{247}{310}{$\therefore w(\alpha)$ 单调下降.}
\handline[Teal]{72}{276}{480}{\circled{2} 若 $m(E)<+\infty$，$f$ 在 $E$ 上几乎处处有限}
\handline[Teal]{72}{308}{600}{随着 $\alpha\uparrow+\infty$，$\{x\in E\mid f(x)>\alpha\}\downarrow\{x\in E\mid f(x)=+\infty\}$.}
\handline[Teal]{72}{339}{340}{于是 $\lim_{\alpha\to+\infty}w(\alpha)=0$.}
\handline[Teal]{72}{369}{685}{类似可得 $\lim_{\alpha\to-\infty}w(\alpha)=m(E)$.于是 $w(\alpha)$ 在 $\R$ 上从 $m(E)$ 降到 $0$.}
\handline[Teal]{72}{401}{770}{$\therefore w(\alpha)$ 是 $(-\infty,+\infty)$ 上的有界变差函数.且 $V[w;-\infty,+\infty]=m(E)$}
\handline[NoteRed]{72}{434}{625}{本节都假设 $m(E)<+\infty$，$f$ 在 $E$ 上几乎处处有限.}
\handline{72}{472}{790}{Lemma.\quad$\alpha,\beta\in\R$，$\alpha<\beta$，则 $m(\{x\in E\mid\alpha<f(x)\le\beta\})=w(\alpha)-w(\beta)$.}
\handline{72}{506}{635}{Def.\quad$w(\alpha+):=\lim_{\varepsilon\to0+}w(\alpha+\varepsilon)$.$w(\alpha-):=\lim_{\varepsilon\to0+}w(\alpha-\varepsilon)$}
\handline{72}{540}{595}{Lemma.\quad\circled{1} $w(\alpha)=w(\alpha+)$，$\forall\alpha\in\R$\quad\ann{即 $w$ 右连续}}
\handline{72}{574}{650}{\circled{2} $w(\alpha-)=m(\{x\in E\mid f(x)\ge\alpha\})=m(\{x\in E\mid f(x)=\alpha\})+w(\alpha)$.}
\handline{72}{608}{605}{Pf.\quad\circled{1} 设 $\varepsilon_k\downarrow0+$ 则 $\{f>\alpha+\varepsilon_k\}\uparrow\{f>\alpha\}$.}
\handline{72}{645}{525}{$\therefore\lim_{k\to\infty}m(\{f>\alpha+\varepsilon_k\})=m(\{f>\alpha\})=w(\alpha)$}
\handline{72}{681}{300}{$\therefore w(\alpha+)=w(\alpha)$.}
\handline{72}{715}{550}{\circled{2} 同样设 $\varepsilon_k\downarrow0$.则 $\{f>\alpha-\varepsilon_k\}\downarrow\{f\ge\alpha\}$}
\handline{72}{754}{595}{由于 $m(E)<+\infty\quad\therefore\lim_{k\to\infty}m(\{f>\alpha-\varepsilon_k\})=m(\{f\ge\alpha\})$}
\handline{72}{787}{355}{$\therefore w(\alpha-)=m(\{f\ge\alpha\})$.}
\handline[AnnotationBlue]{72}{824}{550}{$w$ 的所有间断点都是跳跃间断点.}
\handline{72}{863}{856}{Cor.\quad\circled{1} $w(\alpha-)-w(\alpha)=m(\{f=\alpha\})$\quad\ann{$w$ 在 $\alpha$ 连续 $\iff m(\{f=\alpha\})=0$.}}
\handline{72}{904}{805}{\circled{2} $\alpha<\beta\in\R$.$w$ 在 $(\alpha,\beta)$ 上为常数 $\iff m(\{\alpha<f<\beta\})=0$.}
\end{handpage}

% SOURCE PAGE 041
\begin{handpage}{41}
\handline{72}{132}{665}{Thm.\quad$a<b\in\R$.$x\in E$，$a<f(x)\le b$.则 $\int_E f=-\int_a^b\alpha\,dw(\alpha)$}
\handline[AnnotationBlue]{451}{151}{150}{（Lebesgue）}
\handline[AnnotationBlue]{611}{142}{105}{（R-S）}
\handline{72}{196}{735}{Pf.\quad$m(E)<+\infty$.$f$ 可测 $\Rightarrow\int_E f^+,\int_E f^-$ 有限 $\Rightarrow\int_E f$ 存在.}
\handline{72}{221}{660}{$\alpha$ 是 $[a,b]$ 上连续函数.$w(\alpha)$ 是 B.V.的 $\Rightarrow\int_a^b\alpha\,dw(\alpha)$ 存在.}
\handline{72}{253}{570}{考虑分划.$\{a=\alpha_0<\alpha_1<\cdots<\alpha_n=b\}$.}
\handline{72}{285}{580}{令 $E_j=\{x\in E\mid\alpha_j<f(x)\le\alpha_{j+1}\}$.$0\le j\le n-1$.}
\handline{72}{316}{405}{则 $E_j$ 两两不交.且 $ E=\bigcup_{j=0}^{n-1}E_j$.}
\handline{72}{351}{560}{考虑 $ S_\Gamma=\sum_{j=0}^{n-1}\alpha_jm(E_j)$，$ S_\Gamma'=\sum_{j=0}^{n-1}\alpha_{j+1}m(E_j)$.}
\handline{72}{382}{350}{则 $S_\Gamma\le\int_E f\le S_\Gamma'$.}
\handline{72}{417}{715}{而 $ S_\Gamma=\sum_{j=0}^{n-1}\alpha_j(w(\alpha_j)-w(\alpha_{j+1}))=-\sum_{j=0}^{n-1}\alpha_j(w(\alpha_{j+1})-w(\alpha_j))$.}
\handline{72}{451}{565}{$ S_\Gamma'=-\sum_{j=0}^{n-1}\alpha_{j+1}(w(\alpha_{j+1})-w(\alpha_j))$.}
\handline{72}{485}{790}{由 $-\int_a^b\alpha\,dw(\alpha)$ 存在 故 令 $\|\Gamma\|\to0\Rightarrow S_\Gamma,S_\Gamma'\to-\int_a^b\alpha\,dw(\alpha)$.}
\handline{72}{521}{465}{$\therefore\int_E f=-\int_a^b\alpha\,dw(\alpha)$.}
\handline{72}{561}{710}{Thm.\quad$f$ 在 $E$ 上可测，$a<b\in\R$.令 $E_{ab}=\{x\in E\mid a<f(x)\le b\}$.}
\handline{72}{594}{425}{则 $\int_{E_{ab}}f=-\int_a^b\alpha\,dw(\alpha)$.}
\handline{72}{637}{630}{Pf.\quad 考虑 $f\mathbf1_{E_{ab}}$.则 $\int_{E_{ab}}f=-\int_a^b\alpha\,dw_{ab}(\alpha)$.}
\handline[NoteRed]{72}{674}{490}{$w_{ab}$ 是 $f$ 在 $E_{ab}$ 上的分布函数.}
\handline{72}{713}{690}{但由于 $\forall a\le\alpha<\beta\le b$.$w_{ab}(\alpha)-w_{ab}(\beta)=w(\alpha)-w(\beta)$}
\handline{72}{753}{695}{故 结合 R-S 积分定义 $\Rightarrow\int_a^b\alpha\,dw(\alpha)=\int_a^b\alpha\,dw_{ab}(\alpha)$}
\handline{72}{790}{430}{$\therefore\int_{E_{ab}}f=-\int_a^b\alpha\,dw(\alpha)$.}
\end{handpage}

% SOURCE PAGE 042
\begin{handpage}{42}
\handline{72}{120}{780}{Def.\quad 若 $\lim_{\substack{a\to-\infty\\b\to+\infty}}\int_a^b\alpha\,dw(\alpha)$ 存在，则定义 $\int_{-\infty}^{+\infty}\alpha\,dw(\alpha):=\lim_{\substack{a\to-\infty\\b\to+\infty}}\int_a^b\alpha\,dw(\alpha)$.}
\handline{72}{151}{575}{Thm.\quad$\int_E f$，$\int_{-\infty}^{+\infty}\alpha\,dw(\alpha)$ 若其中一个存在且有限}
\handline{72}{180}{610}{则另一个也存在且有限.且 $\int_E f=-\int_{-\infty}^{+\infty}\alpha\,dw(\alpha)$.}
\handline{72}{210}{510}{Pf.\quad\circled{1} 若 $\int_E f$ 存在且有限 即 $f\in L^1(E)$.}
\handline{72}{240}{460}{$\forall a<b\in\R$.$\int_{E_{ab}}f=-\int_a^b\alpha\,dw(\alpha)$.}
\handline{72}{270}{515}{$a\to-\infty,b\to+\infty\quad E_{ab}\uparrow(E\setminus Z)$\quad 其中 $m(Z)=0$.}
\handline{72}{301}{430}{$\therefore\int_{E_{ab}}f\to\int_{E\setminus Z}f=\int_E f$}
\handline{72}{337}{620}{$\therefore\lim_{\substack{a\to-\infty\\b\to+\infty}}\int_a^b\alpha\,dw(\alpha)$ 存在.且 $\int_E f=-\int_{-\infty}^{+\infty}\alpha\,dw(\alpha)$}
\handline{72}{389}{420}{\circled{2} 若 $\int_{-\infty}^{+\infty}\alpha\,dw(\alpha)$ 存在且有限，}
\handline{72}{422}{400}{则 $\int_0^{+\infty}\alpha\,dw(\alpha)$ 存在且有限}
\handline{72}{458}{665}{$b>0$.$\int_{E_{0b}}f=-\int_0^b\alpha\,dw(\alpha)\to-\int_0^{+\infty}\alpha\,dw(\alpha)\quad(b\to+\infty)$.}
\handline{72}{490}{480}{而 $\int_{E_{0b}}f=\int_{E_{0b}}f^+\to\int_E f^+\ (b\to+\infty)$}
\handline{72}{523}{395}{$\therefore\int_E f^+=-\int_0^{+\infty}\alpha\,dw(\alpha)$.}
\handline{72}{557}{345}{考虑 $-a>0$，$a<0$.}
\handline{72}{590}{515}{$\int_{E_{a0}}f=-\int_a^0\alpha\,dw(\alpha)\to-\int_{-\infty}^0\alpha\,dw(\alpha)$}
\handline{72}{624}{630}{$\int_{E_{a0}}f=\int_{E_{a0}}-f^-=\int_E-f^-\quad\therefore\int_E-f^-=-\int_{-\infty}^0\alpha\,dw(\alpha)$}
\handline{72}{657}{700}{$\therefore\int_E f=\int_E f^+-f^-=-\int_0^{+\infty}\alpha\,dw(\alpha)-\int_{-\infty}^0\alpha\,dw(\alpha)=-\int_{-\infty}^{+\infty}\alpha\,dw(\alpha)$.}
\handline{72}{694}{700}{Def.\quad$f,g$ 在 $E$ 上可测.若 $\forall\alpha\in\R$.$w_{f,E}(\alpha)=w_{g,E}(\alpha)$}
\handline{72}{728}{480}{则称 $f,g$ 是等可测或等分布的.}
\handline{72}{769}{745}{Cor.\quad$f,g$ 在 $E$ 上等可测，且 $f\in L(E)$，则 $g\in L(E)$ 且 $\int_E f=\int_E g$.}
\handline{72}{804}{620}{Pf.\quad$\int_E f$ 存在且有限 $\Rightarrow-\int_{-\infty}^{+\infty}\alpha\,dw_{f,E}(\alpha)=\int_E f$.}
\handline{72}{839}{610}{$\therefore-\int_{-\infty}^{+\infty}\alpha\,dw_{g,E}(\alpha)=-\int_{-\infty}^{+\infty}\alpha\,dw_{f,E}(\alpha)=\int_E f$}
\handline{72}{878}{520}{$\therefore\int_E g=-\int_{-\infty}^{+\infty}\alpha\,dw_{g,E}(\alpha)=\int_E f$.}
\handline{72}{910}{100}{Ex.}
\handdiagram{713}{874}{215}{119}{
\draw[->] (0,0)--(5,0);\draw[->] (0,0)--(0,4);\draw[Teal] (0,3)--(3,0);\draw[AnnotationBlue] (0,0)--(3,3);\draw[Teal] (0,2.2)--(3,2.2);\node[left,Teal] at (0,2.2) {$\alpha$};\draw[dashed,Teal] (.8,0)--(.8,2.2);\draw[dashed,Teal] (2.2,0)--(2.2,2.2);\draw[dashed,AnnotationBlue] (3,0)--(3,3);\foreach \x in {.05,.15,.25,.35,.45,.55,.65,.75,2.25,2.35,2.45,2.55,2.65,2.75,2.85,2.95}{\draw[Teal] (\x,0)--(\x+.05,.2);}
}
\handline{72}{928}{470}{$f(x)=1-x$，$g(x)=x$.$E=[0,1]$.}
\handline{72}{968}{385}{$f,g$ 在 $E$ 上等分布.}
\end{handpage}

% SOURCE PAGE 043
\begin{handpage}{43}
\handline{72}{147}{790}{\underline{Def.} 设 $T=\{0=\alpha_0<\alpha_1<\cdots<\alpha_k\cdots\}$, $\lim_{n\to\infty}a_n=\infty$, $m(E)<\infty$, 且 $0\le f<+\infty$ a.e.}
\handline{72}{175}{760}{$\|T\|:=\sup_k(\alpha_{k+1}-\alpha_k)$ 令 $E_k=\{x\in E\mid\alpha_k\le f(x)<\alpha_{k+1}\}$}
\handline{72}{200}{755}{于是 $E=\left(\bigcup_{k=0}^{\infty}E_k\right)\cup Z$, $Z=\{x\in E\mid f(x)=+\infty\}$, $m(Z)=0$}
\handline{72}{227}{745}{定义 $\phi_T:=\sum_{k=0}^{\infty}\alpha_k\chi_{E_k}$, $\psi_T:=\sum_{k=0}^{\infty}\alpha_{k+1}\chi_{E_k}$}
\handline{72}{252}{650}{则 $\phi_T\le f\le\psi_T$}
\handline{72}{278}{680}{$\Rightarrow\int_E\phi_T\le\int_E f\le\int_E\psi_T$}
\handline{72}{305}{770}{定义 $s_T:=\sum_{k=0}^{\infty}\alpha_km(E_k)$. $S_T:=\sum_{k=0}^{\infty}\alpha_{k+1}m(E_k)$ 则 $s_T\le\int_Ef\le S_T$.}
\handline{72}{340}{790}{\underline{Thm.} $E\in\mathcal M$. $m(E)<\infty$, $0\le f<+\infty$ a.e. 则}
\handline{72}{368}{725}{$\int_Ef=\lim_{\|T\|\to0}s_T=\lim_{\|T\|\to0}S_T$.}
\handline{72}{396}{810}{Pf. $0\le S_T-s_T=\sum_{k=0}^{\infty}(\alpha_{k+1}-\alpha_k)m(E_k)\le\|T\|\sum_{k=0}^{\infty}m(E_k)=\|T\|\cdot m(E)\to0\ (\|T\|\to0)$.}
\handline{72}{424}{650}{\circled{1} $\int_Ef<+\infty$}
\handline{72}{453}{735}{于是 $\overline{\lim}_{\|T\|\to0}s_T\le\int_Ef$, $\underline{\lim}_{\|T\|\to0}S_T\ge\int_Ef$.}
\handline{72}{481}{760}{另一方面 $\overline{\lim}_{\|T\|\to0}S_T\le\lim_{\|T\|\to0}(S_T-s_T)+\overline{\lim}_{\|T\|\to0}s_T\le0+\int_Ef=\int_Ef$}
\handline{72}{511}{720}{$\therefore\lim_{\|T\|\to0}S_T=\int_Ef$, 同理 $\lim_{\|T\|\to0}s_T=\int_Ef$.}
\handline{72}{540}{705}{\circled{2} $\int_Ef=+\infty$. 则 $S_T=+\infty$.}
\handline{72}{569}{760}{又 $\lim_{\|T\|\to0}(S_T-s_T)=0\quad\therefore S_T-s_T$ 有界 $\therefore s_T=+\infty\quad\therefore s_T=S_T=\int_Ef=+\infty$.}
\handline{72}{603}{790}{\underline{Thm.} $E\in\mathcal M$, $m(E)<+\infty$, $a<f(x)\le b$ a.e.}
\handline{72}{633}{760}{令 $\phi:[a,b]\to\R$ 连续, 则 $\int_E\phi\circ f=-\int_a^b\phi(\alpha)\,dw(\alpha)$.}
\handline{72}{665}{810}{Pf. 作分划 $T_k=\{a=\alpha_0^{(k)}<\alpha_1^{(k)}<\cdots<\alpha_n^{(k)}=b\}$, 且令 $\|T_k\|\to0\ (k\to\infty)$.}
\handline{72}{697}{765}{$E_j=\{x\in E\mid\alpha_{j-1}^{(k)}<f(x)\le\alpha_j^{(k)}\}$. 定义 $S_k=\sum_{j=1}^n\alpha_j^{(k)}\chi_{E_j}$ 则 $S_k\to f$}
\handline{72}{726}{730}{由于 $\phi$ 连续 $\therefore\phi\circ S_k\to\phi\circ f$.}
\handline{72}{755}{730}{$\phi$ 有界. $\therefore|\phi\circ S_k|\le M$. 由于 $m(E)<\infty$}
\handline{72}{786}{750}{故由有界收敛定理 $\Rightarrow\int_E\phi\circ S_k\to\int_E\phi\circ f.\quad(k\to\infty)$.}
\handline{72}{817}{780}{而 $\int_E\phi\circ S_k=\sum_{j=1}^n\phi(\alpha_j^{(k)})m(E_j)=\sum_{j=1}^n\phi(\alpha_j^{(k)})(w(\alpha_{j-1}^{(k)})-w(\alpha_j^{(k)}))$.}
\handline{72}{852}{590}{$=-\sum_{j=1}^n\phi(\alpha_j^{(k)})(w(\alpha_j^{(k)})-w(\alpha_{j-1}^{(k)}))$.}
\handline{72}{887}{760}{$\therefore$ 当 $\|T_k\|\to0$ 时. $\int_E\phi\circ S_k\to-\int_a^b\phi(\alpha)\,dw(\alpha)$}
\handline{72}{918}{740}{$\therefore\int_E\phi\circ f=-\int_a^b\phi(\alpha)\,dw(\alpha)$.}
\end{handpage}

% SOURCE PAGE 044
\begin{handpage}{44}
\handline{72}{110}{660}{\underline{Def.} $\phi:\R\to\R$}
\handline{72}{140}{755}{若 $\lim_{\substack{a\to-\infty\\b\to+\infty}}\int_a^b\phi(\alpha)\,dw(\alpha)$ 存在则定义 $\int_{-\infty}^{+\infty}\phi(\alpha)\,dw(\alpha):=\lim_{\substack{a\to-\infty\\b\to+\infty}}\int_a^b\phi(\alpha)\,dw(\alpha)$.}
\handline{72}{166}{790}{\underline{Thm.} $E\in\mathcal M$, $m(E)<+\infty$ $f$ 在 $E$ 上几乎处处有限}
\handline{72}{194}{760}{$\phi:\R\to\R$ 连续. 若 $\phi\circ f\in L^1(E)$, 则 $\int_{-\infty}^{+\infty}\phi(\alpha)\,dw(\alpha)$ 存在}
\handline{72}{222}{720}{且 $\int_E\phi\circ f=-\int_{-\infty}^{+\infty}\phi(\alpha)\,dw(\alpha)$.}
\handline{72}{252}{800}{Pf. $a<b\in\R$. $E_{ab}=\{x\in E\mid a<f(x)\le b\}$.}
\handline{72}{284}{750}{$\therefore\int_{E_{ab}}\phi\circ f=-\int_a^b\phi(\alpha)\,dw_{ab}(\alpha)=-\int_a^b\phi(\alpha)\,dw(\alpha)$}
\handline{72}{308}{730}{由 $\phi\circ f\in L^1(E)$}
\handline{72}{339}{790}{$\Rightarrow\int_E\phi\circ f=\lim_{\substack{a\to-\infty\\b\to+\infty}}\int_{E_{ab}}\phi\circ f=-\lim_{\substack{a\to-\infty\\b\to+\infty}}\int_a^b\phi(\alpha)\,dw(\alpha)=-\int_{-\infty}^{+\infty}\phi(\alpha)\,dw(\alpha)$.}
\handline[Teal]{72}{369}{810}{特别地, 取 $\phi(\alpha)=|\alpha|^p\quad0<p<\infty$. 若 $f\in L^p(E)$, $|f|^p\in L^1(E)$}
\handline[Teal]{72}{400}{755}{则 $\int_E|f|^p=-\int_{-\infty}^{+\infty}|\alpha|^p\,dw_f(\alpha)=-\int_0^{+\infty}\alpha^p\,dw_{|f|}(\alpha)$.}
\handline[NoteRed]{72}{429}{487}{$w_f(\alpha)$ 是 $f$ 的分布函数. $w_{|f|}(\alpha)$ 是 $|f|$ 的分布函数.}
\handline{72}{461}{790}{\underline{Thm.} \ann{（$L^p$ 版本的 Tchebyshev's Inequality）}}
\handline{72}{494}{800}{$0<p<\infty$, $\forall\alpha>0$, $m(\{x\in E\mid f(x)>\alpha\})\le\dfrac1{\alpha^p}\int_{\{f>\alpha\}}f^p$}
\handline{72}{526}{790}{\underline{Thm.} $0<p<\infty$, $f\in L^p(E)$. 则 $\lim_{\alpha\to+\infty}\alpha^pw(\alpha)=0$.}
\handline{72}{557}{800}{Pf. $\forall k\in\N$. 令 $f_k=\begin{cases}f&f>k.\\0&\text{其他}\end{cases}$ 由于 $f\in L^p(E)\Rightarrow|f|<\infty$ a.e.}
\handline{72}{585}{385}{$\therefore f_k\to0$ a.e.}
\handline{72}{620}{745}{$|f_k|^p\le|f|^p.\quad |f|^p\in L(E)\quad\therefore$ 由 L.D.C.T. $\Rightarrow\int_Ef_k^p\to\int_E0=0$}
\handline{72}{654}{740}{即 $\int_{\{f>k\}}f_k^p\to0$. 而 $k^pw(k)\le\int_{\{f>k\}}f_k^p$ \ann{（Tchebyshev）}}
\handline{72}{685}{750}{$\therefore\lim_{k\to+\infty}k^pw(k)=0\quad\therefore\lim_{\alpha\to+\infty}\alpha^pw(\alpha)=0$.}
\handline{72}{718}{835}{\underline{Thm.} $0<p<\infty$, $f\ge0$, $f\in L^p(E)$. 则 $\int_Ef^p=-\int_0^{\infty}\alpha^p\,dw(\alpha)=p\int_0^{\infty}\alpha^{p-1}w(\alpha)\,d\alpha$.}
\handline{72}{751}{810}{Pf. $\forall a<b\in\R$. $\int_a^b\alpha^p\,dw(\alpha)=\left.\alpha^pw(\alpha)\right|_a^b-\int_a^bpw(\alpha)\alpha^{p-1}\,d\alpha$.}
\handline{72}{784}{775}{$b\to+\infty.\quad b^pw(b)=0.\quad a\to0+,\quad a^pw(a)\le a^pm(E)\to0$.}
\handline{72}{821}{765}{$\therefore\int_{-\infty}^{+\infty}\alpha^p\,dw(\alpha)=-\int_{-\infty}^{+\infty}p\alpha^{p-1}w(\alpha)\,d\alpha$.}
\handline{72}{858}{740}{$\therefore\int_Ef^p=-\int_0^{\infty}\alpha^p\,dw(\alpha)=p\int_{-\infty}^{+\infty}\alpha^{p-1}w(\alpha)\,d\alpha$.}
\end{handpage}

% SOURCE PAGE 045
% SECTION: 5.5 | Riemann and Lebesgue Integrals
\begin{handpage}{45}
\handtitle{72}{117}{800}{$\S4.5\quad$ Riemann and Lebesgue Integrals}
\handline{72}{146}{800}{\underline{Thm.} $f:[a,b]\to\R$ 且 $f$ 有界. 设 $f$ 在 $[a,b]$ 上 Riemann 可积.}
\handline{72}{172}{765}{则 $f$ 在 $[a,b]$ 上 Lebesgue 可积且 $\int_{[a,b]}f=(R)\int_a^bf(x)\,dx$}
\handline{72}{202}{840}{Pf. 考虑分划 $T_k=\{a=\alpha_0^{(k)}<\cdots<\alpha_{n_k}^{(k)}=b\}$. 不妨设 $T_k\subseteq T_{k+1}$, 且 $\|T_k\|\to0\ (k\to\infty)$.}
\handline{72}{234}{780}{设 $l_k=\sum_{j=1}^{n_k}\left(\inf_{[\alpha_{j-1}^{(k)},\alpha_j^{(k)}]}f\right)\chi_{[\alpha_{j-1}^{(k)},\alpha_j^{(k)}]}$, $u_k=\sum_{j=1}^{n_k}\left(\sup_{[\alpha_{j-1}^{(k)},\alpha_j^{(k)}]}f\right)\chi_{[\alpha_{j-1}^{(k)},\alpha_j^{(k)}]}$.}
\handline{72}{284}{785}{则 $l_k$, $u_k$ 都是简单函数, 故均可测. 且 $l_k\uparrow l$, $u_k\downarrow u$.}
\handline{72}{319}{720}{于是 $l\le f\le u$.}
\handline{72}{344}{775}{对于 $l_k\uparrow l$. $l_k\le l\le f\le u\le u_k$. 而 $u_k\in L^1([a,b])$.}
\handline{72}{377}{805}{$\therefore$ 由 L.D.C.T. $\int_{[a,b]}l=\lim_{k\to\infty}\int_{[a,b]}l_k=\lim_{k\to\infty}\sum_{j=1}^{n_k}\left(\inf_{[\alpha_{j-1}^{(k)},\alpha_j^{(k)}]}f\right)(\alpha_j^{(k)}-\alpha_{j-1}^{(k)})$}
\handline{494}{409.037}{360}{$=(R)\int_a^bf(x)\,dx$.}
\handline{72}{411}{400}{同理 $\int_{[a,b]}u=(R)\int_a^bf(x)\,dx$.}
\handline{72}{443}{745}{$\int_{[a,b]}(u-l)=0$ 且 $u-l\ge0\quad\Rightarrow u=l$ a.e.}
\handline{72}{472}{740}{又 $l\le f\le u\quad\therefore f=u=l$ a.e. $\therefore f$ 可测.}
\handline{72}{505}{755}{由于 $f$ 有界, $m([a,b])<+\infty\quad\therefore\int_{[a,b]}f$ 存在.}
\handline{72}{536}{750}{$\therefore\int_{[a,b]}f=\int_{[a,b]}u=(R)\int_a^bf(x)\,dx$.}
\handline[NoteRed]{72}{572}{800}{一个细节, 为什么存在 $l,u$, s.t. $l_k\uparrow l$, $u_k\downarrow u$. 且 $l,u$ 可测.}
\handline[NoteRed]{72}{602}{815}{可以对每个点, $x\in[a,b]$ 考虑. 固定 $x$. $\{l_k(x)\}_{k=1}^{\infty}$ 单增有上界, $\{u_k(x)\}_{k=1}^{\infty}$ 单减有下界.}
\handline[NoteRed]{72}{632}{785}{从而 $l_k(x)\to l(x)$, $u_k(x)\to u(x)$, 从而定义好 $l,u$.}
\handline[NoteRed]{72}{666}{790}{同时, 由于可测函数的上下极限仍可测, 故 $l,u$ 可测.}
\handline{72}{706}{856}{\underline{Thm.} $f\ge0$, $E=[a,b]$. 若 \circled{1} $f\in R[a+\varepsilon,b]$. $\forall\varepsilon\in(0,b-a)$ \circled{2} $\lim_{\varepsilon\to0+}(R)\int_{a+\varepsilon}^bf(x)\,dx=I$ 存在}
\handline{72}{737}{770}{则 $f$ 在 $[a,b]$ 上可测, 且 $\int_a^bf=I$.}
\handline{72}{770}{856}{Pf. $f\in R[a+\dfrac1n,b]$. $\Rightarrow f$ 在 $[a+\dfrac1n,b]$ 上可测. 则 $\int_{[a+\frac1n,b]}f=(R)\int_{a+\frac1n}^bf(x)\,dx$.}
\handline{72}{803}{805}{令 $f_n=f\cdot\chi_{[a+\frac1n,b]}$. 则 $f_n\uparrow f$. \redann{（$f_n$ 可测, 故极限 $f$ 也可测）}}
\handline{72}{838}{820}{且 $\lim_{n\to\infty}\int_{[a,b]}f_n=\lim_{n\to\infty}\int_{[a+\frac1n,b]}f=\lim_{n\to\infty}(R)\int_{a+\frac1n}^bf(x)\,dx=I$}
\handline{72}{875}{780}{由 Levi 引理 $\Rightarrow\int_{[a,b]}f=\lim_{n\to\infty}\int_{[a,b]}f_n=I$.}
\end{handpage}

% SOURCE PAGE 046
\begin{handpage}{46}
\handline{72}{98}{750}{\underline{Thm.} \ann{(Lebesgue)}}
\handline{72}{124}{755}{$f:[a,b]\to\R$ 且有界}
\handline{72}{151}{785}{则 $f\in R[a,b]\Longleftrightarrow f$ 在 $[a,b]$ 上几乎处处连续.}
\handline{72}{181}{830}{Pf. “$\Rightarrow$” $f\in R[a,b]$. 与之前一样定义 $T_k$, $l_k$, $u_k$, $l,u$. $\|T_k\|\to0$}
\handline{72}{208}{770}{则 $l_k\uparrow l$, $u_k\downarrow u$. 且 $l=f=u$ a.e.}
\handline{72}{237}{780}{令 $Z_1=\{x\in[a,b]\mid u(x)\ne l(x)\}$, $Z_2=\bigcup_{k=1}^{\infty}T_k$. $Z=Z_1\cup Z_2$}
\handline{72}{267}{735}{则 $m(Z)=0$. 下证 $\forall x\notin Z$. $f$ 在 $x$ 处连续, 即 $\omega_f(x)=0$}
\handline{72}{295}{770}{若 $f$ 在 $x$ 处不连续. 则 $\exists\varepsilon_0$. s.t. $\forall\delta>0$. $\sup\{f(s)-f(t)\mid s,t\in(x-\delta,x+\delta)\}\ge\varepsilon_0$.}
\handline{72}{324}{765}{由于 $\|T_k\|\to0$. $\therefore\forall\delta>0$. $\exists k_0$ s.t. $k>k_0$ 时. $\|T_k\|<\delta$}
\handline{72}{355}{780}{由于 $x\notin Z$, 故 $x$ 不是分划点. $\therefore\forall k>k_0$. $\exists j(k)$ s.t. $x\in I_{j(k)}^{(k)}=[x_{j(k)}^{(k)},x_{j(k)+1}^{(k)}]$}
\handline{72}{385}{745}{由于 $x_{j(k)+1}^{(k)}-x_{j(k)}^{(k)}\le\|T_k\|<\delta\quad\therefore I_{j(k)}^{(k)}\subseteq(x-\delta,x+\delta)$.}
\handline{72}{417}{750}{$\therefore u_k(x)-l_k(x)=\operatorname{osc}(f,I_{j(k)}^{(k)})\ge\omega_f(x)\ge\varepsilon_0$.}
\handline{72}{447}{770}{令 $k\to\infty$. $u(x)-l(x)\ge\varepsilon_0\quad\therefore x\in Z_1$, 与 $x\notin Z$ 矛盾.}
\handline{72}{478}{790}{“$\Leftarrow$” 若 $f$ 在 $[a,b]$ 上几乎处处连续.}
\handline{72}{511}{720}{由于 $f$ 有界 $\therefore f\in L[a,b]$.}
\handline{72}{541}{745}{考虑分划 $T_k'$, $\|T_k'\|\to0$. 同样定义 $l_k',u_k'$. 则 $l_k'\le f\le u_k'$.}
\handline{72}{572}{760}{令 $Z_1=\{x\mid f\text{ 在 }x\text{ 处不连续}\}$, $Z_2=\bigcup_{k=1}^{\infty}T_k'$. $Z=Z_1\cup Z_2$, 则 $m(Z)=0$.}
\handline{72}{605}{775}{若 $x\notin Z$ 则 $f$ 在 $x$ 处连续. $\forall\varepsilon>0$, $\exists\delta>0$, $x_1,x_2\in(x-\delta,x+\delta)$, $|f(x_1)-f(x_2)|<\varepsilon$}
\handline{72}{640}{780}{$\|T_k'\|\to0$ 且 $x\notin Z_2\Rightarrow\exists I_k'{}^{(j)}\subseteq(x-\delta,x+\delta)$. 从而 $0\le u_k'(x)-l_k'(x)\le\varepsilon$.}
\handline{72}{675}{760}{而 $l_k'(x)\le f\le u_k'(x)\quad\therefore k\to\infty$. $u_k'(x)\to f(x)$, $l_k'(x)\to f(x)$.}
\handline{72}{709}{765}{由于 $f$ 有界 $\Rightarrow u_k',l_k'$ 一致有界. 而 $m([a,b])=b-a<\infty$}
\handline{72}{745}{780}{故由有界收敛定理 $\Rightarrow\int_{[a,b]}u_k'=\int_{[a,b]\setminus Z}u_k'\to\int_{[a,b]\setminus Z}f=\int_{[a,b]}f$.}
\handline{72}{777}{720}{同理 $\int_{[a,b]}l_k'\to\int_{[a,b]}f$.}
\handline{72}{815}{780}{而 $\int_{[a,b]}u_k'$ 为 Riemann 上和 $U_k$, $\int_{[a,b]}l_k'$ 为 Riemann 下和 $L_k$}
\handline{72}{852}{715}{$\therefore U_k,L_k\to\int_{[a,b]}f$.}
\handline{72}{888}{775}{从而 $f\in R[a,b]$. 且 $(R)\int_a^bf(x)\,dx=\int_{[a,b]}f$.}
\end{handpage}

% SOURCE PAGE 047
% CHAPTER: 6 | Repeated Integration
% SECTION: 6.1 | Fubini's Theorem
\begin{handpage}{47}
\handtitle{72}{140}{780}{Chapter 5 Repeated Integration}
\handtitle{72}{167}{740}{$\S5.1\quad$ Fubini's Theorem}
\handline{72}{191}{750}{\underline{Thm.} \ann{(Fubini)}}
\handline{72}{214}{795}{$x\in\R^n$, $x=(x_1,x_2,\cdots,x_n)$, $y\in\R^m$, $y=(y_1,\cdots,y_m)$. $(x,y)\in\R^{n+m}$}
\handline{72}{240}{790}{$I_1=\{x\in\R^n\mid a_j\le x_j\le b_j,\ 1\le j\le n\}$, $I_2=\{y\in\R^m\mid c_k\le y_k\le d_k,\ 1\le k\le m\}$.}
\handline{72}{265}{790}{设 $I=I_1\times I_2\subseteq\R^{m+n}$, $f\in L^1(I)$ 记为 $L^1(dxdy)$.}
\handline{72}{292}{800}{则 \circled{1} 对 a.e. $x\in I_1$. $f(x,y)$ 在 $I_2$ 上可测且 $\in L^1(I_2)=L^1(dy)$}
\handline{72}{319}{760}{\circled{2} $\int_c^df(x,y)\,dy$ 在 $I_1$ 上可测且 $\in L^1(I_1)=L^1(dx)$}
\handline{72}{348}{780}{\circled{3} $\iint_If(x,y)\,dxdy=\int_{I_1}\left(\int_{I_2}f(x,y)\,dy\right)\,dx$}
\handline[Teal]{72}{373}{790}{这个叙述是对称的, 互换 $x,y$, 结论也成立.}
\handline[Teal]{72}{399}{800}{在 $I$ 之外, 补充定义 $f(x,y)=0$, 故不妨设 $I_1=\R^n$, $I_2=\R^m$, $I=\R^{m+n}$}
\handline[Teal]{72}{426}{790}{先从简单情形出发, 最后回到一般情形.}
\handline[Teal]{72}{453}{780}{为了叙述方便, 先作如下定义}
\handline{72}{482}{820}{\underline{Def.} 若 $f\in L(dxdy)$ 具有 Fubini 定理中的性质.}
\handline{72}{509}{740}{则称 $f$ 具有性质 $\mathcal F$.}
\handline{72}{540}{830}{\underline{Lemma.}\circled{1} 任何有限个具有性质 $\mathcal F$ 的函数的线性组合仍具有性质 $\mathcal F$.}
\handline{72}{568}{825}{\underline{Lemma.}\circled{2} 设 $f_k\in L(dxdy)$ 且有性质 $\mathcal F$.}
\handline{72}{596}{740}{若 $f_k\uparrow f$ 或 $f_k\downarrow f$ 且 $f\in L(dxdy)$, 则 $f$ 也具有性质 $\mathcal F$.}
\handline{72}{628}{856}{Pf. Step1: $\forall k\in\N$. 设 $Z_k$ 表示所有不满足性质 $\mathcal F$ 中 \circled{1} 的 $x$ 构成的集合 则 $m(Z_k)=0$.}
\handline{72}{656}{775}{令 $Z=\bigcup_{k=1}^{\infty}Z_k$ 则 $m(Z)=0$. $x\notin Z\Rightarrow x\in Z_k\ (\forall k\in\N)$, $\therefore f_k(x,y)\in$ 在 $I_2$ 可测且 $\in L(dy)$.}
% 原稿 x不属于Z 后写 x属于Z_k，照录.
\handline{72}{684}{203}{由 $f_k(x,y)\uparrow f(x,y)$}
\handline{516}{690}{315}{$\Rightarrow f(x,y)$ 在 $I_2$ 可测 a.e. $x\in I_1$.}
\handline{72}{718}{650}{Step2: $f_k\uparrow f$. $f_k\ge f_1\in L(dxdy)$, $f\in L(dxdy)$}
\handline{72}{752}{750}{$\therefore$ 由 Levi 引理 $\int_{\R^{m+n}}f_k(x,y)\,dxdy\uparrow\int_{\R^{m+n}}f(x,y)\,dxdy<\infty$.}
\handline{72}{781}{770}{由于 $f_k$ 满足性质 $\mathcal F$. 故由 \circled{3} 知. 设 $h_k(x)=\int_{\R^m}f_k(x,y)\,dy$}
\handline{72}{815}{856}{$\int_{\R^n}h_k(x)\,dx=\int_{\R^n}\left(\int_{\R^m}f_k(x,y)\,dy\right)\,dx=\int_{\R^{m+n}}f_k(x,y)\,dxdy\uparrow\int_{\R^{m+n}}f(x,y)\,dxdy<\infty$.}
\handline{72}{845}{845}{Step3: a.e. $x\in\R^n$. $f_k(x,y)\in L(dy)$ 且 $f_k(x,y)\uparrow f(x,y)$. $f_k(x,y)\ge f_1(x,y)\in L(dy)$}
\handline{72}{878}{600}{故由 Levi 引理 $\int_{\R^m}f_k(x,y)\,dy\uparrow\int_{\R^m}f(x,y)\,dy$.}
\handline{72}{903}{423}{设 $h(x)=\int_{\R^m}f(x,y)\,dy$ 则}
\handline{72}{924}{446}{$h_k(x)\uparrow h(x)$ a.e. $\Rightarrow h(x)$ 在 $\R^n$ 上可测.}
\end{handpage}

% SOURCE PAGE 048
\begin{handpage}{48}
\handline{72}{106}{790}{Step4. 由性质 $\mathcal F$ \circled{2} $\Rightarrow h_k(x)=\int_{\R^m}f_k(x,y)\,dy\in L(dx)$.}
\handline{72}{133}{685}{$h_k(x)\uparrow h(x)$ a.e. 且 $h_k(x)\ge h_1(x)\in L(dx)$.}
\handline{72}{161}{760}{$\therefore$ 由 Levi 引理 $\int_{\R^n}h_k(x)\,dx\uparrow\int_{\R^n}h(x)\,dx$.}
\handline{72}{189}{740}{而由 Step2 可知 $\int_{\R^n}h_k(x)\,dx\uparrow\int_{\R^{m+n}}f(x,y)\,dxdy<\infty$}
\handline{72}{214}{735}{于是 $\int_{\R^n}h(x)\,dx=\int_{\R^{m+n}}f(x,y)\,dxdy<\infty$}
\handline{72}{239}{350}{$\therefore h(x)=\int_{\R^m}f(x,y)\,dy\in L(dx)$}
\handline{72}{265}{462}{且 $\int_{\R^{m+n}}f(x,y)\,dxdy=\int_{\R^n}\left(\int_{\R^m}f(x,y)\,dy\right)\,dx$.}
\handline{72}{293}{800}{Step5. 由于 $h(x)\in L(dx)\Rightarrow h(x)$ 几乎处处有限.}
\handline{72}{319}{650}{即 $\int_{\R^m}f(x,y)\,dy$ 几乎处处有限}
\handline{72}{345}{485}{$\therefore$ 对于所有 $x\in\R^n$, $f(x,y)$ 作为 $y$ 的函数 $\in L(dy)$}
\handline{72}{401}{820}{\underline{Lemma.}\circled{3} $E$ 是 $G_\delta$ 集. $E=\bigcap_{k=1}^{\infty}G_k$, $G_k$ 是开集}
\handline{72}{428}{710}{若 $m(G_1)<\infty$. 则 $\chi_E$ 具有性质 $\mathcal F$.}
\handline{72}{457}{830}{Pf. Case I: $E=J_1\times J_2$. $J_1,J_2$ 分别是 $\R^n$ 和 $\R^m$ 中的有限开区间}
\handline{72}{485}{740}{\circled{1}: $\forall x\in\R^n$ $\chi_E(x,y)$ 作为 $\R^m$ 上的函数, 显然可测.}
\handline{72}{516}{670}{且 $\int_{\R^m}\chi_E(x,y)\,dy=m(J_2)$ 或 $0<\infty\quad\therefore\chi_E(x,y)\in L(dy)$.}
\handline{72}{546}{735}{\circled{2} $\forall x\in\R^n$. $x\in J_1\quad\int_{\R^m}\chi_E(x,y)\,dy=m(J_2)$.}
\handline{72}{576}{590}{$x\notin J_1\quad\int_{\R^m}\chi_E(x,y)\,dy=0$}
\handline{72}{606}{710}{均为常值函数. 故 $\int_{\R^m}\chi_E(x,y)\,dy$ 是 $\R^n$ 上的可测函数.}
\handline{72}{638}{700}{由于 $\int_{\R^n}\left(\int_{\R^m}\chi_E(x,y)\,dy\right)\,dx=\int_{J_1}m(J_2)\,dx=m(J_1)m(J_2)<\infty$}
\handline{72}{669}{690}{$\therefore\int_{\R^m}\chi_E(x,y)\,dy\in L(dx)$.}
\handline{72}{698}{760}{\circled{3}. 由乘积测度 $\Rightarrow m(E)=m(J_1)m(J_2)$}
\handline{72}{730}{690}{$\therefore\int_{\R^{m+n}}\chi_E(x,y)\,dxdy=\int_{\R^n}\left(\int_{\R^m}\chi_E(x,y)\,dy\right)\,dx$.}
\handline{72}{765}{720}{故 $\chi_E$ 具有性质 $\mathcal F$.}
\handline{72}{798}{770}{\underline{Case II:} $E\subseteq\partial(J_1\times J_2)\Rightarrow m(E)=0\Rightarrow\int_{\R^{m+n}}\chi_E(x,y)\,dxdy=0$.}
\handline{72}{830}{745}{$\forall x\in\operatorname{int}(J_1)$. $\{y\mid(x,y)\in E\}\subseteq\partial(J_2)$. 故为 $\R^m$ 中零测集.}
\handline{72}{864}{740}{$\therefore h(x)=\int_{\R^m}\chi_E(x,y)\,dy=0$.}
\handline{72}{894}{760}{而 $m(\operatorname{int}(J_1))=m(J_1)$. $\therefore$ 对于 a.e. $x\in\R^n$. $h(x)=0$.}
\handline{72}{924}{720}{$\therefore\int_{\R^n}h(x)\,dx=0=\int_{\R^{m+n}}\chi_E(x,y)\,dxdy.\quad\therefore\chi_E$ 具有性质 $\mathcal F$.}
\end{handpage}

% SOURCE PAGE 049
\begin{handpage}{49}
\handline{72}{133}{760}{Case III: $E=(J_1\times J_2)\cup F,\quad F\subseteq\partial(J_1\times J_2)$.}
\handline{72}{161}{690}{则 $\chi_E=\chi_{J_1\times J_2}+\chi_F\Rightarrow\chi_E$ 具有性质 $\mathcal F$.}
\handline{72}{189}{750}{Case IV: $E\subseteq\R^{n+m}$ 为开集，且 $m(E)<\infty$.}
\handline{72}{215}{690}{于是 $\chi_E\in L(dxdy)$.}
\handline{72}{241}{700}{设 $E=\bigcup_{k=1}^{\infty}I_k$，$I_k$ 半开半闭且两两不交.}
\handline{72}{268}{720}{再令 $E_m=\bigcup_{k=1}^m I_k$，则 $E_m\uparrow E$，$\therefore\chi_{E_m}\uparrow\chi_E$.}
\handline{72}{295}{710}{由于 $E_m$ 具有 Case III 的形式，$\therefore\chi_{E_m}$ 具有性质 $\mathcal F$.}
\handline{72}{321}{720}{由引理 \circled{3} $\Rightarrow\chi_E$ 具有性质 $\mathcal F$.}
\handline{72}{349}{780}{Case V: $E$ 是一 $G_\delta$ 集，$E=\bigcap_{k=1}^{\infty}G_k$ 且 $m(G_1)<\infty$.}
\handline{72}{377}{720}{令 $\widehat G_m=\bigcap_{j=1}^mG_k$，则 $\widehat G_m\downarrow E$，$\therefore\chi_{\widehat G_m}\downarrow\chi_E$.}
\handline{72}{405}{720}{由 Case IV，$m(\widehat G_m)\le m(G_1)<\infty$.}
\handline{72}{434}{720}{$\therefore\chi_{\widehat G_m}$ 具有性质 $\mathcal F$. 而 $\chi_E\le\chi_{G_1}$，$\therefore\chi_E\in L(dxdy)$.}
\handline{72}{459}{730}{由引理 \circled{3} $\Rightarrow\chi_E$ 具有性质 $\mathcal F$.}
\handline{72}{487}{850}{Lemma \circled{4}. $E\subseteq\R^{n+m}$，$m(E)=0$，则 a.e. $x\in\R^n$，$E_x=\{y\in\R^m\mid(x,y)\in E\}$}
\handline{72}{515}{760}{$E_x$ 在 $\R^m$ 中零测且 $\chi_E$ 具有性质 $\mathcal F$.}
\handline{72}{543}{835}{Pf. $m(E)=0$，$\therefore\exists\,G_\delta$ 集 $H$，s.t. $E\subseteq H$，$m(E)=m(H)=0$.}
\handline{72}{572}{770}{设 $H=\bigcap_{k=1}^{\infty}G_k$，$G_k$ 开集，不妨设 $m(G_1)<\infty$，则由引理 \circled{3} $\chi_H$ 具有性质 $\mathcal F$.}
\handline{72}{601}{760}{$\therefore\quad0=\int_{\R^{n+m}}\chi_H(x,y)\,dxdy=\int_{\R^n}\bigl(\int_{\R^m}\chi_H(x,y)\,dy\bigr)dx$.}
\handline{72}{630}{735}{于是 $h(x)=\int_{\R^m}\chi_H(x,y)\,dy$ 几乎处处为零（在 $\R^n$ 上）.}
\handline{72}{659}{760}{$\therefore$ a.e. $x\in\R^n$，$\{y\in\R^m\mid(x,y)\in H\}$ 在 $\R^m$ 中零测.}
\handline{72}{687}{745}{而 $E_x\subseteq\{y\in\R^m\mid(x,y)\in H\}$，$\therefore E_x$ 在 $\R^m$ 中零测.}
\handline{72}{715}{770}{最后，验证 $\chi_E$ 具有性质 $\mathcal F$.}
\handline{72}{746}{770}{\circled{1} a.e. $x\in\R^n$，$\chi_E(x,y)=\chi_{E_x}(y)$.}
\handline{72}{777}{740}{由 $m(E_x)=0\Rightarrow\int_{\R^m}\chi_{E_x}(y)\,dy=0$.}
\handline{72}{807}{765}{$\therefore\chi_E(x,y)$ 在 $\R^m$ 上可测，且 $\in L(dy)$.}
\handline{72}{837}{785}{\circled{2} $h(x)=\int_{\R^m}\chi_E(x,y)\,dy$ 几乎处处为零，故可测，且 $\in L(dx)$.}
\handline{72}{868}{790}{\circled{3} $\int_{\R^{n+m}}\chi_E(x,y)\,dxdy=0=\int_{\R^n}\bigl(\int_{\R^m}\chi_E(x,y)\,dy\bigr)dx$.}
\handline{72}{924}{770}{$\therefore\chi_E$ 具有性质 $\mathcal F$.}
\end{handpage}

% SOURCE PAGE 050
\begin{handpage}{50}
\handline{72}{128}{780}{Lemma \circled{5}. $E\subseteq\R^{n+m}$，$E\in\mathcal M$，$m(E)<\infty$，则 $\chi_E$ 具有性质 $\mathcal F$.}
\handline{72}{157}{780}{Pf. $\exists H$ 为 $G_\delta$ 集 s.t. $m(E)=m(H)<\infty$，$E\subseteq H$.}
\handline{72}{185}{710}{设 $Z=H\setminus E$，则 $E=H\setminus Z$，$m(Z)=0$.}
\handline{72}{214}{720}{$\therefore\chi_E=\chi_H-\chi_Z$ 具有性质 $\mathcal F$.}
\handline{72}{240}{790}{Thm \ann{(Fubini)}}
\handline{72}{269}{800}{Pf. $f\in L(dxdy)$，不妨设 $f\ge0$.（否则作分解 $f=f^+-f^-$，分别考虑.）}
\handline{72}{298}{770}{$\exists$ 非负简单函数列，s.t. $S_k\uparrow f$，且 $S_k\in L(dxdy)$.}
\handline{72}{328}{740}{而简单函数 $S_k=\sum_{j=1}^{n_k}a_{kj}\chi_{E_{kj}}$.}
\handline{72}{358}{740}{而 $E_{kj}$（$1\le j\le n_k$）两两不交且 $m(E_{kj})<\infty$.}
\handline{72}{387}{750}{由引理 \circled{5} $\Rightarrow\chi_{E_{kj}}$ 具有性质 $\mathcal F$.}
\handline{72}{415}{735}{$\therefore S_k$ 具有性质 $\mathcal F$.}
\handline{72}{443}{745}{$\therefore$ 由引理 \circled{2} $\Rightarrow f$ 具有性质 $\mathcal F$. Q.E.D.}
\handline{72}{501}{840}{Thm. (i) $f$ 在 $\R^{n+m}$ 上可测，则 a.e. $x\in\R^n$，$f(x,y)$ 在 $\R^m$ 上可测.}
\handline{72}{533}{800}{(ii) $E\subseteq\R^{n+m}$ 可测，则 a.e. $x\in\R^n$，$E_x=\{y\mid(x,y)\in E\}$ 为 $\R^m$ 上可测集.}
\handline{72}{562}{840}{Pf. 先证 (ii)：$E$ 可测，$\therefore\exists F_\sigma$ 集 $F\subseteq E$，s.t. $E=F\cup Z$，$m(Z)=0$.}
\handline{72}{592}{760}{$\therefore E_x=F_x\cup Z_x$，而 $F_x$ 仍是 $F_\sigma$ 集，且对 a.e. $x\in\R^n$，$m(Z_x)=0$.}
\handline{72}{622}{745}{故 $E_x$ 为 $\R^m$ 上可测集（a.e. $x\in\R^n$）.}
\handline{72}{653}{780}{再证 (i) 由于 $f$ 在 $\R^{n+m}$ 上可测.}
\handline{72}{686}{790}{故 $\forall a\in\Q$，$E_a=\{(x,y)\mid f(x,y)>a\}$ 是 $\R^{n+m}$ 中的可测集.}
\handline{72}{718}{785}{由 (ii)，对 a.e. $x\in\R^n$，$E_{a,x}=\{y\mid(x,y)\in E_a\}$ 是 $\R^m$ 上可测集.}
\handline{72}{750}{780}{设 $Z_a$ 为使得 $E_{a,x}$ 不可测的 $x$ 的全体，则 $Z_a\subseteq\R^n$ 为零测集.}
\handline{72}{781}{785}{令 $Z=\bigcup_{a\in\Q}Z_a$，则 $Z\subseteq\R^n$ 为零测集.}
\handline{72}{812}{810}{$\forall x\notin Z:\Rightarrow\forall a\in\Q$，$x\notin Z_a$，$\therefore E_{a,x}=\{y\mid(x,y)\in E_a\}=\{y\mid f(x,y)>a\}$ 可测.}
\handline{72}{843}{785}{由于 $\overline\Q=\R$，$\therefore f(x,y)$ 作为 $y$ 的函数在 $\R^m$ 上可测.}
\handline{72}{876}{785}{$\therefore$ a.e. $x\in\R^n$，$f(x,y)$ 在 $\R^m$ 上可测.}
\end{handpage}

% SOURCE PAGE 051
\begin{handpage}{51}
\handline{72}{139}{856}{Cor. $E\subseteq\R^{n+m}$ 为可测集，$f$ 在 $E$ 上可测，则}
\handline{72}{167}{800}{(i) a.e. $x\in\R^n$，$f(x,y)$ 在 $E_x$ 上可测.}
\handline{72}{195}{810}{(ii) 若 $f(x,y)\in L(E)$，则 a.e. $x\in\R^n$，$f(x,y)\in L(E_x)$.}
\handline{72}{224}{800}{且 $h(x)=\int_{E_x}f(x,y)\,dy\in L(dx)$，$\int_E f(x,y)\,dxdy=\int_{\R^n}\bigl(\int_{E_x}f(x,y)\,dy\bigr)dx$.}
\handtext{72}{253}{856}{9}{black}{Pf. (i) 令 $\bar f(x,y)=\begin{cases}f(x,y)&(x,y)\in E,\\0&(x,y)\notin E.\end{cases}$ \ann{（即 $\bar f=f\cdot\chi_E$）}}
\handline{72}{280}{780}{则 $\bar f$ 是 $\R^{n+m}$ 上可测函数.}
\handline{72}{309}{800}{$\therefore$ a.e. $x\in\R^n$，$\bar f(x,y)$ 在 $\R^m$ 上可测，$E_x$ 在 $\R^m$ 上为可测集.}
\handline{72}{338}{760}{而 $f(x,y)|_{E_x}=\bar f(x,y)|_{E_x}$.}
\handline{72}{366}{780}{$\therefore$ a.e. $x\in\R^n$，$f(x,y)$ 在 $E_x$ 上可测.}
\handline{72}{394}{850}{(ii) $f(x,y)\in L(E)\Rightarrow\bar f(x,y)\in L(\R^{n+m})$.}
\handline{72}{423}{800}{$\therefore$ 由 Fubini 定理}
\handline{72}{453}{760}{$\int_{\R^n}\bigl(\int_{\R^m}\bar f(x,y)\,dy\bigr)dx=\int_{\R^{n+m}}\bar f(x,y)\,dxdy<\infty$.}
\handline{72}{483}{800}{$\therefore h(x)=\int_{\R^m}\bar f(x,y)\,dy=\int_{E_x}f(x,y)\,dy\in L(dx)$.}
\handline{72}{514}{810}{而 $\int_E f(x,y)\,dxdy=\int_{\R^{n+m}}\bar f(x,y)\,dxdy=\int_{\R^n}\bigl(\int_{\R^m}\bar f(x,y)\,dy\bigr)dx=\int_{\R^n}\bigl(\int_{E_x}f(x,y)\,dy\bigr)dx$.}
\handline{72}{544}{770}{而由 $h(x)\in L(dx)$}
\handline{72}{573}{755}{$\Rightarrow h(x)<\infty$，a.e. $x\in\R^n$.}
\handline{72}{604}{755}{$\therefore$ 即 $\int_{E_x}f(x,y)\,dy<\infty$，a.e. $x\in\R^n$.}
\handline{72}{636}{700}{$\therefore$ a.e. $x\in\R^n$，$f(x,y)\in L(E_x)$.}
\end{handpage}

% SOURCE PAGE 052
% SECTION: 6.2 | Tonelli's Theorem
\begin{handpage}{52}
\handtitle{72}{107}{800}{\S5.2 Tonelli's Theorem}
\handline[AnnotationBlue]{72}{135}{820}{Fubini 定理表明，若 $f\in L^1(dxdy)$ 则 $f$ 具有性质 $\mathcal F$.}
\handline[AnnotationBlue]{72}{163}{820}{反之，若 $f$ 具有性质 $\mathcal F$，则 $f\in L^1(dxdy)$ 吗？}
\handline[NoteRed]{72}{191}{140}{Ex.}
\handdiagram{686}{192}{242}{156}{
\draw[Teal,thick] (0,0) rectangle (4,4);\draw[Teal] (0,2)--(4,2) (2,0)--(2,4) (0,1)--(2,1) (1,0)--(1,2) (2,3)--(4,3) (3,2)--(3,4) (2,2.5)--(3,2.5) (2.5,2)--(2.5,3) (3,3.5)--(4,3.5) (3.5,3)--(3.5,4);
\node[below left] at(0,0){$0$};\node[below] at(4,0){$1$};\node[left] at(0,4){$1$};
\foreach \x/\y/\t in {.5/.5/11,1.5/.5/12,1.5/1.5/13,.5/1.5/14,2.25/2.25/21,2.75/2.25/22,2.75/2.75/23,2.25/2.75/24}{\node[NoteRed,font=\tiny] at(\x,\y){$I_{\t}$};}
\node[AnnotationBlue] at(1,3){$0$};\node[AnnotationBlue] at(3,1){$0$};\node[AnnotationBlue,font=\tiny] at(2.5,3.5){$0$};\node[AnnotationBlue,font=\tiny] at(3.5,2.5){$0$};\node[AnnotationBlue,font=\tiny] at(3.25,3.75){$0$};\node[AnnotationBlue,font=\tiny] at(3.75,3.25){$0$};\node[NoteRed,font=\tiny] at(3.25,3.25){$I_3$};\node[NoteRed,font=\tiny] at(3.75,3.75){$\cdots$};
}
\handline{72}{222.617}{490}{$f(x,y)=\begin{cases}\frac1{m(I_{kj})}&j=1,3,\\-\frac1{m(I_{kj})}&j=2,4\end{cases}\quad((x,y)\in\operatorname{int}(I_{kj})), $}
\handline{72}{254}{370}{$0\qquad\text{其他}.$}
\handline{72}{294}{500}{则 $\forall x,y\in[0,1]$，$\int_0^1f(x,y)\,dx=\int_0^1f(x,y)\,dy=0$.}
\handline{72}{330}{490}{于是 $\int_0^1\bigl(\int_0^1 f(x,y)\,dx\bigr)dy=\int_0^1\bigl(\int_0^1f(x,y)\,dy\bigr)dx=0$.}
\handline{72}{390}{780}{但是 $\int_I|f(x,y)|\,dxdy=\sum_{k=1}^{+\infty}1=+\infty$，$\therefore f\notin L(I)$.}
\handline[NoteRed]{72}{421}{820}{故 Fubini 定理的逆命题不成立.}
\handline{72}{451}{850}{Thm. \ann{(Tonelli)} \redann{（将 Fubini 中 $f\in L^1(dxdy)$ 的条件换成可测且恒正）}}
\handline{72}{482}{825}{$f$ 在 $\R^{n+m}$ 上可测且 $f\ge0$，则 $\int_{\R^{n+m}}f(x,y)\,dxdy=\int_{\R^n}\bigl(\int_{\R^m}f(x,y)\,dy\bigr)dx$.}
\handline{72}{513}{855}{Pf. $f$ 在 $\R^{n+m}$ 上可测，$\therefore$ a.e. $x\in\R^n$，$f(x,y)$ 在 $\R^m$ 上可测.}
\handline{72}{541}{565}{a.e. $y\in\R^m$，$f(x,y)$ 在 $\R^n$ 上可测.}
\handline{72}{579}{765}{令 $f_k(x,y)=\begin{cases}\min(k,f(x,y))&|(x,y)|\le k,\\0&|(x,y)|>k.\end{cases}$ 则 $f_k\in L(dxdy)$ 且 $f_k\uparrow f$.}
\handline{72}{612}{775}{$\therefore$ 由 Fubini 定理和 Levi 引理，}
\handline{72}{649}{720}{$\int_{\R^{n+m}}f_k(x,y)\,dxdy\uparrow\int_{\R^{n+m}}f(x,y)\,dxdy$,}
\handline{72}{684}{725}{$\int_{\R^n}\bigl(\int_{\R^m}f_k(x,y)\,dy\bigr)dx\uparrow\int_{\R^n}\bigl(\int_{\R^m}f(x,y)\,dy\bigr)dx$.}
\handline{72}{715}{750}{$\therefore\int_{\R^{n+m}}f(x,y)\,dxdy=\int_{\R^n}\bigl(\int_{\R^m}f(x,y)\,dy\bigr)dx$.}
\handline[Teal]{72}{746}{856}{推广 Fubini：若 $f$ 可测，且 $\int_{\R^{n+m}}f$ 存在（不一定有限）令 $f=f^+-f^-$.}
\handline[Teal]{72}{780}{820}{则 $f^+,f^-\ge0$，可测且 $\int_{\R^{n+m}}f^+$，$\int_{\R^{n+m}}f^-$ 不能同时为 $+\infty$.}
\handline[Teal]{72}{810}{815}{于是由 Tonelli 和积分线性性}
\handline[Teal]{72}{850}{760}{$\Rightarrow\quad\int_{\R^{n+m}}f(x,y)\,dxdy=\int_{\R^n}\bigl(\int_{\R^m}f(x,y)\,dy\bigr)dx$.}
\handline[NoteRed]{72}{882}{856}{Tonelli 由 Fubini 推出，又反过来推广 Fubini.}
\handline[NoteRed]{72}{922}{810}{只要 $\int_{\R^{n+m}}f(x,y)\,dxdy$ 存在就可以换序.}
\end{handpage}

% SOURCE PAGE 053
% SECTION: 6.3 | Applications of Fubini's Theorem
\begin{handpage}{53}
\handtitle{72}{108}{850}{\S5.3 Applications of Fubini's Theorem}
\handline{72}{137}{840}{一. 定理补遗}
\handline{72}{167}{856}{Thm. $E\subseteq\R^n$，可测，$f\ge0$，则 $f$ 在 $E$ 上可测 $\Leftrightarrow\int_E f$ 存在 \ann{（即 $R(f,E)$ 可测）}}
\handline{72}{196}{850}{Pf. 在 \S4.1 中已证 ``$\Rightarrow$''，下证 ``$\Leftarrow$''.}
\handline{72}{224}{790}{若 $R(f,E)$ 是 $\R^{n+1}$ 中的可测集.}
\handline{72}{253}{790}{\circled{1} $y<0$，则 $\{x\in E\mid f(x)\ge y\}=E$，可测.}
\handline{72}{281}{800}{\circled{2} $0\le y<\infty$，$\{x\in E\mid f(x)\ge y\}=\{x\in E\mid(x,y)\in R(f,E)\}=E\cap R(f,E)_y$.}
\handline{72}{311}{780}{由于 $R(f,E)$ 可测，故 a.e. $0\le y<\infty$，$R(f,E)_y$ 可测.}
\handline{72}{341}{790}{$\therefore\{x\in E\mid f(x)\ge y\}$ 可测，a.e. $0\le y<\infty$.}
\handline{72}{372}{800}{设 $G$ 为使得 $\{x\in E\mid f(x)\ge y\}$ 可测的 $y$ 的全体.}
\handline{72}{401}{795}{则 $\overline G$ 在 $\R$ 中稠密，$\therefore\{x\in E\mid f(x)\ge y\}$ 可测（$\forall y\in\R$）.}
\handline{72}{432}{800}{$\therefore f(x)$ 在 $E$ 上可测.}
\handline[Teal]{72}{461}{800}{零测集的补集一定稠密.}
\handline{72}{497}{856}{二. 卷积 \ann{(Convolution)}}
\handline{72}{526}{856}{Def. $f,g:\R^n\to\R$.}
\handline{72}{557}{810}{$f*g(x):=\int_{\R^n}f(x-t)g(t)\,dt$ 称为 $f$ 和 $g$ 的卷积.}
\handline[Teal]{72}{590}{810}{$g*f(x)=\int_{\R^n}f(t)g(x-t)\,dt\ \overset{\text{变量代换}}{=}\ \int_{\R^n}f(x-t)g(t)\,dt=f*g(x)$.}
\handline{72}{622}{856}{Lemma. $f(x)$ 在 $\R^n$ 上可测，则 $F(x,t)=f(x-t)$ 在 $\R^{2n}$ 上可测.}
\handline{72}{653}{856}{Pf. 设 $g(x,t)=f(x)$，则 $g$ 在 $\R^{2n}$ 上可测.}
\handline{72}{688}{765}{作变量代换 $\begin{cases}x=\xi-\eta\\t=\xi+\eta.\end{cases}$ 这是线性变换 $\Rightarrow$ Lipschitz.}
\handline{72}{754}{770}{由于 Lipschitz 变换保持可测性.}
\handline{72}{792}{770}{故 $f(\xi-\eta)=g(\xi-\eta,\xi+\eta)$ 在 $\R^{2n}$ 中可测.}
\end{handpage}

% SOURCE PAGE 054
\begin{handpage}{54}
\handline{72}{92}{800}{Thm. $f,g\in L^1(\R^n)$ 则 $f*g\in L^1(\R^n)$ 且 $\|f*g\|_1\le\|f\|_1\cdot\|g\|_1$.}
\handline{72}{122}{750}{若 $f,g\ge0$，则 $\|f*g\|_1=\|f\|_1\cdot\|g\|_1$.}
\handline{72}{151}{820}{Pf. $f\ge0$，$g\ge0$，$f*g(x)=\int_{\R^n}f(x-t)g(t)\,dt$.}
\handline{72}{182}{760}{由 Tonelli 定理 $\int_{\R^n}\bigl(\int_{\R^n}f(x-t)g(t)\,dt\bigr)dx=\int_{\R^n}\bigl(\int_{\R^n}f(x-t)\,dx\bigr)g(t)\,dt$}
\handline{72}{212}{640}{$=\int_{\R^n}\bigl(\int_{\R^n}f(x)\,dx\bigr)g(t)\,dt=\int_{\R^n}\|f\|_1g(t)\,dt=\|f\|_1\|g\|_1$.}
\handline{72}{241}{765}{$\therefore\|f*g\|_1=\int_{\R^n}f*g(x)\,dx=\int_{\R^n}\bigl(\int_{\R^n}f(x-t)g(t)\,dt\bigr)dx=\|f\|_1\cdot\|g\|_1$.}
\handline{72}{267}{745}{一般情况：}
\handline{72}{298}{740}{$\|f*g\|_1=\int_{\R^n}|f*g(x)|\,dx=\int_{\R^n}\bigl|\int_{\R^n}f(x-t)g(t)\,dt\bigr|dx$}
\handline{72}{328}{610}{$\le\int_{\R^n}\bigl(\int_{\R^n}|f(x-t)|\,|g(t)|\,dt\bigr)dx=\|f\|_1\cdot\|g\|_1<\infty$}
\handline{72}{360}{750}{$\therefore f*g\in L^1(\R^n)$.}
\handline{72}{389}{830}{三. 闭区间的几何结构}
\handline{72}{422}{850}{Def. $F\subseteq\R$，闭集，$x\in\R$，定义 $\delta(x)=\delta(x,F)=\inf\{|x-y|\mid y\in F\}$.}
\handline[Teal]{72}{449}{780}{首先 $\delta(x)=0\Leftrightarrow x\in F$.}
\handline[Teal]{72}{479}{790}{$\delta(x)=0\Leftrightarrow\forall\varepsilon>0$，$\exists y\in F$ s.t. $|x-y|<\varepsilon\Leftrightarrow\forall\varepsilon>0$，$B(x,\varepsilon)\cap F\ne\varnothing$}
\handline[Teal]{72}{511}{780}{$\Leftrightarrow x$ 是 $F$ 的接触点，即 $x\in\overline F=F$.}
\handline[Teal]{72}{546}{795}{其次 $\forall x,y\in\R$，$|\delta(x)-\delta(y)|\le|x-y|\Rightarrow\delta(x)$ 是 Lipschitz 函数 $\Rightarrow$ 一致连续.}
\handline{72}{580}{840}{Ex.}
\handline{72}{610}{780}{$F=[0,1]\cup[2,4]\cup[6,10]$.}
\handdiagram{72}{635}{779}{112}{
\draw[Teal,thick](-1.6,1.6)--(0,0)--(1,0)--(1.5,.5)--(2,0)--(4,0)--(5,1)--(6,0)--(10,0)--(11.3,1.6);
\draw[black] (0,0)--(10.1,0);\foreach \k in {0,...,10}{\node[below] at(\k,0){$\k$};}
\draw[AnnotationBlue,dashed](1.5,0)--(1.5,.5)(5,0)--(5,1);\node[AnnotationBlue,above right]at(1.5,.5){$.5$};\node[AnnotationBlue,above right]at(5,1){$1$};
\foreach \a/\b in {0/1,2/4,6/10}{\foreach \t in {0,.12,...,1}{\draw[NoteRed]({\a+(\b-\a)*\t},0)--++(.05,.1);}}
}
\handline{72}{782}{770}{$\delta(x)$ 在 $F$ 上取值为 $0$.}
\handline{72}{817}{770}{在 $F$ 之外为斜率是 $\pm1$ 的直线.}
\end{handpage}

% SOURCE PAGE 055
\begin{handpage}{55}
\handline{72}{150}{850}{Thm. \ann{(Marcinkiewicz)}}
\handline{72}{176}{820}{$(a,b)$ 有界开区间，$F\subseteq(a,b)$ 为闭集，$\delta(x)=\delta(x,F)$.}
\handline{72}{206}{825}{$\forall\lambda>0$，令 $M_\lambda(x)=M_\lambda(x,F)=\int_a^b\frac{\delta^\lambda(y)}{|x-y|^{1+\lambda}}\,dy$.}
\handline{72}{232}{825}{则 $M_\lambda(x)\in L^1(F)$ 且 $\int_FM_\lambda\le2\lambda^{-1}m(G)$，其中 $G=(a,b)\setminus F$.}
\handline{72}{260}{856}{Pf.}
\handline{72}{290}{790}{$\int_FM_\lambda=\int_F\bigl(\int_a^b\frac{\delta^\lambda(y)}{|x-y|^{1+\lambda}}\,dy\bigr)dx\ \overset{\text{Tonelli}}{=}\ \int_a^b\bigl(\int_F\frac1{|x-y|^{1+\lambda}}\,dx\bigr)\delta^\lambda(y)\,dy$.}
\handline{72}{340}{790}{$\overset{y\in F\Leftrightarrow\delta(y)=0}{=}\quad\int_G\bigl(\int_F\frac1{|x-y|^{1+\lambda}}\,dx\bigr)\delta^\lambda(y)\,dy$.}
\handline{72}{375}{800}{$\forall y\in G$，$y\notin F$.}
\handdiagram{302}{387}{286}{43}{
\draw[thick](0,0)--(5,0);\node[AnnotationBlue,below]at(0,0){$a$};\node[AnnotationBlue,below]at(5,0){$b$};\node[AnnotationBlue,left]at(0,0){$(\,$};\node[AnnotationBlue,right]at(5,0){$\,)$};\node[NoteRed,below]at(1,0){$y$};\draw[NoteRed](1,-.05)--(1,.1);\draw[NoteRed](1,-.15)--(1,-.25)--(2,-.25)--(2,-.15);\node[NoteRed,below]at(1.5,-.25){$\delta(y)$};\foreach\x in {2,2.15,...,3.5}{\draw[AnnotationBlue](\x,0)--++(.1,.2);}\node[AnnotationBlue,below]at(2.8,0){$F$};
}
\handline{72}{469}{790}{$\int_F\frac1{|x-y|^{1+\lambda}}\,dx\le\int_{|x-y|\ge\delta(y)}\frac1{|x-y|^{1+\lambda}}\,dx=2\int_{\delta(y)}^{+\infty}\frac1{t^{1+\lambda}}\,dt=\frac2\lambda(\delta(y))^{-\lambda}$.}
\handline{72}{523}{800}{$\therefore\int_F M_\lambda(x)\,dx\le\int_G\frac2\lambda(\delta(y))^{-\lambda}\cdot\delta(y)^\lambda\,dy=\frac2\lambda m(G)$.}
\handline{72}{552}{795}{$\therefore M_\lambda\in L^1(F)$ \ann{（从而 $M_\lambda$ 在 $F$ 上几乎处处有限）}}
\handline[NoteRed]{72}{582}{856}{Ex. $M_\lambda$ 在 $F$ 的端点会发散.}
\handline[AnnotationBlue]{72}{613}{820}{$(-1,1)$，$F=[-\frac12,0]$.}
\handline[AnnotationBlue]{72}{653}{820}{$M_\lambda(0)=\int_{-1}^1\frac{\delta^\lambda(y)}{|y|^{1+\lambda}}\,dy\ge\int_{-\frac14}^{\frac14}\frac{\delta^\lambda(y)}{|y|^{1+\lambda}}\,dy=\int_0^{\frac14}\frac{\delta^\lambda(y)}{|y|^{1+\lambda}}\,dy=\int_0^{\frac14}\frac{y^\lambda}{|y|^{1+\lambda}}\,dy$}
\handline[AnnotationBlue]{72}{709}{345}{$=\int_0^{\frac14}\frac1y\,dy=+\infty$.}
\end{handpage}

% SOURCE PAGE 056
% CHAPTER: 7 | Differentiation
% SECTION: 7.1 | The Indefinite Integral
\begin{handpage}{56}
\handtitle{72}{133}{820}{Chapter 6 Differentiation}
\handtitle{72}{164}{815}{\S6.1 The Indefinite Integral}
\handline{72}{192}{850}{Def. 设 $\Sigma$ 是一 $\sigma$-代数，若 $\varphi:\Sigma\to\R$（或 $\C$）满足可数可加性.}
\handline{72}{221}{775}{即 $\varphi(E)=\varphi(\bigcup_{j=1}^\infty E_j)=\sum_{j=1}^\infty\varphi(E_j)$，其中 $E=\bigcup_{j=1}^\infty E_j$，$E_j$ 两两不交 $\in\Sigma$，称 $\varphi$ 是 $\Sigma$ 上的集函数.}
\handline[Teal]{72}{248}{845}{注：\circled{1} $\varphi:\Sigma\to\R$，不包含 $\pm\infty$，故 $\varphi$ 是有限的.}
\handline[Teal]{72}{278}{785}{\circled{2} 由于 $\sum_{j=1}^\infty\varphi(E_j)=\varphi(E)$，故任意重排之后的值相同，从而必须绝对收敛.}
\handline[Teal]{72}{306}{790}{\circled{3} 这样的 $\varphi$ 事实上也定义了广义测度，且它的全变差是有限的.}
\handline{72}{335}{850}{Ex. $f\in L^1(\R^n)$，$E\in\mathcal M$，令 $F(E)=\int_E f$.}
\handline{72}{363}{780}{则 $F$ 是 $\mathcal M$ 上的集函数.}
\handline{72}{394}{856}{Def. 若 $\forall\varepsilon>0$，$\exists\delta>0$，s.t. 当 $\operatorname{diam}(E)<\delta$ 时，$|\varphi(E)|<\varepsilon$，则称 $\varphi$ 是连续的.}
\handline{72}{423}{780}{即当 $\operatorname{diam}(E)\to0$ 时，$\varphi(E)\to0$.}
\handline{72}{452}{856}{Ex. $\varphi:\mathcal M\to\R$，$\varphi(E)=\begin{cases}1&0\in E\\0&0\notin E.\end{cases}$ 不是连续的. 因为当 $\{0\}\subseteq E$ 时 $\varphi(E)=1$，而 $\operatorname{diam}(\{0\})=0$.}
\handline{72}{481}{856}{Def. 若 $\forall\varepsilon>0$，$\exists\delta>0$，s.t. $m(E)<\delta$ 时，$|\varphi(E)|<\varepsilon$，则称 $\varphi$ 关于 Lebesgue 测度绝对连续.}
\handline[Teal]{72}{507}{790}{绝对连续 $\Rightarrow$ 连续 \ann{（$\operatorname{diam}(E)\to0\Rightarrow m(E)\to0$）}}
\handline{72}{542}{850}{Ex. $\Omega=[0,1]\times[0,1]$，$D=\{(x,y)\mid x=y,\ 0\le x\le1\}$.}
\handline{72}{571}{741}{$\Sigma=\{E\mid E\subseteq\Omega,\ E\in\mathcal M_2,\ E\cap D\in\mathcal M_1\}$，$\varphi(E):=m_1(E\cap D)$ \ann{（一维测度）}}
\handline{72}{600}{570}{则 $\varphi$ 是连续的. $\operatorname{diam}(E)\to0$，$m_1(E\cap D)\to0$.}
\handline{72}{631}{570}{但 $\varphi$ 不是绝对连续的，$m_2(E)\to0$，$m_1(E\cap D)\not\to0$.}
\handdiagram{831}{573}{97}{67}{
\draw[Teal,thick](0,0)rectangle(2,2);\draw[Teal,thick](0,0)--(2,2);\draw[AnnotationBlue] (1,1)ellipse(.7 and .5);\node[AnnotationBlue,above left]at(.3,1.5){$E$};\draw[NoteRed,thick](.6,.6)--(1.4,1.4);\draw[NoteRed,->](1.1,1)--(1.2,-.2)node[below]{$E\cap D$};
}
\handline{72}{662}{856}{Thm. $f\in L(A)$，$F(E):=\int_E f$，$E\subseteq A$，$E\in\mathcal M$.}
\handline{72}{691}{780}{则 $F$ 绝对连续.}
\handline{72}{723}{856}{Pf. 不妨设 $f\ge0$，否则考虑 $f=f^+-f^-$. $f\in L(A)\Rightarrow0\le\int_A f<+\infty$.}
\handline{72}{765}{795}{$\therefore\forall\varepsilon>0$，$\exists k_0\in\N$，s.t. $0\le\int_{A\cap\{f>k_0\}}f<\frac12\varepsilon$.}
\handline{72}{794}{775}{当 $E\subseteq A$，且 $m(E)<\frac\varepsilon{2k_0}$ 时，}
\handline{72}{831}{765}{$0\le\int_E f=\int_{E\cap\{f\le k_0\}}f+\int_{E\cap\{f>k_0\}}f\le k_0\cdot m(E)+\frac12\varepsilon<\varepsilon$.}
\handline{72}{861}{780}{$\therefore F$ 绝对连续.}
\end{handpage}

% SOURCE PAGE 057
% SECTION: 7.2 | Lebesgue's Differentiation Theorem
\begin{handpage}{57}
\handtitle{72}{106}{795}{\S6.2 Lebesgue's Differentiation Theorem}
\handline{72}{138}{800}{\textbf{Def.}\quad $\Omega\subseteq\R^n$为开集，$f:\Omega\to\R$ 可测.}
\handline{72}{164}{805}{若$\forall$紧集$K\subseteq\Omega$，$\int_K|f|<+\infty$，则称$f$局部Lebesgue可积，记为$f\in L^1_{\mathrm{loc}}(\Omega)$.}
\handline[Teal]{72}{193}{820}{由于可微性是一个局部概念，故可在$L^1_{\mathrm{loc}}(\Omega)$中研究可微.}
\handline{72}{218}{856}{\textbf{Def.}\quad $Q$：$n$维立方体且边平行坐标轴.\quad $f$是$\R^n$上可测函数. $F(E)$为$f$的不定积分.}
\handline{72}{249}{790}{若$\lim_{Q\downarrow x}\frac{F(Q)}{m(Q)}=f(x)$，则称$F$在$x$处可微且具有导数$f(x)$.}
\handline{72}{279}{856}{\textbf{Thm.}\quad \ann{(Lebesgue's Differentiation Theorem)}}
\handline{72}{305}{850}{$f\in L^1_{\mathrm{loc}}(\R^n)$，则$f$的不定积分在$\R^n$上几乎处处可微，且导数为$f(x)$.}
\handline[Teal]{72}{335}{800}{若$f\in L^1\cap C(\R^n)$}
\handline[Teal]{72}{363}{815}{则$\left|\frac{F(Q)}{m(Q)}-f(x)\right|=\left|\frac1{m(Q)}\int_Qf(y)\,dy-f(x)\right|=\frac1{m(Q)}\left|\int_Q(f(y)-f(x))\,dy\right|$}
\handline[Teal]{72}{391}{700}{$\le\frac1{m(Q)}\int_Q|f(y)-f(x)|\,dx$.}
\handline[Teal]{72}{419}{710}{当$m(Q)\to0$时，$|f(y)-f(x)|\le\varepsilon$，$\forall y\in Q$.}
\handline[Teal]{72}{450}{770}{$\therefore\quad\left|\frac{F(Q)}{m(Q)}-f(x)\right|\le\frac1{m(Q)}\cdot m(Q)\cdot\varepsilon=\varepsilon$.}
\handline[Teal]{72}{483}{810}{于是这个定理对连续的$L^1$函数成立.}
\handline[Teal]{72}{511}{805}{并希望在$L^1$范数下用$C$函数逼近$L^1$函数从而完成证明.}
\handline{72}{548}{856}{\textbf{Thm.}\quad $f\in L^1(\R^n)$，则$f$可以用紧支连续函数在$L^1$范数下逼近.}
\handline{72}{583}{820}{即$\forall\varepsilon>0$，$\exists\varphi_\varepsilon\in C_0(\R^n)$，s.t. $\int_{\R^n}|f-\varphi_\varepsilon|<\varepsilon$.\quad\ann{$\bigl(\overline{C_0(\R^n)}=L^1(\R^n)\bigr)$}.}
\handline{72}{614}{856}{\textbf{Pf.}\quad 为了叙述方便，称具有上述性质的$f$为具有性质$A$.}
\handline{72}{647}{856}{断言：\circled{1}\ 有限个具有性质$A$的函数的线性组合仍具有性质$A$.\ \ann{(trivial)}}
\handline{72}{678}{815}{\circled{2}\ $f_k\in L^1$，且具有性质$A$，若$f_k\to f$\ \ann{(依$L^1$范数)}}
\handline{72}{712}{790}{则$f\in L^1(\R^n)$，且$f$具有性质$A$.}
\handline{72}{742}{760}{不妨设$\|f_k-f\|_1\le1$}
\handline{72}{777}{800}{$\therefore\quad\int_{\R^n}|f|\le\int_{\R^n}|f-f_k|+\int_{\R^n}|f_k|\le1+\int_{\R^n}|f_k|<+\infty$，$\therefore f\in L^1(\R^n)$}
\handline{72}{813}{800}{$\forall\varepsilon>0$，选得一个$f_{k_1}$ s.t. $\|f_{k_1}-f\|_1<\frac\varepsilon2$.}
\handline{72}{847}{780}{对$f_{k_1}$，$\exists\varphi\in C_0(\R^n)$，s.t. $\|f_{k_1}-\varphi\|_1<\frac\varepsilon2$}
\handline{72}{882}{740}{$\therefore\quad\|f-\varphi\|_1\le\|f-f_{k_1}\|_1+\|f_{k_1}-\varphi\|_1<\varepsilon$.}
\end{handpage}

% SOURCE PAGE 058
\begin{handpage}{58}
\handline{72}{102}{775}{有了两个断言之后，回到原定理：}
\handline{72}{132}{740}{不妨设$f\ge0$，否则令$f=f^+-f^-$，分别考虑.}
\handline{72}{162}{740}{$\exists$简单函数$S_k\uparrow f$.}
\handline{72}{191}{730}{$S_k=\sum_{j=1}^{n_k}a_{kj}\chi_{E_{kj}}$，$a_{kj}>0$，由于$f\in L^1$，$\therefore m(E_{kj})<\infty$.}
\handline{72}{217}{740}{由$0\le f-S_k=|f-S_k|\le|f|+|S_k|=2f\in L^1$.}
\handline{72}{245}{740}{$\therefore$由LDCT $\Rightarrow\quad\int_{\R^n}|f-S_k|\to0$，于是$S_k\to f$\ \ann{(依$L^1$范数)}}
\handline{72}{274}{740}{最后由断言\circled{1}\circled{2}，只需证明$\chi_E$（$m(E)<\infty$）具有性质$A$.}
\handline{72}{303}{735}{$\forall\varepsilon>0$，$\exists$开集$G\subseteq\R^n$，s.t. $E\subseteq G$，$m(G\setminus E)<\varepsilon$.}
\handline{72}{331}{700}{而$\int_{\R^n}|\chi_G-\chi_E|=m(G\setminus E)<\varepsilon$.}
\handline{72}{360}{730}{故只需证明$\chi_G$，（$m(G)<\infty$，$G$开集），具有性质$A$.}
\handline{72}{391}{730}{而$G=\bigcup_{j=1}^\infty I_j$，$I_j$：两两不交，半开区间. $m(I_j)<\infty$.}
\handline{72}{422}{730}{$\therefore\chi_G=\sum_{j=1}^\infty\chi_{I_j}=\lim_{N\to\infty}\sum_{j=1}^N\chi_{I_j}$.}
\handline{72}{452}{735}{故只需证明$\chi_I$（$I$：半开区间，$m(I)<\infty$）具有性质$A$.}
\handline{72}{484}{725}{$\exists I^*$：闭区间，s.t. $\overline I\subseteq\operatorname{int}(I^*)$，且$m(I^*\setminus I)<\varepsilon$.}
\handline{72}{516}{730}{$\therefore$取$\varphi\in C(\R^n)$，且$\varphi|_I\equiv1$ 且$\supp(\varphi)\subseteq\operatorname{int}(I^*)$}
\handline{72}{546}{675}{则$\int_{\R^n}|\chi_I-\varphi|\le m(I^*\setminus I)<\varepsilon$\qquad Q.E.D.}
\handline[Teal]{72}{606}{815}{支集：$\varphi:\R^n\to\R$，$\supp(\varphi):=\overline{\{x\in\R^n\mid\varphi(x)\ne0\}}$}
\handline[Teal]{72}{638}{720}{即$\varphi$不为零的点构成集合的闭包.}
\handline[Teal]{72}{670}{715}{函数的支集一定是闭集}
\handline[Teal]{72}{703}{830}{紧支函数：在$\R^n$的一个有界集之外，函数值为$0$.}
\handline[Teal]{72}{737}{705}{$\Longleftrightarrow\quad\supp(\varphi)$有界}
\handline[Teal]{72}{771}{710}{由于$\supp(\varphi)$本身是闭集，故还等价于$\supp(\varphi)$是紧集.}
\handline[NoteRed]{72}{805}{760}{此定理中的$C_0(\R^n)$还可以加强为$C_0^\infty(\R^n)$.}
\end{handpage}

% SOURCE PAGE 059
\begin{handpage}{59}
\handline{72}{120}{856}{\textbf{Thm.}\quad (Vitali Covering Lemma)}
\handline{72}{147}{856}{$E\subseteq\R^n$，$m^*(E)<\infty$，$K$是某些闭立方体构成的族，且$E\subseteq\bigcup_{Q\in K}Q$.}
\handline{72}{176}{856}{则$\exists$一个正数$\beta=\beta(n)$，以及$K$中有限个互不相交的立方体$Q_1,Q_2,\ldots,Q_N$.}
\handline{72}{204}{835}{s.t.\quad $\sum_{j=1}^Nm(Q_j)\ge\beta m^*(E)$.}
\handline{72}{234}{856}{\textbf{Pf.}\quad $E\subseteq\bigcup_{Q\in K}Q$，设$K=\{Q(t)\mid t\text{为}Q\text{的边长}\}:=K_1$.}
\handline{72}{259}{795}{设$t_1^*=\sup\{t\mid Q(t)\in K_1\}$.}
\handline{72}{287}{810}{\circled{1}\ $t_1^*=+\infty$，用一个$t$足够大的$Q$，即可. $m(Q)=t^n>\beta m^*(E)$.}
\handline{72}{314}{780}{$\beta$取任意正数，$t$都可以取到.}
\handline{72}{343}{810}{\circled{2}\ $t_1^*<+\infty$.}
\handline{72}{372}{790}{选取$Q_1=Q_1(t_1)$，s.t. $t_1>\frac{t_1^*}2$.}
\handline{72}{402}{790}{设$K_2=\{Q(t)\cap Q_1=\varnothing\}$，$K'_2=\{Q(t)\cap Q_1\ne\varnothing\}$，则$K_1=K_2\cup K'_2$.}
\handline{72}{430}{790}{设$Q_1^*$为与$Q_1$同心，且边长为$5t_1$的立方体.}
\handline{72}{459}{785}{若$Q(t)\in K'_2$，则$t\le t_1^*<2t_1<\frac52t_1\Rightarrow Q(t)\subseteq Q_1^*$.}
\handline{72}{488}{785}{设$t_2^*=\sup\{t\mid Q(t)\in K_2\}$.}
\handline{72}{518}{780}{选取$Q_2=Q_2(t_2)\in K_2$，s.t. $t_2>\frac{t_2^*}2$}
\handline{72}{550}{780}{同理\quad$K_2=K_3\cup K'_3$，}
\handline{72}{580}{785}{$K_3=\{Q(t)\in K_2\mid Q(t)\cap Q_2=\varnothing\}$，$K'_3=\{Q(t)\in K_2\mid Q(t)\cap Q_2\ne\varnothing\}$.}
\handline{72}{610}{780}{将$Q_2$同心放大$5$倍为$Q_2^*$.}
\handline{72}{642}{780}{若$Q(t)\in K'_3$，则$Q(t)\subseteq Q_2^*$.}
\handline{72}{676}{810}{$\cdots\cdots$重复上述过程，$K_1=K_2\cup K'_2=K_3\cup K'_3\cup K'_2=\cdots=K_{j+1}\cup K'_{j+1}\cup\cdots\cup K'_2\cdots$}
\handline{72}{708}{800}{(I)若$\exists j\in\N$，s.t. $K_{j+1}=\varnothing$.}
\handline{72}{741}{795}{$\forall Q\in K_1=K'_{j+1}\cup K'_j\cup\cdots\cup K'_2$}
\handline{72}{774}{800}{根据构造，$Q_1,Q_2,\ldots,Q_j$两两不交.}
\handline{72}{806}{790}{且$Q\subseteq Q_1^*\cup Q_2^*\cdots\cup Q_j^*$.}
\handline{72}{839}{800}{$\therefore E\subseteq\bigcup_{Q\in K}Q\subseteq Q_1^*\cup Q_2^*\cdots\cup Q_j^*$.}
\handline{72}{874}{800}{$\therefore m^*(E)\le\sum_{j=1}^Nm(Q_j^*)=\sum_{j=1}^N5^n m(Q_j)$}
\handline{72}{910}{780}{则$\beta=\beta(n)=\frac1{5^n}$.}
\end{handpage}

% SOURCE PAGE 060
\begin{handpage}{60}
\handline{72}{122}{790}{(II)\ $\forall j\in\N$，$K_j\ne\varnothing$.}
\handline{72}{153}{745}{则$t_1^*\ge t_2^*\ge\cdots\ge t_j^*\ge\cdots\ge0$.}
\handline{72}{182}{750}{(i)\ $t_j^*\downarrow\delta$，$\delta>0$.}
\handline{72}{211}{700}{对于$Q_j=Q_j(t_j)$，$t_j>\frac{t_j^*}2\ge\frac\delta2$，$\therefore m(Q_j)=t_j^n>(\frac\delta2)^n$.}
\handline{72}{242}{730}{于是$\sum_{j=1}^Nm(Q_j)>N\cdot(\frac\delta2)^n\to+\infty$（$N\to+\infty$）.}
\handline{72}{270}{730}{$\therefore\forall\beta>0$，$\exists N\in\N$，s.t. $\sum_{j=1}^Nm(Q_j)>\beta m^*(E)$.}
\handline{72}{299}{750}{(ii)\ $t_j^*\downarrow0$.}
\handline{72}{328}{750}{断言：$\forall Q\in K_1$，$\exists j\in\N$，s.t. $Q\subseteq Q_j^*$. 从而$Q\subseteq\bigcup_{j=1}^\infty Q_j^*$}
\handline{72}{357}{745}{若否，$\forall j\in\N$，$Q\nsubseteq Q_j^*$，$\Rightarrow Q=Q(t)\in K_{j+1}$}
\handline{72}{386}{730}{由$t_{j+1}^*$的定义知，$t\le t_{j+1}^*$（$j\in\N$）.}
\handline{72}{414}{710}{有$t_j^*\downarrow0\Rightarrow t\le0$与$t>0$矛盾!}
\handline{72}{444}{750}{$\therefore E\subseteq\bigcup_{Q\in K}Q\subseteq\bigcup_{j=1}^\infty Q_j^*$}
\handline{72}{475}{750}{$\therefore m^*(E)\le\sum_{j=1}^\infty5^n m(Q_j)\Rightarrow\frac1{5^n}m^*(E)\le\sum_{j=1}^\infty m(Q_j)$.}
\handline{72}{505}{730}{\circled{1}\ 若$\sum_{j=1}^\infty m(Q_j)=+\infty$，则$\forall\beta>0$，$\exists N\in\N$，s.t. $\beta m^*(E)<\sum_{j=1}^Nm(Q_j)$.}
\handline{72}{536}{735}{\circled{2}\ 若$\sum_{j=1}^\infty m(Q_j)<+\infty$，则$\forall\beta\in(0,\frac1{5^n})$，$\exists N\in\N$，s.t. $\beta m^*(E)<\sum_{j=1}^Nm(Q_j)$.}
\handline[Teal]{72}{568}{856}{注：\circled{1}\ $E$没有要求可测，只有外测度有限这一个条件.}
\handline[Teal]{72}{601}{805}{\circled{2}\ $\beta$只依赖于维数$n$，不依赖于具体的$E,K$.}
\handline[Teal]{72}{632}{805}{\circled{3}\ $Q_1,\ldots,Q_N$为有限个从$K$中抽出的立方体且要求两两不交}
\handline[Teal]{72}{669}{805}{\circled{4}\ 这个结果表明$m^*(E)$可以被$Q_1,\ldots,Q_N$的“体积”控制.}
\handline[Teal]{72}{703}{805}{\circled{5}\ 这个结果只对欧氏空间和Lebesgue测度成立，因为它具有良好的几何特性.}
\handline[Teal]{72}{737}{780}{比如$m(Q_j^*)=5^n m(Q_j)$，这是普遍的测度不具有的性质.}
\handline[Teal]{72}{771}{790}{在一般的测度空间中也有类似结果：\ann{Besicovich Covering Lemma}}
\handline[Teal]{72}{804}{805}{\circled{6}\ $K$中“闭立方体”改成“开立方体”结论一样成立}
\end{handpage}

% SOURCE PAGE 061
\begin{handpage}{61}
\handline{72}{133}{856}{若$E\subseteq\R^n$是可测集，且$m(E)<\infty$，则证明过程可以极大地简化.}
\handline{72}{163}{856}{设$K$中立方体都是开立方体.}
\handline{72}{193}{850}{$\forall\varepsilon>0$，$\exists$闭集$F\subseteq E$，s.t. $m(E\setminus F)<\varepsilon$.}
\handline{72}{224}{830}{取$\varepsilon=\frac12m(E)$，则$m(F)\ge\frac12m(E)$.}
\handline{72}{263}{855}{$F\subseteq E\subseteq\bigcup_{Q\in K}Q$，由Heine--Borel定理知，$F\subseteq\bigcup_{j=1}^mR_j$}
\handline{72}{294}{830}{在$R_1,\ldots,R_m$中，选边长最大的记为$Q_1$.}
\handline{72}{325}{830}{删去$Q_1$，在剩下的立方体中选一个与$Q_1$不交，且边长最大的，记为$Q_2$.}
\handline{72}{356}{825}{若剩下的每一个都与$Q_1$相交，则在这一步停止.}
\handline{72}{391}{830}{重复上述步骤，有限次之后必定停止.}
\handline{72}{423}{830}{于是得到$Q_1,\ldots,Q_N$，两两不交. 且$N\le m$.}
\handline{72}{457}{825}{若$R_i\notin\{Q_1,\ldots,Q_N\}$，则$\exists1\le j\le N$，s.t. $R_i$与$Q_j$相交.}
\handline{72}{491}{820}{由于删去$R_i$时，$Q_j$是剩下的立方体中边长最大的.}
\handline{72}{527}{820}{于是$R_i$的边长$\le Q_j$的边长.\quad $\Rightarrow R_i\subseteq5Q_j=Q_j^*$.}
\handline{72}{565}{815}{$\therefore F\subseteq\bigcup_{j=1}^mR_j\subseteq\bigcup_{j=1}^NQ_j^*$.}
\handline{72}{600}{825}{$\therefore\frac12m(E)\le m(F)\le\sum_{j=1}^Nm(Q_j^*)=5^n\sum_{j=1}^Nm(Q_j)$}
\handline{72}{636}{815}{$\therefore\frac12\cdot\frac1{5^n}m(E)\le\sum_{j=1}^Nm(Q_j)$.}
\handline{72}{675}{780}{$\beta=\beta(n)=\frac12\cdot\frac1{5^n}$.}
\end{handpage}

% SOURCE PAGE 062
\begin{handpage}{62}
\handline{72}{106}{835}{\textbf{Def.}\quad $f\in L^1_{\mathrm{loc}}(\R^n)$.}
\handline{72}{136}{770}{$f^*(x):=\sup_Q\frac1{m(Q)}\int_Q|f(y)|\,dy$，$Q$是以$x$为中心的立方体.}
\handline{72}{163}{765}{$f^*$称为Hardy--Littlewood极大函数\ \ann{(Hardy--Littlewood maximal function)}}
\handline{72}{194}{840}{\textbf{Prop.}\quad (i)\ $0\le f^*\le+\infty$}
\handline{72}{223}{775}{(ii)\ $(f+g)^*(x)\le f^*(x)+g^*(x)$}
\handline{72}{250}{775}{(iii)\ $(cf)^*(x)=|c|f^*(x)$}
\handline{72}{277}{775}{(iv)\ $f^*$是下半连续函数，从而$f^*$可测.}
\handline{72}{307}{850}{\textbf{Pf.}\quad 只证(iv).\quad $f^*$是l.s.c. $\Leftrightarrow\forall\alpha>0$，$\{x\mid f^*(x)>\alpha\}$是开集.}
\handline{72}{338}{720}{$\forall x_0\in\{x\mid f^*(x)>\alpha\}$.}
\handline{72}{367}{720}{$f^*(x_0)=\sup_Q\frac1{m(Q)}\int_Q|f(y)|\,dy>\alpha$.}
\handline{72}{397}{725}{$\therefore\exists Q_0$，s.t. $\frac1{m(Q_0)}\int_{Q_0}|f(y)|\,dy>\alpha$.}
\handline{72}{425}{710}{考虑$x_0$的邻域中的$x$，$Q_x=Q_0+(x-x_0)$.}
\handline{72}{453}{710}{于是$m(Q_x)=m(Q_0)$.}
\handline{72}{483}{710}{$f^*(x)\ge\frac1{m(Q_x)}\int_{Q_x}|f|=\frac1{m(Q_0)}\int_{Q_x}|f|$.}
\handline{72}{513}{715}{由$f\in L^1_{\mathrm{loc}}(\R^n)$，$\Rightarrow\int_E|f|$绝对连续，}
\handline{72}{546}{720}{$\therefore\frac1{m(Q_0)}\int_{Q_x}|f|\to\frac1{m(Q_0)}\int_{Q_0}|f|=f^*(x_0)>\alpha$.\quad ($x\to x_0$).}
\handline{72}{577}{720}{$\therefore\exists\delta>0$，s.t. $|x-x_0|<\delta$，$\frac1{m(Q_0)}\int_{Q_x}|f|>\alpha+\frac12(f^*(x_0)-\alpha)>\alpha$.}
\handline{72}{609}{695}{$\Rightarrow\quad\forall x\in B(x_0,\delta)$，$f^*(x)>\alpha$.}
\handline{72}{641}{710}{$\therefore B(x_0,\delta)\subseteq\{x\mid f^*(x)>\alpha\}$，故为开集}
\handline{72}{672}{710}{$\therefore f^*$是下半连续函数.}
\handline{72}{706}{856}{\textbf{Ex.}\quad $E\in\mathcal M$，$\chi_E^*(x)=\sup_Q\frac1{m(Q)}\int_Q|\chi_E|=\sup_Q\frac{m(E\cap Q)}{m(Q)}$}
\handline{72}{740}{673}{若$E$有界，$|x|$很大，考虑两个特殊的$Q$，$Q_1$：刚接触$E$，$Q_2$：即将离开$E$.}
\handline{72}{771}{673}{设$Q_1=Q_1(t_1)$，$Q_2=Q_2(t_2)$，$Q=Q(t)$.}
\handline{72}{802}{673}{当$0\le t<t_1$时，$E\cap Q=\varnothing$，$m(E\cap Q)=0$；当$t>t_2$时，$E\cap Q=E$，$m(E\cap Q)=m(E)$}
\handline{72}{843}{673}{$\Rightarrow0\le t<t_1$时，$\frac1{m(Q)}\int_Q|f|=0$，$t>t_2$时，$\frac1{m(Q)}\int_Q|f|$随$t$增大而减小.}
\handline{72}{876}{673}{$\therefore$最大值出现在$t_1\le t\le t_2$时，设$t_1(x),t_2(x)$为$x$的函数，}
\handline{72}{906}{673}{则$t_2(x)-t_1(x)\le\operatorname{diam}(E)<+\infty$.}
\handdiagram{763}{747}{165}{174}{
 \draw[AnnotationBlue] (1,0.9)--(1,4.2); \draw[AnnotationBlue] (1,0.9)--(3.35,0.9);
 \draw[AnnotationBlue,dashed] (1,2.85)--(2.85,2.85); \fill (2.85,2.85) circle(.035); \node[right] at(2.85,2.85){$x$};
 \draw[NoteRed] (.15,.05)--(.15,4.4); \draw[NoteRed] (.15,.05)--(3.25,.05);
 \draw[NoteRed,dashed] (2.85,2.85)--(2.85,.05);
 \draw plot[smooth cycle] coordinates{(.3,1.9)(.7,2.5)(1,2.4)(1,1.5)(.75,1)(.4,1.3)}; \node at(.6,1.65){$E$};
 \draw[AnnotationBlue,->] (1,3.7) to[bend left] (1.5,4.0); \node[AnnotationBlue,right] at(1.5,4.0){$Q_1$};
 \draw[NoteRed,->] (.85,0) --(.7,-.5); \node[NoteRed,below] at(.7,-.5){$Q_2$};
 \draw[AnnotationBlue] plot[smooth] coordinates{(1.2,2.6)(1.3,2.45)(1.7,2.45)(1.9,2.3)(2.1,2.45)(2.55,2.45)(2.65,2.6)}; \node[AnnotationBlue,below] at(1.9,2.55){$\frac{t_1}2$};
 \draw[NoteRed] plot[smooth] coordinates{(3.05,2.7)(3.2,2.6)(3.2,1.7)(3.35,1.45)(3.2,1.2)(3.2,.3)(3.05,.2)}; \node[NoteRed,right] at(3.05,1.7){$\frac{t_2}2$};
}
\end{handpage}

% SOURCE PAGE 063
\begin{handpage}{63}
\handline{72}{131}{830}{$\therefore$当$|x|\to\infty$时，$m(Q(t))=t^n\sim|x|^n$.}
\handline{72}{162}{805}{$\therefore\chi_E^*(x)\sim\frac{m(E)}{|x|^n}$.}
\handline{72}{188}{805}{即$\exists c_1,c_2>0$，s.t. $\frac{c_1m(E)}{|x|^n}\le\chi_E^*(x)\le\frac{c_2m(E)}{|x|^n}$\quad ($|x|\to\infty$).}
\handline{72}{215}{805}{而$\frac1{|x|^n}\notin L^1(\R^n)$.}
\handline{72}{242}{800}{$\therefore\chi_E^*(x)\notin L^1(\R^n)$.}
\handline[Teal]{72}{270}{820}{事实上，若$f$可测，$f$在$\R^n$的一个正测度集上不为$0$，则$f^*\notin L^1(\R^n)$}
\handline{72}{300}{856}{\textbf{Def.}\quad 若$f$可测，$\exists c>0$，s.t. $\forall\alpha>0$，}
\handline{72}{328}{825}{$m(\{x\mid f(x)>\alpha\})\le\frac c\alpha$，则称$f$是弱$L^1$的\ \ann{(weak $L^1$)}，记为$f\in WL^1(\R^n)$.}
\handline{72}{353}{820}{这种估计称为弱估计\ \ann{(weak type estimate)}}
\handline[Teal]{72}{383}{850}{\circled{1}\ 若$f\in L^1$，则由Tchebyshev不等式}
\handline[Teal]{72}{412}{820}{$m(\{x\mid f(x)>\alpha\})\le\frac1\alpha\int_{\R^n}|f|=\frac{\|f\|_1}\alpha$，$\therefore f\in WL^1$}
\handline[Teal]{72}{441}{820}{于是$L^1\subseteq WL^1$}
\handline[Teal]{72}{471}{850}{\circled{2}\ $f(x)=\frac1{|x|^n}:\R^n\to\R$，则$f\in WL^1$，但$f\notin L^1$.}
\handline[Teal]{72}{503}{810}{以$n=1$为例，\quad$f(x)=\frac1{|x|}\notin L^1$.}
\handline[Teal]{72}{530}{815}{$\forall\alpha>0$，$m(\{x\mid f(x)>\alpha\})=m(\{x\mid|x|<\frac1\alpha\})=\frac2\alpha$，$\therefore f(x)\in WL^1$.}
\handline{72}{565}{856}{\textbf{Thm.}\ \ann{(Hardy--Littlewood)}}
\handline{72}{595}{856}{若$f\in L^1(\R^n)$，则$f^*\in WL^1(\R^n)$，$\exists C>0$，s.t. $m(\{x\mid f^*(x)>\alpha\})\le\frac C\alpha\int_{\R^n}|f|=\frac C\alpha\|f\|_1$（$\forall\alpha>0$）}
\handline{72}{625}{856}{\textbf{Pf.}\quad\circled{1}\ 若$f$是紧支撑的，则当$|x|$充分大时，$\exists r>0$，s.t. $|x|>r$时，$\exists c>0$，s.t. $f^*(x)\le\frac c{|x|^n}$.}
\handline{72}{656}{835}{$\forall\alpha>0$，令$E=\{x\mid f^*(x)>\alpha\}$，当$|x|>r$时，若$x\in E$则$\alpha<f^*(x)\le\frac c{|x|^n}\Rightarrow|x|^n\le\frac c\alpha$.}
\handline{72}{686}{830}{$\therefore E$是有界的，$\therefore m(E)<\infty$.}
\handline{72}{717}{830}{$\forall x\in E$，$f^*(x)>\alpha$，$\therefore\exists Q_x$，s.t. $\frac1{m(Q_x)}\int_{Q_x}|f|>\alpha$.}
\handline{72}{748}{830}{令$K=\bigcup_{x\in E}Q_x$，则$E\subseteq\bigcup_{x\in E}Q_x$.}
\handline{72}{780}{830}{由Vitali覆盖引理，$\exists\beta>0$，$Q_1,\ldots,Q_N$两两不交，s.t. $\beta m(E)\le\sum_{j=1}^Nm(Q_j)$.}
\handline{72}{812}{830}{$\therefore m(E)\le\frac1\beta\sum_{j=1}^Nm(Q_j)\le\frac1{\alpha\beta}\sum_{j=1}^N\int_{Q_j}|f|=\frac1{\alpha\beta}\int_{\bigcup_{j=1}^NQ_j}|f|\le\frac1{\alpha\beta}\|f\|_1$.}
\handline{72}{843}{850}{\circled{2}\ 一般情形，$f\in L^1(\R^n)$，不妨设$f\ge0$. 令$f_k=f\cdot\chi_{\overline{B_k}}$，则$f_k$是紧支撑的且$f_k\uparrow f$.}
\handline{72}{874}{835}{由\circled{1}$\Rightarrow m(\{x\mid f_k^*(x)>\alpha\})\le\frac C\alpha\int_{\R^n}|f_k(x)|\,dx\le\frac C\alpha\int_{\R^n}|f(x)|\,dx=\frac C\alpha\|f\|_1$}
\handline{72}{907}{856}{由$f_k\uparrow f\Rightarrow f_k^*\uparrow f\Rightarrow\{x\mid f_k^*(x)>\alpha\}\uparrow\{x\mid f(x)>\alpha\}\Rightarrow m(\{x\mid f(x)>\alpha\})\le\frac C\alpha\|f\|_1$.}
\end{handpage}

% SOURCE PAGE 064
\begin{handpage}{64}
\handline{72}{105}{750}{回到 Lebesgue 微分定理的证明：}
\handline{72}{133}{795}{$ f\in L^1(\R^n)$，则 $\lim_{Q\downarrow x}\frac1{m(Q)}\int_Q f(y)\,dy=f(x)$ a.e.}
\handline{72}{162}{755}{\textbf{Pf.}\quad 由于 $C_0(\R^n)$ 在 $L^1(\R^n)$ 中稠密.}
\handline{72}{190}{730}{选定 $\{C_k\}\subseteq C_0(\R^n)$，s.t. $\int_{\R^n}|f-C_k|\to0$.}
\handline{72}{218}{715}{令 $ F(Q)=\int_Q f(y)\,dy$，$ F_k(Q)=\int_Q C_k(y)\,dy$.}
\handline{72}{256}{760}{$\left|\frac{F(Q)}{m(Q)}-f(x)\right|\leq\left|\frac{F(Q)}{m(Q)}-\frac{F_k(Q)}{m(Q)}\right|+\left|\frac{F_k(Q)}{m(Q)}-C_k(x)\right|+|C_k(x)-f(x)|$.}
\handline{72}{299}{780}{$\therefore\limsup_{Q\downarrow x}\left|\frac{F(Q)}{m(Q)}-f(x)\right|\leq\limsup_{Q\downarrow x}\left|\frac{F(Q)}{m(Q)}-\frac{F_k(Q)}{m(Q)}\right|+\limsup_{Q\downarrow x}\left|\frac{F_k(Q)}{m(Q)}-C_k(x)\right|+|C_k(x)-f(x)|$.}
\handline{72}{330}{765}{由于 $ C_k\in C_0(\R^n)$，故 $\lim_{Q\downarrow x}\frac1{m(Q)}\int_Q C_k(y)\,dy=C_k(x)$.}
\handline{72}{361}{720}{$\therefore\limsup_{Q\downarrow x}\left|\frac{F_k(Q)}{m(Q)}-C_k(x)\right|=0$.}
\handline{72}{392}{785}{$\therefore\limsup_{Q\downarrow x}\left|\frac{F(Q)}{m(Q)}-f(x)\right|\leq\limsup_{Q\downarrow x}\left|\frac{F(Q)}{m(Q)}-\frac{F_k(Q)}{m(Q)}\right|+|C_k(x)-f(x)|$}
\handline{72}{429}{580}{$=\limsup_{Q\downarrow x}\left|\frac{\int_Q(f(y)-C_k(y))\,dy}{m(Q)}\right|+|C_k(x)-f(x)|$}
\handline{72}{458}{555}{$\leq(f-C_k)^*(x)+|C_k(x)-f(x)|$.}
\handline{72}{488}{810}{$\therefore\forall\varepsilon>0,\quad\{x\mid\mathrm{LHS}(x)>\varepsilon\}\subseteq\{x\mid(f-C_k)^*(x)>\dfrac\varepsilon2\}\cup\{x\mid|C_k(x)-f(x)|>\dfrac\varepsilon2\}$.}
\handline{72}{520}{790}{由 H--L 定理，$\exists C>0$，s.t. $ m(\{x\mid(f-C_k)^*(x)>\frac\varepsilon2\})\leq\frac{2C}\varepsilon\int_{\R^n}|f-C_k|$.}
\handline{72}{554}{790}{由 Tchebyshev $\Rightarrow m(\{x\mid|C_k(x)-f(x)|>\frac\varepsilon2\})\leq\frac2\varepsilon\int_{\R^n}|f-C_k|$.}
\handline{72}{588}{775}{$\therefore m(\{x\mid\mathrm{LHS}(x)>\varepsilon\})\leq\frac{2C+2}\varepsilon\int_{\R^n}|f-C_k|\to0\ (k\to\infty)$.}
\handline{72}{618}{720}{$\therefore m(\{x\mid\mathrm{LHS}(x)>\varepsilon\})=0$.}
\handline{72}{650}{770}{$\therefore m(\{x\mid\mathrm{LHS}(x)>0\})\leq\sum_{k=1}^{\infty}m(\{x\mid\mathrm{LHS}(x)>\frac1k\})=0$.}
\handline{72}{684}{760}{$\therefore\mathrm{LHS}(x)=0$ a.e. $\therefore\lim_{Q\downarrow x}\frac1{m(Q)}\int_Q f(y)\,dy=f(x)$ a.e.}
\handline[AnnotationBlue]{72}{717}{805}{Lebesgue 微分定理的证明用到了以下工具，是一个很深刻的定理.}
\handline[AnnotationBlue]{72}{750}{795}{\circled{1} $C_0(\R^n)$ 在 $L^1(\R^n)$ 中稠密 \circled{2} Hardy--Littlewood 极大函数}
\handline[AnnotationBlue]{72}{782}{795}{\circled{3} Hardy--Littlewood 定理 \circled{4} Vitali 覆盖引理 \circled{5} Tchebyshev 不等式.}
\handline[Teal]{72}{816}{795}{如果要证明 $L^1_{\mathrm{loc}}(\R^n)$ 的版本，需要做局部逼近.}
\handline[Teal]{72}{849}{795}{取紧集 $B_k=\overline B(0,k)$，令 $f_k=f\cdot\chi_{B_k}$，则 $f_k\in L^1(\R^n)$.}
\handline[Teal]{72}{880}{640}{再利用 LDCT 即可.}
\end{handpage}

% SOURCE PAGE 065
\begin{handpage}{65}
\handline{72}{131}{825}{下面是 Lebesgue 微分定理的应用.}
\handline{72}{158}{856}{\textbf{Def.}\quad$E\subseteq\R^n$ 可测，则 $\chi_E\in L^1_{\mathrm{loc}}(\R^n)$，故 $\lim_{Q\downarrow x}\frac1{m(Q)}\int_Q\chi_E=\lim_{Q\downarrow x}\frac{m(Q\cap E)}{m(Q)}=\chi_E(x)$ a.e.}
\handline{72}{186}{795}{若 $\lim_{Q\downarrow x}\frac{m(Q\cap E)}{m(Q)}=1$，则称 $x$ 是 $E$ 的密度点. \ann{（point of density）}}
\handline{72}{216}{795}{若 $\lim_{Q\downarrow x}\frac{m(Q\cap E)}{m(Q)}=0$，则称 $x$ 是 $E$ 的稀疏点. \ann{（point of dispersion）}}
\handline[Teal]{72}{244}{760}{\circled{1} 若 $x$ 是 $E$ 的密度点，则 $x$ 是 $E^c$ 的稀疏点.}
\handline[Teal]{72}{271}{735}{若 $x$ 是 $E$ 的稀疏点，则 $x$ 是 $E^c$ 的密度点.}
\handline[Teal]{72}{302}{790}{\circled{2} 如果 $x$ 是 $E$ 的密度点，那么表明 $x$ 周围聚集了很多 $E$ 中的点.}
\handline{72}{330}{856}{\textbf{Cor.}\quad$E\subseteq\R^n$ 可测，则几乎每个 $E$ 中的点都是 $E$ 的密度点. \ann{（Lebesgue 密度定理）}}
\handline{72}{362}{856}{\textbf{Def.}\quad$f$ 是 $\R^n$ 上可测函数. 若 $\lim_{Q\downarrow x}\frac1{m(Q)}\int_Q|f(y)-f(x)|\,dy=0$}
\handline{72}{389}{740}{则称 $x$ 是 $f$ 的 Lebesgue 点.}
\handline{72}{418}{790}{$f$ 的 Lebesgue 点构成的集合称为 $f$ 的 Lebesgue 集，记作 $L_f$.}
\handline[Teal]{72}{448}{805}{显然 Lebesgue 点满足 $\lim_{Q\downarrow x}\frac1{m(Q)}\int_Q f(y)\,dy=f(x)$，但反之不一定!}
\handline[Teal]{72}{479}{815}{下面加强 Lebesgue 定理，不仅“不可导”点构成零测集，非 Lebesgue 点也构成零测集.}
\handline{72}{512}{856}{\textbf{Thm.}\quad$f\in L^1_{\mathrm{loc}}(\R^n)$，则 $m(\R^n\setminus L_f)=0$.}
\handline{72}{540}{840}{\textbf{Pf.}\quad$r_k\in\Q$，则 $f-r_k\in L^1_{\mathrm{loc}}(\R^n)\Rightarrow|f-r_k|\in L^1_{\mathrm{loc}}(\R^n)$.}
\handline{72}{571}{785}{$\therefore\lim_{Q\downarrow x}\frac1{m(Q)}\int_Q|f-r_k|=|f(x)-r_k|$ a.e. $x\in\R^n$. \ann{$(*)$}}
\handline{72}{601}{800}{设 $Z_k$ 为使得 \ann{$(*)$} 不成立的点构成的集合，则 $m(Z_k)=0$.}
\handline{72}{630}{755}{令 $ Z=\bigcup_{k=1}^{\infty}Z_k$，则 $m(Z)=0$.}
\handline{72}{663}{815}{$\forall x\notin Z\Rightarrow\forall k\in\N,\ x\notin Z_k\Rightarrow\lim_{Q\downarrow x}\frac1{m(Q)}\int_Q|f-r_k|=|f(x)-r_k|$.}
\handline{72}{697}{795}{$\therefore\frac1{m(Q)}\int_Q|f(y)-f(x)|\,dy\leq\frac1{m(Q)}\int_Q|f(y)-r_k|\,dy+\frac1{m(Q)}\int_Q|f(x)-r_k|\,dy$.}
\handline{72}{729}{585}{$=\frac1{m(Q)}\int_Q|f(y)-r_k|\,dy+|f(x)-r_k|$.}
\handline{72}{764}{810}{$\Rightarrow\limsup_{Q\downarrow x}\frac1{m(Q)}\int_Q|f(y)-f(x)|\,dy\leq\lim_{Q\downarrow x}\frac1{m(Q)}\int_Q|f(y)-r_k|\,dy+|f(x)-r_k|=2|f(x)-r_k|$.}
\handline{72}{795}{810}{由 $\Q$ 在 $\R$ 中的稠密性，故可不妨令 $r_k\to f(x)\ (k\to\infty)$.}
\handline{72}{831}{820}{$\Rightarrow\lim_{Q\downarrow x}\frac1{m(Q)}\int_Q|f(y)-f(x)|\,dy$，$\forall x\notin Z$. 即 $x$ 是 $f$ 的 Lebesgue 点.}
\end{handpage}

% SOURCE PAGE 066
\begin{handpage}{66}
\handline{72}{125}{775}{下面从另一个角度来推广 Lebesgue 微分定理，把立方体换成更一般的集合.}
\handline{72}{152}{687}{\textbf{Def.}\quad$\{S_\alpha\}_{\alpha\in\Lambda}$，$S_\alpha\subseteq\R^n$ 可测，}
\handline{72}{181}{540}{若 \circled{1} $\operatorname{diam}(S_\alpha)\to0$}
\handline{72}{208}{540}{\circled{2} 若 $Q_\alpha$ 是以 $x$ 为中心，包含 $S_\alpha$ 的最小立方体.}
\handline{72}{235}{535}{$\exists$ 常数 $K$ s.t. $m(Q_\alpha)\leq K\,m(S_\alpha)$.}
\handline{72}{261}{660}{则称 $S_\alpha$ 正则收缩到 $x$ \ann{（shrink regularly）}}
\handdiagram{777}{159}{151}{72}{
\draw[AnnotationBlue](0,0)rectangle(2,1.7);\fill(.9,.9)circle(.025);\node[below]at(.9,.9){$x$};\draw(1.7,1.15)ellipse(.35 and .2);\node at(1.7,1.15){$S_\alpha$};\node[AnnotationBlue,right]at(2,0){$Q_\alpha$};
}
\handline{72}{292}{760}{\textbf{Thm.}\quad$f\in L^1_{\mathrm{loc}}(\R^n)$，$\{S_\alpha\}_{\alpha\in\Lambda}$ 正则收缩到 $x$.}
\handline{72}{320}{730}{若 $x\in L_f$，则 $\frac1{m(S_\alpha)}\int_{S_\alpha}f(y)\,dy\to f(x)$.}
\handline{72}{345}{130}{\textbf{Pf.}}
\handline{72}{378}{675}{$\frac1{m(S_\alpha)}\int_{S_\alpha}|f(y)-f(x)|\,dy=\frac{m(Q_\alpha)}{m(S_\alpha)}\cdot\frac1{m(Q_\alpha)}\int_{S_\alpha}|f(y)-f(x)|\,dy$.}
\handline{72}{422}{525}{$\leq K\cdot\frac1{m(Q_\alpha)}\int_{Q_\alpha}|f(y)-f(x)|\,dy$.}
\handline{72}{468}{685}{$\therefore Q_\alpha\to x$ 时，$\frac1{m(S_\alpha)}\int_{S_\alpha}|f(y)-f(x)|\,dy\to0$.}
\handline{72}{501}{630}{$\therefore\frac1{m(S_\alpha)}\int_{S_\alpha}f(y)\,dy\to f(x)$.}
\handline[Teal]{72}{532}{740}{由于 $m(\R^n\setminus L_f)=0$，于是 $\frac1{m(S_\alpha)}\int_{S_\alpha}f(y)\,dy\to f(x)$ a.e. $x\in\R^n$.}
\handline[Teal]{72}{563}{810}{特别地，$n=1$ 时，令 $Q_x=[x-h,x+h]$.}
\handline[Teal]{72}{597}{665}{$\frac1{m(Q_x)}\int_{Q_x}f(y)\,dy=\frac1{2h}\int_{x-h}^{x+h}f(y)\,dy$}
\handline[Teal]{72}{631}{735}{由 Lebesgue 微分定理，$\lim_{h\to0+}\frac1{2h}\int_{x-h}^{x+h}f(y)\,dy=f(x)$ a.e.}
\handline[Teal]{72}{659}{700}{再取 $S_x=[x,x+h]\subseteq Q_x$，则 $m(Q_x)=2m(S_x)$，}
\handline[Teal]{72}{692}{620}{且 $\operatorname{diam}(S_x)\to0\ (h\to0+)$.}
\handline[Teal]{72}{723}{660}{故 $S_x$ 正则收缩到 $x$.}
\handline[Teal]{72}{759}{705}{$\therefore\lim_{h\to0+}\frac1h\int_x^{x+h}f(y)\,dy=f(x)$ a.e.}
\end{handpage}

% SOURCE PAGE 067
% SECTION: 7.3 | Vitali Covering Lemma
\begin{handpage}{67}
\handtitle{72}{130}{680}{§6.3 Vitali Covering Lemma}
\handline{72}{157}{830}{\textbf{Def.}\quad$K$ 是一族立方体，若 $\forall x\in E,\ \eta>0$，$\exists Q\in K$ s.t. $x\in Q$ 且 $\operatorname{diam}(Q)<\eta$.}
\handline{72}{181}{720}{则称 $K$ 在 Vitali 意义下覆盖了 $E$.}
\handline{72}{209}{740}{\textbf{Thm.}\quad\ann{（Vitali Covering Lemma）.}}
\handline{72}{236}{800}{$E\subseteq\R^n$，$0<m^*(E)<\infty$，$K$ 是一族闭立方体，在 Vitali 意义下覆盖了 $E$.}
\handline{72}{263}{825}{则 $\forall\varepsilon>0$，$\exists\{Q_j\}_{j=1}^{\infty}\subseteq K$，两两不交，s.t. $ m(E\setminus(\bigcup_{j=1}^{\infty}Q_j))=0$，且 $\sum_{j=1}^{\infty}m(Q_j)\leq(1+\varepsilon)m^*(E)$.}
\handline{72}{292}{810}{\textbf{Pf.}\quad$\forall\varepsilon>0$，$\exists$ 开集 $G\subseteq\R^n$，s.t. $E\subseteq G$，$m(G)\leq(1+\varepsilon)m^*(E)$}
\handline{72}{319}{760}{由简单版本，$\exists0<\beta<1$，$Q_1,\ldots,Q_{N_1}$ 两两不交，}
\handline{72}{345}{760}{s.t. $\sum_{j=1}^{N_1}m(Q_j)\geq\beta\,m^*(E)$.}
\handline{72}{373}{780}{令 $ E_1=E\setminus(\bigcup_{j=1}^{N_1}Q_j)$，$\therefore m^*(E_1)\leq m(G\setminus(\bigcup_{j=1}^{N_1}Q_j))\leq m(G)-\sum_{j=1}^{N_1}m(Q_j)\leq(1+\varepsilon-\beta)m^*(E)$.}
\handline{72}{401}{735}{不妨设 $0<\varepsilon<\dfrac\beta2$，则 $m^*(E_1)\leq(1-\dfrac\beta2)m^*(E)$.}
\handline{72}{427}{760}{设 $ F_1=\bigcup_{j=1}^{N_1}Q_j$，$K_1=\{Q\in K\mid Q\cap F_1=\varnothing\}$.}
\handline{72}{453}{730}{下证 $K_1$ 在 Vitali 意义下覆盖了 $E_1$.}
\handline{72}{483}{765}{$\forall x\in E_1=E\setminus F_1$，由于 $F_1$ 是闭集，$x\notin F_1$，$\therefore d(x,F_1)>0$.}
\handline{72}{512}{760}{由于 $K$ 在 Vitali 意义下覆盖了 $E$，$\therefore\exists Q\in K$ s.t. $x\in Q$ 且 $\operatorname{diam}(Q)<d(x,F_1)$.}
\handline{72}{541}{755}{由于 $\operatorname{diam}(Q)<d(x,F_1)$ 且 $x\in Q\Rightarrow Q\cap F_1=\varnothing\quad\therefore Q\in K_1$.}
\handline{72}{570}{745}{$\therefore K_1$ 在 Vitali 意义下覆盖了 $E_1$.}
\handline{72}{601}{760}{故 $\exists Q_{N_1+1},\ldots,Q_{N_2}$，两两不交，s.t. $ E_2=E_1\setminus(\bigcup_{j=N_1+1}^{N_2}Q_j)=E\setminus\bigcup_{j=1}^{N_2}Q_j$.}
\handline{72}{630}{720}{$ m^*(E_2)\leq(1-\frac\beta2)m^*(E_1)\leq(1-\frac\beta2)^2m^*(E)$.}
\handline{72}{658}{720}{$\ldots\ldots$ 一直进行下去，得到 $\{Q_j\}_{j=1}^{\infty}$.}
\handline{72}{686}{730}{而 $ m^*(E_n)\leq(1-\frac\beta2)^n m^*(E)\to0\ (n\to\infty)$}
\handline{72}{718}{780}{$\Rightarrow m^*(E\setminus(\bigcup_{j=1}^{\infty}Q_j))=0$，$\therefore m(E\setminus(\bigcup_{j=1}^{\infty}Q_j))=0$.}
\handline{72}{747}{770}{最后考虑 $K'=\{Q\in K\mid Q\subseteq G\}$.}
\handline{72}{777}{780}{$\forall x\in E$，由于 $G^c$ 是闭集，$\therefore d(x,G^c)>0$，$\therefore\exists Q\in K$，s.t. $x\in Q$ 且 $\operatorname{diam}(Q)<d(x,G^c)$.}
\handline{72}{806}{760}{由于 $\operatorname{diam}(Q)<d(x,G^c)$，$x\in Q$ 故 $Q\subseteq G\quad\therefore Q\in K'$.}
\handline{72}{838}{775}{故 $K'$ 在 Vitali 意义下覆盖了 $E$，故不妨设 $\{Q_j\}_{j=1}^{\infty}$ 全部包含于 $G$.}
\handline{72}{867}{765}{$\therefore\sum_{j=1}^{\infty}m(Q_j)\leq m(G)\leq(1+\varepsilon)m^*(E)$}
\handline[AnnotationBlue]{72}{893}{800}{\textbf{Cor.}\quad$ m(E\setminus\bigcup_{j=1}^{\infty}Q_j)=0.\ \therefore\forall\varepsilon>0,\ \exists N$, s.t. $n\geq N\quad m^*(E\setminus\bigcup_{j=1}^N Q_j)<\varepsilon$.}
\handline[AnnotationBlue]{72}{919}{800}{由 Carathéodory $\Rightarrow m^*(E)=m^*(E\cap\bigcup_{j=1}^N Q_j)+m^*(E\setminus\bigcup_{j=1}^N Q_j)\Rightarrow m^*(E\cap\bigcup_{j=1}^N Q_j)>m^*(E)-\varepsilon$}
\end{handpage}

% SOURCE PAGE 068
% SECTION: 7.4 | Differentiation of Monotone Functions
\begin{handpage}{68}
\handtitle{72}{129}{735}{§6.4 Differentiation of Monotone Functions}
\handline{72}{157}{755}{\textbf{Def.}\quad$f:\R\to\R$，$x\in\R$}
\handline{72}{185}{265}{$ D_1f(x)=\limsup_{h\to0+}\frac{f(x+h)-f(x)}h$}
\handline{72}{211}{265}{$ D_2f(x)=\liminf_{h\to0+}\frac{f(x+h)-f(x)}h$}
\handline{355}{213}{430}{称为 Dini 数 \ann{（Dini numbers 或 Dini derivatives）}}
\handline{72}{238}{265}{$ D_3f(x)=\limsup_{h\to0-}\frac{f(x+h)-f(x)}h$}
\handline[Teal]{354}{240}{430}{显然 $D_1f(x)\geq D_2f(x)$，$D_3f(x)\geq D_4f(x)$.}
\handline{72}{263}{270}{$ D_4f(x)=\liminf_{h\to0-}\frac{f(x+h)-f(x)}h$}
\handline{72}{292}{780}{\textbf{Thm.}\quad$f:(a,b)\to\R$ 单增，则在 $(a,b)$ 上 $f'(x)$ 几乎处处存在}
\handline{72}{320}{710}{且 $f'$ 是 $(a,b)$ 上的非负可测函数}
\handline{72}{348}{710}{且 $0\leq\int_a^b f'\leq f(b-)-f(a+)$.}
\handline{72}{374}{795}{\textbf{Pf.}\quad 令 $ A=\underset{r,s\in\Q}{A_{r,s}}=\{x\in(a,b)\mid D_1f(x)>r>s>D_4f(x)\}$，则 $m^*(A)\leq b-a<+\infty$.}
\handline{72}{428}{700}{若 $m^*(A)>0$，$x\in A$，$D_4f(x)<s$.}
\handline{72}{458}{720}{$\therefore$ 存在充分小的 $h_x>0$，s.t. $\frac{f(x-h_x)-f(x)}{-h_x}<s$.}
\handline{72}{487}{730}{令 $K=\{[x-h_x,x]\mid x\in A\}$，则 $K$ 在 Vitali 意义下覆盖了 $A$}
\handline{72}{516}{725}{由 Vitali 覆盖引理，$\forall\varepsilon>0$，}
\handline{72}{546}{725}{$\exists[x_j-h_j,x_j]$，$1\leq j\leq N$，}
\handline{72}{576}{730}{s.t. \circled{1} $f(x_j)-f(x_j-h_j)<sh_j$}
\handline{72}{608}{660}{\circled{2} $ m^*(A\cap(\bigcup_{j=1}^N[x_j-h_j,x_j]))>m^*(A)-\varepsilon$ \ann{（由前一页的 Cor 得到）.}}
\handline{72}{639}{660}{\circled{3} $\sum_{j=1}^N h_j\leq(1+\varepsilon)m^*(A)$.}
\handline{72}{691}{710}{$\therefore\sum_{j=1}^N(f(x_j)-f(x_j-h_j))\leq s\sum_{j=1}^N h_j\leq s(1+\varepsilon)m^*(A)$.}
\handline{72}{723}{730}{令 $ B=A\cap(\bigcup_{j=1}^N[x_j-h_j,h_j])$，则 $0<m^*(B)\leq m^*(A)<\infty$.}
\handline{72}{752}{730}{若 $y\in B$ 且不为 $[x_j-h_j,h_j]\quad1\leq j\leq N$ 中的端点，}
\handline{72}{785}{730}{则 $D_1(y)>r\quad\therefore\exists$ 充分小的 $k_y>0$，s.t. $\frac{f(y+k_y)-f(y)}{k_y}>r$.}
\handline{72}{817}{745}{令 $K'=\{[y,y+k_y]\mid y\in B\}$，则 $K'$ 在 Vitali 意义下覆盖了 $B$.}
\handline{72}{847}{730}{同理\quad$\exists[y_i,y_i+k_i]\quad1\leq i\leq M$.}
\handline{72}{881}{715}{s.t. \circled{1} $[y_i,y_i+k_i]\subseteq[x_j-h_j,x_j]$ for some $j$. \ann{（需要解释）.}}
\end{handpage}

% SOURCE PAGE 069
\begin{handpage}{69}
\handline{72}{126}{740}{\circled{2} $f(y_i+k_i)-f(y_i)>rk_i$}
\handline{72}{154}{735}{\circled{3} $\sum_{i=1}^M k_i>m^*(B)-\varepsilon>m^*(A)-2\varepsilon$}
\handline{72}{182}{740}{\circled{4} $\sum_{i=1}^M[f(y_i+k_i)-f(y_i)]\geq r\sum_{i=1}^M k_i\geq r(m^*(A)-2\varepsilon)$.}
\handline{72}{210}{785}{由于 $f$ 单调增.}
\handline{72}{240}{750}{$\therefore\sum_{i=1}^M[f(y_i+k_i)-f(y_i)]\leq\sum_{j=1}^N[f(x_j)-f(x_j-h_j)]$.}
\handline{72}{267}{735}{$\therefore r(m^*(A)-2\varepsilon)\leq s(1+\varepsilon)m^*(A)$.}
\handline{72}{295}{715}{令 $\varepsilon\to0+\Rightarrow r\,m^*(A)\leq s\,m^*(A)$.}
\handline{72}{323}{735}{当 $m^*(A)>0$ 时 $\Rightarrow r\leq s$，与 $r>s$ 矛盾!}
\handline{72}{349}{735}{$\therefore m^*(A)=0$.}
\handline{72}{378}{765}{令 $\widetilde A=\{x\in(a,b)\mid D_1f(x)>D_4f(x)\}$}
\handline{72}{406}{745}{则 $\widetilde A\subseteq\bigcup_{r,s\in\Q}A_{r,s}\qquad(r,s)\in\Q^2$ 为可数集}
\handline{72}{435}{725}{$\therefore m^*(\widetilde A)\leq\sum_{r,s\in\Q}m^*(A_{r,s})=0\Rightarrow m(\widetilde A)=0$.}
\handline{72}{464}{745}{$\therefore D_1f(x)=D_4f(x)$ a.e.}
\handline{72}{490}{765}{同理可证 $D_1f(x)=D_2f(x)=D_3f(x)=D_4f(x)$ a.e.}
\handline{72}{520}{760}{$\therefore f'(x)$ 在 $(a,b)$ 上几乎处处存在.}
\handline{72}{547}{765}{将 $f$ 延拓到 $(a,+\infty)$，令 $f(x)\equiv f(b-)\ (x\geq b)$.}
\handline{72}{581}{745}{令 $ f_k=\frac{f(x+\frac1k)-f(x)}{\frac1k}$，则 $f_k\to f'$ a.e.}
\handline{72}{640}{745}{由于 $f_k\geq0\Rightarrow f'\geq0$ a.e.，且 $f'$ 可测.}
\handline{72}{675}{755}{由 Fatou 引理}
\handline{72}{711}{660}{$\int_a^b\underline{\lim}\,f_k\leq\underline{\lim}\int_a^b f_k\quad$ 即 $\int_a^b f'\leq\underline{\lim}\int_a^b f_k$.}
\handline{72}{744}{750}{而 $\int_a^b f_k=k\int_a^b[f(x+\frac1k)-f(x)]\,dx$}
\handline{72}{782}{665}{$=k\left(\int_{a+\frac1k}^{b+\frac1k}f(x)\,dx-\int_a^b f(x)\,dx\right)=\frac{\int_b^{b+\frac1k}f(x)\,dx}{\frac1k}-\frac{\int_a^{a+\frac1k}f(x)\,dx}{\frac1k}$}
\handline{72}{833}{655}{$=f(b-)-\frac{\int_a^{a+\frac1k}f(x)\,dx}{\frac1k}$}
\handline{72}{869}{790}{而 $ f(a+)\cdot\frac1k\leq\int_a^{a+\frac1k}f(x)\,dx\leq f(a+\frac1k)\frac1k$}
\handline{72}{906}{780}{$\Rightarrow f(a+)\leq\frac{\int_a^{a+\frac1k}f(x)\,dx}{\frac1k}\leq f(a+\frac1k)$，令 $k\to\infty$ 则 $f(a+\frac1k)\to f(a+)$.}
\handline{72}{951}{850}{$\therefore\frac{\int_a^{a+\frac1k}f(x)\,dx}{\frac1k}\to f(a+)\ (k\to\infty)\quad\therefore\lim\int_a^b f_k=f(b-)-f(a+)\quad\therefore0\leq\int_a^b f'\leq f(b-)-f(a+)$}
\end{handpage}

% SOURCE PAGE 070
\begin{handpage}{70}
\handline{72}{117}{750}{$\int_a^b f'\leq f(b-)-f(a+)$ 的等号不一定能取到.}
\handline[Teal]{72}{142}{755}{例如 Cantor 函数 $C(x)$，在 $[0,1]$ 上单增且连续.}
\handline[Teal]{72}{171}{750}{$ C'(x)=0$ a.e. $\therefore\int_0^1 C'(x)\,dx=0.<C(1)-C(0)=1-0=1$.}
\handline{72}{198}{810}{\textbf{Cor.}\quad 若 $f\in\mathrm{BV}[a,b]$，则 $f'$ 在 $[a,b]$ 上几乎处处存在，且 $f'\in L[a,b]$.}
\handline[AnnotationBlue]{72}{227}{760}{（BV 函数可写为两个增函数之差）}
\handline[Teal]{72}{257}{745}{由于 BV 函数的变差函数 $V(x)=V[a,x]\geq0$ 且在 $[a,b]$ 上单增.}
\handline[Teal]{72}{285}{750}{故 $V'(x)$ 也几乎处处存在，自然会问 $V'(x)$ 与 $f'(x)$ 有什么关系？}
\handline{72}{316}{805}{\textbf{Lemma.}\quad\ann{（Fubini）}}
\handline{72}{347}{790}{$ f_k$ 在 $[a,b]$ 上单增，若 $\sum_{k=1}^{\infty}f_k(x)\to S(x)$ 逐点.}
\handline{72}{375}{785}{则 $ S'(x)=\sum_{k=1}^{\infty}f'_k(x)$ a.e.}
\handline{72}{403}{825}{\textbf{Pf.}\quad 令 $S(x)=S_m(x)+r_m(x)$.}
\handline{72}{433}{745}{$ S_m(x)=\sum_{k=1}^m f_k(x)$，$ r_m(x)=\sum_{k=m+1}^{\infty}f_k(x)$.}
\handline{72}{462}{740}{$S(x)$，$S_m(x)$，$r_m(x)$ 均单增，故导数几乎处处存在.}
\handline{72}{491}{750}{$\therefore\exists Z_m$，$m(Z_m)=0$，s.t. $x\notin Z_m$ 时：}
\handline{72}{520}{730}{$ S'(x)=S'_m(x)+r'_m(x)=\sum_{k=1}^m f'_k(x)+r'_m(x)$}
\handline{72}{552}{740}{令 $ Z=\bigcup_{m=1}^{\infty}Z_m$，则 $m(Z)=0$.}
\handline{72}{582}{745}{$\forall x\notin Z$，由于 $r'_m(x)\geq0$，$\therefore\sum_{k=1}^m f'_k(x)\leq S'(x)\Rightarrow\sum_{k=1}^{\infty}f'_k(x)\leq S'(x)$.}
\handline{72}{610}{740}{下面证明 $r'_m(x)\to0\ (m\to\infty)$ a.e.}
\handline{72}{642}{735}{$r_m(x)=f_m(x)+r_{m+1}(x)$ 且 $r_m(x),f_m(x),r_{m+1}(x)$ 均单增}
\handline{72}{672}{740}{$\therefore r'_m(x)=f'_m(x)+r'_{m+1}(x)$ a.e. 且 $f'_m(x)\geq0$ a.e.}
\handline{72}{700}{740}{$\Rightarrow r'_m(x)\geq r'_{m+1}(x)$ a.e.}
\handline{72}{732}{735}{$ r_m(a)=\sum_{k=m+1}^{\infty}f_k(a)$，$ r_m(b)=\sum_{k=m+1}^{\infty}f_k(b)\to0\ (m\to\infty)$.}
\handline{72}{764}{740}{$\therefore\exists$ 子列 $\{r_{m_j}\}_{j=1}^{\infty}$ s.t. $\sum_j r_{m_j}(a),\ \sum_j r_{m_j}(b)<\infty$.}
\handline{72}{793}{720}{考虑 $\int_a^b\sum_{j=1}^{\infty}r'_{m_j}(x)$，$r'_{m_j}(x)\geq0$ a.e.}
\handline{72}{825}{750}{$\therefore$ 由 Levi 引理 $\Rightarrow0\leq\int_a^b\sum_{j=1}^{\infty}r'_{m_j}(x)=\sum_{j=1}^{\infty}\int_a^b r'_{m_j}(x)\leq\sum_{j=1}^{\infty}r_{m_j}(b)-r_{m_j}(a)<\infty$.}
\handline{72}{855}{745}{$\Rightarrow\sum_{j=1}^{\infty}r'_{m_j}(x)\in L^1[a,b]\quad\therefore$ 在 $[a,b]$ 上几乎处处有限.}
\handline{72}{888}{745}{$\Rightarrow r'_{m_j}(x)\to0$ a.e.，而 $r'_m(x)\downarrow\quad\Rightarrow r'_m(x)\to0$ a.e. $\ (m\to\infty)$.}
\handline{72}{918}{415}{从而 $ S'(x)=\sum_{k=1}^{\infty}f'_k(x)$ a.e.}
\end{handpage}

% SOURCE PAGE 071
\begin{handpage}{71}
\handline{72}{137}{840}{Thm. $f\in\mathrm{BV}[a,b]$，则 $V'(x)=|f'(x)|$ a.e.}
\handline{72}{162}{840}{Pf. $f\in\mathrm{BV}[a,b]$.}
\handline{72}{189}{850}{$\forall k\in\N$，$\varepsilon_k=\frac1{2^k}>0$，$\exists$ 分划 $\Gamma_k=\{a=x_0^k<\cdots<x_{m_k}^k=b\}$}
\handline{72}{214}{630}{$S_{\Gamma_k}=\sum_{j=1}^{m_k}|f(x_j^k)-f(x_{j-1}^k)|$}
\handline{72}{239}{635}{s.t. $0\le V(b)-S_{\Gamma_k}<\frac1{2^k}$.}
\handline{72}{277}{856}{在 $[x_{j-1}^k,x_j^k]$ 上令 $f_k=\begin{cases}f(x)+C_j^k&\text{若 }f(x_j^k)\ge f(x_{j-1}^k),\\-f(x)+C_j^k&\text{若 }f(x_j^k)<f(x_{j-1}^k).\end{cases}$}
\handline{72}{317}{720}{并且选择合适的 $C_j^k$，使得 $f_k(a)=0$.}
\handline{72}{348}{720}{则 $S_{\Gamma_k}=\sum_{j=1}^{m_k}|f(x_{j-1}^k)-f(x_j^k)|=\sum_{j=1}^{m_k}f_k(x_j^k)-f_k(x_{j-1}^k)$}
\handline{72}{377}{420}{$=f_k(b)-f_k(a)=f_k(b)$.}
\handline{72}{407}{720}{$\therefore0\le V(b)-f_k(b)<\frac1{2^k}$.}
\handline{72}{436}{760}{断言：$V(x)-f_k(x)$ 在 $[a,b]$ 上单增.}
\handline{72}{472}{765}{\circled{1} $\forall s,t\in[x_{j-1}^k,x_j^k]$，$s<t$.}
\handline{72}{501}{750}{$f_k(t)-f_k(s)\le|f_k(t)-f_k(s)|\le V(t)-V(s)$.}
\handline{72}{528}{745}{$\therefore V(t)-f_k(t)\ge V(s)-f_k(s)$.}
\handline{72}{563}{775}{\circled{2} $s\in[x_{j_1}^k,x_{j_1+1}^k]$，$t\in[x_{j_2}^k,x_{j_2+1}^k]$，$j_1\le j_2-1$}
\handline{72}{592}{745}{则 $f_k(t)-f_k(s)=f_k(t)-f_k(x_{j_2}^k)+f_k(x_{j_2}^k)-f_k(x_{j_2-1}^k)\cdots-f_k(s)$}
\handline{72}{622}{705}{$\le|f_k(t)-f_k(x_{j_2}^k)|+\cdots+|f_k(x_{j_1+1}^k)-f_k(s)|\le V(t)-V(s)$}
\handline{72}{650}{745}{$\therefore V(t)-f_k(t)\ge V(s)-f_k(s)$.}
\handline{72}{685}{790}{而 $0\le\sum_{k=1}^\infty V(x)-f_k(x)\le\sum_{k=1}^\infty\frac1{2^k}\to1$，$\therefore$ 由 Weierstrass 定理 $\sum_{k=1}^\infty V(x)-f_k(x)$ 收敛.}
\handline{72}{718}{785}{由 Fubini 引理 $\bigl(\sum_{k=1}^\infty V(x)-f_k(x)\bigr)'=\sum_{k=1}^\infty V'(x)-f_k'(x)$ a.e.}
\handline{72}{751}{770}{$\therefore\sum_{k=1}^\infty V'(x)-f_k'(x)$ 收敛，$\therefore V'(x)-f_k'(x)\to0$ a.e.}
\handline{72}{782}{750}{$\therefore\lim_{k\to\infty}f_k'(x)=V'(x)$ a.e.}
\handline{72}{815}{755}{由 $f_k$ 的定义 $\Rightarrow|f_k'(x)|=|f'(x)|$ a.e.}
\handline{72}{849}{760}{$\therefore V'(x)=|V'(x)|=\lim_{k\to\infty}|f_k'(x)|=|f'(x)|$ a.e.}
\end{handpage}

% SOURCE PAGE 072
% SECTION: 7.5 | Absolutely continuous and singular functions
\begin{handpage}{72}
\handtitle{72}{96}{820}{\S6.5 Absolutely continuous and singular functions}
\handline{72}{126}{825}{Def. $f:[a,b]\to\R$}
\handline{72}{153}{765}{若 $\forall\varepsilon>0$，$\exists\delta>0$，s.t. $\{[a_i,b_i]\}_{i=1}^\infty$ 两两不交，$\subseteq[a,b]$.}
\handline{72}{181}{750}{当 $\sum_{i=1}^\infty(b_i-a_i)<\delta$ 时，$\sum_{i=1}^\infty|f(b_i)-f(a_i)|<\varepsilon$.}
\handline{72}{210}{765}{则称 $f$ 在 $[a,b]$ 上绝对连续 \ann{（A.C. absolutely continuous）}}
\handline{72}{239}{830}{Ex. \circled{1} $g\in L^1[a,b]$，$G(x)=\int_a^x g(t)\,dt$.}
\handline{72}{270}{760}{$\sum_{i=1}^\infty|G(b_i)-G(a_i)|=\sum_{i=1}^\infty\bigl|\int_{a_i}^{b_i}g(t)\,dt\bigr|\le\sum_{i=1}^\infty\int_{a_i}^{b_i}|g(t)|\,dt=\int_{\bigcup_{i=1}^\infty[a_i,b_i]}|g(t)|\,dt$.}
\handline{72}{295}{770}{由于 $h(E)=\int_E|g(t)|\,dt$ 作为集函数是绝对连续的.}
\handline{72}{324}{770}{$\therefore\sum_{i=1}^\infty(b_i-a_i)\to0$ 时，$h(\bigcup_{i=1}^\infty[a_i,b_i])=\int_{\bigcup_{i=1}^\infty[a_i,b_i]}|g(t)|\,dt\to0$.}
\handline{72}{352}{765}{$\therefore G(x)\in\mathrm{AC}[a,b]$.}
\handline{72}{383}{790}{\circled{2} Lipschitz 函数，$|f(x)-f(y)|\le L|x-y|$，$L>0$.}
\handline{72}{414}{765}{$\sum_{i=1}^\infty|f(b_i)-f(a_i)|\le\sum_{i=1}^\infty L|b_i-a_i|<L\delta$，取 $\delta=\frac\varepsilon L$ 即可.}
\handline{72}{442}{760}{故 Lipschitz 函数 $\in\mathrm{AC}[a,b]$.}
\handline[Teal]{72}{472}{770}{特别地 $g\in C^1[a,b]$ 则 $g\in\mathrm{AC}[a,b]$.}
\handline[Teal]{72}{503}{770}{$|g(x)-g(y)|=|g'(\xi)||x-y|\le(\max_{[a,b]}|g'(t)|)\cdot|x-y|$}
\handline[Teal]{72}{533}{770}{$\therefore g$ 是 $[a,b]$ 上 Lipschitz 函数 $\Rightarrow g\in\mathrm{AC}[a,b]$.}
\handline{72}{566}{785}{\circled{3} Cantor 函数 $C(x)\notin\mathrm{AC}[0,1]$.}
\handline{72}{597}{760}{令 $C_k=\bigcup_j[a_j^k,b_j^k]$，为第 $k$ 次操作后剩下的区间之并.}
\handline{72}{626}{750}{$a_j^k,b_j^k$ 为第 $k$ 次操作后剩下区间的端点.}
\handline{72}{659}{765}{则 $m(C_k)\to0$（$k\to\infty$）. 但 $\sum_j C(b_j^k)-C(a_j^k)\equiv1$.}
\handline{72}{691}{760}{$\therefore$ Cantor 函数不是绝对连续函数.}
\handline[AnnotationBlue]{72}{723}{760}{但是 Cantor 函数 $C(x)\in C[0,1]$ 且单增，$\therefore C'(x)=0$ a.e.}
\handline{72}{754}{856}{Def. $f:[a,b]\to\R$，若 $f'\equiv0$ a.e. 则称 $f$ 为 $[a,b]$ 上的奇异函数. \ann{(Singular)}}
\handline[Teal]{72}{791}{800}{特别地若 $f\in\mathrm{AC}[a,b]$，$\forall[x,y]\subseteq[a,b]$.}
\handline[Teal]{72}{824}{780}{$\forall\varepsilon>0$，$\exists\delta>0$，$|x-y|<\delta$ 则 $|f(x)-f(y)|<\varepsilon$.}
\handline[Teal]{72}{857}{800}{$\therefore\mathrm{AC}\Rightarrow$ 一致连续 $\Rightarrow$ 连续.}
\end{handpage}

% SOURCE PAGE 073
\begin{handpage}{73}
\handline{72}{150}{856}{Thm. 若 $f\in\mathrm{AC}[a,b]$ 则 $f\in\mathrm{BV}[a,b]$.}
\handline{72}{177}{856}{Pf. $\varepsilon=1$，$\exists\delta>0$，s.t. $\sum(b_i-a_i)<\delta$ 时，$\sum|f(b_i)-f(a_i)|<1$.}
\handline{72}{206}{800}{设 $N\in\N$ 且 $0<\frac{b-a}N<\delta$，将 $[a,b]$ $N$ 等分. $[a,b]=\bigcup_{i=0}^{N-1}[P_i,P_{i+1}]$，其中 $P_i=a+\frac{(b-a)i}N$.}
\handline{72}{233}{795}{对于 $0\le k\le N-1$，$[P_k,P_{k+1}]$，考虑 $V[P_k,P_{k+1}]$.}
\handline{72}{261}{790}{对于 $[P_k,P_{k+1}]$ 的任意分划，$\{P_{k,j}\}_{j=0}^n$.}
\handline{72}{287}{740}{$\sum_{j=0}^{n-1}(P_{k,j+1}-P_{k,j})=P_{k+1}-P_k=\frac{b-a}N<\delta$.}
\handline{72}{317}{775}{$\therefore\sum_{j=0}^{n-1}|f(P_{k,j+1})-f(P_{k,j})|<1\Rightarrow V[P_k,P_{k+1}]\le1$.}
\handline{72}{346}{800}{$\therefore V[a,b]\le\sum_{i=0}^{N-1}V[P_i,P_{i+1}]\le N<+\infty$，$\therefore f\in\mathrm{BV}[a,b]$.}
\handline{72}{376}{856}{Thm. 若 $f$ 在 $[a,b]$ 上绝对连续且奇异，则 $f$ 在 $[a,b]$ 上为常数.}
\handline{72}{405}{856}{Pf. 只需证明 $f(a)=f(b)$. 再考虑区间 $[a,x]$、$[x,b]$，则同理 $\Rightarrow f(x)=f(a)=f(b)$.}
\handline{72}{435}{815}{由于 $f$ 在 $(a,b)$ 上奇异，$\therefore f'(x)=0$ a.e. 令 $E=\{x\in(a,b)\mid f'(x)=0\}$ 则 $m(E)=b-a$.}
\handline{72}{466}{810}{$\forall\varepsilon>0$，$x\in E$，$\exists$ 充分小的 $h_x>0$，s.t. $|f(x+h_x)-f(x)|<\varepsilon h_x$ 且 $[x,x+h_x]\subset(a,b)$.}
\handline{72}{495}{805}{考虑 $K=\bigcup_{x\in E}[x,x+h_x]$，则 $K$ 在 Vitali 意义下覆盖了 $E$.}
\handline{72}{526}{805}{$f$ 绝对连续，$\therefore\exists\delta>0$，s.t. $\sum(b_i-a_i)<\delta$ 则 $\sum|f(b_i)-f(a_i)|<\varepsilon$.}
\handline{72}{555}{805}{由 Vitali 覆盖引理，$\exists[x_1,x_1+h_1],\ldots,[x_n,x_n+h_n]$ 两两不交}
\handline{72}{585}{790}{s.t. $\sum_{j=1}^n h_j>m(E)-\delta=(b-a)-\delta$.}
\handdiagram{123}{594}{532}{42}{
\draw[NoteRed](0,0)--(10,0);\foreach\x/\t in {0/a,1/x_1,2/{x_1+h_1},3/x_2,4/{x_2+h_2},7/x_n,9/{x_n+h_n},10/b}{\draw[NoteRed](\x,0)--++(0,.1);\node[NoteRed,below,font=\tiny]at(\x,0){$\t$};}
\foreach\a/\b in {0/1,2/3,4/4.5,9/10}{\draw[AnnotationBlue,thick](\a,.05)--(\b,.05);}\foreach\x/\t in{0/a_1,1/b_1,2/a_2,3/b_2,4/a_3,9/a_{n+1},10/b_{n+1}}{\node[AnnotationBlue,above,font=\tiny]at(\x,.1){$\t$};}\node[NoteRed,above]at(5.5,0){$\cdots$};
}
\handline{72}{657}{820}{不妨设 $a<x_1<x_1+h_1<x_2<x_2+h_2<\cdots<x_n<x_n+h_n<b$.}
\handline{72}{688}{805}{重新标记，令 $a_1=a$，$b_1=x_1$，$a_2=x_1+h_1$，$b_2=x_2$，$\ldots$，$a_{n+1}=x_n+h_n$，$b_{n+1}=b$.}
\handline{72}{720}{800}{于是 $\sum_{i=1}^{n+1}(b_i-a_i)=b-a-\sum_{j=1}^n h_j<(b-a)-(b-a-\delta)=\delta$.}
\handline{72}{754}{805}{$\therefore|f(b)-f(a)|=|f(b_{n+1})-f(a_1)|\le\sum_{i=1}^{n+1}|f(b_i)-f(a_i)|<\varepsilon$.}
\handline{72}{784}{800}{令 $\varepsilon\to0+\Rightarrow f(b)=f(a)$.}
\handline[Teal]{72}{817}{790}{显然 Cantor 函数 $C(x)$ 是奇异的，但 $C(x)$ 不为常数.}
\handline[Teal]{72}{849}{795}{$\therefore C(x)$ 不是绝对连续的.}
\end{handpage}

% SOURCE PAGE 074
\begin{handpage}{74}
\handline{72}{164}{830}{Thm. $f\in\mathrm{AC}[a,b]\Leftrightarrow f'$ 几乎处处存在，$f'\in L^1[a,b]$ 且 $f(x)-f(a)=\int_a^x f'(t)\,dt$.}
\handline{72}{190}{815}{Pf. ``$\Leftarrow$'' $f'\in L^1[a,b]\Rightarrow\int_{[a,x]}f'$ 作为集函数是绝对连续的.}
\handline{72}{217}{730}{而 $f(x)=f(a)+\int_{[a,x]}f'\Rightarrow f\in\mathrm{AC}[a,b]$. \redann{（一开始的例子中已证）}}
\handline{72}{243}{765}{``$\Rightarrow$'' $f\in\mathrm{AC}[a,b]\Rightarrow f\in\mathrm{BV}[a,b]\Rightarrow f'$ 几乎处处存在且 $f'\in L^1[a,b]$.}
\handline{72}{270}{745}{由于 $f'\in L^1[a,b]$，$\therefore G(x)=\int_{[a,x]}f'\in\mathrm{AC}[a,b]$.}
\handline{72}{296}{740}{$\therefore f(x)-G(x)\in\mathrm{AC}[a,b]$，$\Rightarrow(f(x)-G(x))'$ 几乎处处存在.}
\handline{72}{323}{750}{$\therefore f'(x)-G'(x)=f'(x)-f'(x)=0$ a.e. \redann{（用正则收缩版的 Lebesgue 微分定理）}}
\handline{72}{350}{740}{$\therefore f(x)-G(x)$ 既绝对连续又奇异.}
\handline{72}{376}{740}{$\therefore f(x)-G(x)\equiv f(a)-G(a)=f(a)$}
\handline{72}{404}{730}{$\therefore f(x)-f(a)=G(x)=\int_a^x f'(t)\,dt$.}
\handline[Teal]{72}{433}{770}{\circled{1} Cantor 函数 $C(x)$，$C(x)-C(0)\ne\int_0^x C'(t)\,dt$，从而 $C(x)\notin\mathrm{AC}[0,1]$.}
\handline[Teal]{72}{461}{770}{\circled{2} 若 $f\in\mathrm{BV}[a,b]$，则 $f'$ 几乎处处存在且 $f'\in L^1[a,b]$.}
\handline[Teal]{72}{488}{745}{令 $g(x)=\int_a^x f'(t)\,dt$ 则 $g\in\mathrm{AC}[a,b]$.}
\handline[Teal]{72}{516}{745}{令 $h(x)=f(x)-g(x)$，则 $h'(x)=f'(x)-f'(x)=0$ a.e.，$\therefore h$ 奇异.}
\handline[Teal]{72}{544}{745}{于是 $f(x)=g(x)+h(x)$.}
\handline[NoteRed]{72}{574}{750}{BV 函数可分解为一个 AC 函数与一个奇异函数之和，且在差一个常数的意义下唯一.}
\handline{72}{605}{856}{Thm. $f\in\mathrm{AC}[a,b]$，则 $V(x),P(x),N(x)\in\mathrm{AC}[a,b]$.}
\handline{72}{635}{785}{且 $V(x)=\int_a^x|f'|$，$P(x)=\int_a^x\{f'\}^+$，$N(x)=\int_a^x\{f'\}^-$.}
\handline{72}{664}{856}{Pf. $\forall[\alpha,\beta]\subset[a,b]$，$\Gamma=\{\alpha=t_0<t_1<\cdots<t_n=\beta\}$，$V(\beta)-V(\alpha)=\sup_\Gamma\{\sum_{j=1}^n|f(t_j)-f(t_{j-1})|\}$.}
\handline{72}{693}{800}{$\overset{\ann{\mathrm{A.C.}}}{=}\sup_\Gamma\{\sum_{j=1}^n|\int_{t_{j-1}}^{t_j}f'(t)\,dt|\}\le\sup_\Gamma\{\sum_{j=1}^n\int_{t_{j-1}}^{t_j}|f'(t)|\,dt\}=\int_\alpha^\beta|f'(t)|\,dt$.}
\handline{72}{724}{795}{$\forall\varepsilon>0$，$\{[\alpha_j,\beta_j]\}_{j=1}^\infty$ 两两不交，$\exists\delta>0$，s.t. $\sum(\beta_j-\alpha_j)<\delta$ 时}
\handline{72}{755}{775}{$\sum_{j=1}^\infty V(\beta_j)-V(\alpha_j)\le\sum_{j=1}^\infty\int_{\alpha_j}^{\beta_j}|f'(t)|\,dt=\int_{\bigcup_{j=1}^\infty[\alpha_j,\beta_j]}|f'(t)|\,dt<\varepsilon$.}
\handline[AnnotationBlue]{72}{783}{810}{（这里用的是 $|f'|\in L[a,b]$，其不定积分作为集函数是绝对连续的）}
\handline{72}{813}{810}{$\therefore V(x)\in\mathrm{AC}[a,b]$，$\therefore V(x)-V(a)=\int_a^x V'(t)\,dt=\int_a^x|f'(t)|\,dt$.}
\handline{72}{844}{800}{而 $V(a)=0$，$\therefore V(x)=\int_a^x|f'|$.}
\handline{72}{884}{820}{而 $P(x)=\frac12(V(x)+f(x)-f(a))\Rightarrow P(x)\in\mathrm{AC}[a,b]$，$P(x)=\frac12(\int_a^x|f'|+\int_a^x f')=\int_a^x\{f'\}^+$.}
\handline{72}{923}{800}{$N(x)=\frac12(V(x)-f(x)+f(a))\Rightarrow N(x)\in\mathrm{AC}[a,b]$，$N(x)=\frac12(\int_a^x|f'|-\int_a^x f')=\int_a^x\{f'\}^-$.}
\end{handpage}

% SOURCE PAGE 075
\begin{handpage}{75}
\handline{72}{128}{856}{Thm. \circled{1} $g\in C[a,b]$，$f\in\mathrm{AC}[a,b]$.}
\handline{72}{153}{790}{则 $\int_a^b g\,df=\int_a^b gf'$ \ann{（左边是 R-S 积分，右边是 Lebesgue 积分）}}
\handline{72}{183}{810}{\circled{2} $g,f\in\mathrm{AC}[a,b]$}
\handline{72}{211}{795}{则 $\int_a^b gf'=g(b)f(b)-g(a)f(a)-\int_a^b fg'$ \ann{（Lebesgue 积分）}}
\handline{72}{243}{856}{Pf. \circled{1} $g\in C[a,b]$，$f\in\mathrm{AC}[a,b]\Rightarrow f\in\mathrm{BV}[a,b]$，$\therefore\int_a^b g\,df$ 存在 \ann{（R-S 积分）}}
\handline{72}{270}{765}{$f\in\mathrm{AC}[a,b]\Rightarrow f'\in L^1[a,b]\Rightarrow g\cdot f'$ 可测 $\Rightarrow\int_a^b gf'$ 存在 \ann{（Lebesgue 积分）}}
\handline{72}{296}{790}{$\Gamma=\{a=x_0<x_1<\cdots<x_n=b\}$}
\handline{72}{327}{790}{$\int_a^b gf'=\sum_{j=1}^n\int_{x_{j-1}}^{x_j}gf'=\sum_{j=1}^n\int_{x_{j-1}}^{x_j}(g(x)-g(x_{j-1})+g(x_{j-1}))f'$}
\handline{72}{358}{790}{$=\underbrace{\sum_{j=1}^n g(x_{j-1})\int_{x_{j-1}}^{x_j}f'}_{\ann{\mathrm I}}+\underbrace{\sum_{j=1}^n\int_{x_{j-1}}^{x_j}(g(x)-g(x_{j-1}))f'}_{\ann{\mathrm {II}}}$}
\handline{72}{409}{780}{$g\in C[a,b]\Rightarrow g$ 在 $[a,b]$ 上一致连续.}
\handline{72}{442}{785}{$\forall\varepsilon>0$，$\exists\delta>0$，s.t. $|x-y|<\delta$，$|g(x)-g(y)|<\varepsilon$.}
\handline{72}{477}{780}{当 $\|\Gamma\|<\delta$ 时，$|\mathrm{II}|\le\sum_{j=1}^n\sup_{[x_{j-1},x_j]}|g(x)-g(x_{j-1})|\int_{x_{j-1}}^{x_j}|f'|\le\varepsilon\int_a^b|f'|=\varepsilon\|f'\|_1$.}
\handline{72}{509}{780}{$\therefore\mathrm{II}\to0$（$\|\Gamma\|\to0$）.}
\handline{72}{539}{795}{而 $\mathrm I=\sum_{j=1}^n g(x_{j-1})(f(x_j)-f(x_{j-1}))\to\int_a^b g\,df$（$\|\Gamma\|\to0$）.}
\handline{72}{569}{780}{$\therefore\int_a^b g\,df=\int_a^b gf'$.}
\handline{72}{603}{820}{\circled{2} $g,f\in\mathrm{AC}[a,b]\Rightarrow g,f\in\mathrm{BV}[a,b]$，且 $g,f\in C[a,b]$.}
\handline{72}{634}{790}{$\therefore\int_a^b g\,df=g(b)f(b)-g(a)f(a)+\int_a^b f\,dg$.}
\handline{72}{665}{790}{$\therefore\int_a^b gf'=g(b)f(b)-g(a)f(a)+\int_a^b fg'$.}
\handline[Teal]{72}{698}{810}{在高等微积分中，$g\in C[a,b]$，$f\in C^1[a,b]$ 时}
\handline[Teal]{72}{731}{780}{$\int_a^b g(x)\,df=\int_a^b g(x)f'(x)\,dx$，才会成立.}
\handline[Teal]{72}{763}{805}{但这里 Riemann-Stieltjes 意义下的积分，}
\handline[Teal]{72}{796}{805}{这个定理在 Lebesgue 意义下考虑积分，将 $C^1[a,b]$ 减弱为 $\mathrm{AC}[a,b]$.}
\handline[Teal]{72}{830}{810}{而对于 \circled{2}，分部积分公式对 BV 函数的 R-S 积分成立.}
\handline[Teal]{72}{865}{800}{利用 \circled{1} 可以证明分部积分公式对 AC 函数的 Lebesgue 积分也成立.}
\end{handpage}

% SOURCE PAGE 076
% SECTION: 7.6 | Convex Functions
\begin{handpage}{76}
\handtitle{72}{111}{825}{\S6.6 Convex Functions}
\handline{72}{141}{835}{Def. $\varphi:(a,b)\to\R$}
\handline{72}{165}{775}{若 $\forall[x_1,x_2]\subset(a,b)$，$\forall x\in[x_1,x_2]$，$x=\theta x_1+(1-\theta)x_2$，$0\le\theta\le1$}
\handline{72}{192}{770}{令 $L(x)=\theta\varphi(x_1)+(1-\theta)\varphi(x_2)$ \ann{称为割线（segment line）}}
\handline{72}{218}{775}{有 $\varphi(x)=\varphi(\theta x_1+(1-\theta)x_2)\le L(x)=\theta\varphi(x_1)+(1-\theta)\varphi(x_2)$.}
\handline{72}{245}{780}{则称 $\varphi$ 是 $(a,b)$ 上的凸函数.}
\handline{72}{274}{775}{$\Leftrightarrow\forall a<x_1<x_2<b$，$P_1\ge0$，$P_2\ge0$，$P_1+P_2>0$.}
\handline{72}{305}{715}{$\varphi\bigl(\frac{P_1x_1+P_2x_2}{P_1+P_2}\bigr)\le\frac{P_1\varphi(x_1)+P_2\varphi(x_2)}{P_1+P_2}$.}
\handline{72}{337}{850}{Thm. \ann{(Jensen's Inequality)}}
\handline{72}{365}{825}{$\varphi$ 在 $(a,b)$ 为凸函数，$\{x_j\}_{j=1}^n\subseteq(a,b)$，$\{P_j\}_{j=1}^n$，$P_j\ge0$，$\sum_{j=1}^nP_j>0$.}
\handline{72}{401}{820}{则 $\varphi\biggl(\frac{\sum_{j=1}^nP_jx_j}{\sum_{j=1}^nP_j}\biggr)\le\frac{\sum_{j=1}^nP_j\varphi(x_j)}{\sum_{j=1}^nP_j}$.}
\handline{72}{457}{855}{Prop. \circled{1} $\varphi_1,\varphi_2$ 在 $(a,b)$ 为凸函数，则 $\varphi_1+\varphi_2$ 在 $(a,b)$ 为凸函数.}
\handline{72}{484}{775}{\circled{2} $\varphi$ 在 $(a,b)$ 凸，$C>0$，则 $C\varphi$ 在 $(a,b)$ 凸.}
\handline{72}{512}{775}{\circled{3} $\varphi_k$ 在 $(a,b)$ 凸，$\varphi_k\to\varphi$（逐点）则 $\varphi$ 在 $(a,b)$ 凸.}
\handline{72}{541}{850}{Thm. 若 $\varphi'(x)$ 在 $(a,b)$ 上存在且单增，则 $\varphi$ 在 $(a,b)$ 为凸函数.}
\handline{72}{569}{775}{特别地若 $\varphi''(x)$ 在 $(a,b)$ 上存在且非负，则 $\varphi$ 在 $(a,b)$ 为凸函数.}
\handline{72}{599}{855}{Pf. 用到一个不等式：$b_1,b_2>0$，则 $\min\{\frac{a_1}{b_1},\frac{a_2}{b_2}\}\le\frac{a_1+a_2}{b_1+b_2}\le\max\{\frac{a_1}{b_1},\frac{a_2}{b_2}\}$.}
\handline{72}{634}{770}{$\forall[x_1,x_2]\subset(a,b)$，$x\in[x_1,x_2]$，$L(x)\ge\varphi(x)\Leftrightarrow\frac{\varphi(x_2)-\varphi(x_1)}{x_2-x_1}\ge\frac{\varphi(x)-\varphi(x_1)}{x-x_1}$.}
\handline{72}{676}{770}{而 $\frac{\varphi(x_2)-\varphi(x_1)}{x_2-x_1}=\frac{\varphi(x_2)-\varphi(x)+\varphi(x)-\varphi(x_1)}{(x_2-x)+(x-x_1)}\ge\min\{\frac{\varphi(x_2)-\varphi(x)}{x_2-x},\frac{\varphi(x)-\varphi(x_1)}{x-x_1}\}$.}
\handline{72}{731}{765}{由 Lagrange 中值定理 $\exists\xi_1\in(x_1,x)$，$\xi_2\in(x,x_2)$.}
\handline{72}{766}{760}{s.t. $\frac{\varphi(x_2)-\varphi(x)}{x_2-x}=\varphi'(\xi_2)\ge\varphi'(\xi_1)=\frac{\varphi(x)-\varphi(x_1)}{x-x_1}$.}
\handline{72}{818}{750}{$\therefore\frac{\varphi(x_2)-\varphi(x_1)}{x_2-x_1}\ge\min\{\frac{\varphi(x_2)-\varphi(x)}{x_2-x},\frac{\varphi(x)-\varphi(x_1)}{x-x_1}\}=\frac{\varphi(x)-\varphi(x_1)}{x-x_1}$.}
\handline{72}{864}{856}{Ex. \circled{1} $x^p$，$p\ge1$ 或 $p\le0$ 在 $(0,+\infty)$ 为凸函数.}
\handline{72}{896}{790}{\circled{2} $e^x$ 在 $\R$ 为凸函数.}
\handline{72}{928}{790}{\circled{3} $\log\frac1x=-\log x$ 在 $(0,+\infty)$ 为凸函数.}
\end{handpage}

% SOURCE PAGE 077
\begin{handpage}{77}
\handline{72}{116}{856}{Thm. 若 $\varphi$ 是 $(a,b)$ 上的凸函数，则 $\varphi$ 在 $(a,b)$ 上连续.}
\handline{72}{144}{820}{且 $\varphi'$ 在一个至多可数集以外存在，且 $\varphi'$ 单增.}
\handline{72}{173}{856}{Pf. $\forall x\in(a,b)$，由于 $\varphi$ 是 $(a,b)$ 上的凸函数.}
\handline{72}{201}{810}{$\therefore y\in(x,b)$，当 $y\to x+$ 时，$\frac{\varphi(y)-\varphi(x)}{y-x}$ 递减且 $<\frac{\varphi(b^-)-\varphi(x)}{b-x}$.}
\handline{72}{231}{775}{$y\in(a,x)$，当 $y\to x-$ 时，$\frac{\varphi(y)-\varphi(x)}{y-x}$ 递增且 $>\frac{\varphi(a^+)-\varphi(x)}{a-x}$.}
\handline{72}{261}{795}{$\therefore D^+\varphi(x),D^-\varphi(x)$ 存在 \ann{（左、右导数）}}
\handline{72}{289}{755}{且 $-\infty<D^-\varphi(x)\le D^+\varphi(x)<+\infty$.}
\handline{72}{318}{795}{左右导数存在 $\Rightarrow\varphi$ 在 $x$ 处左右连续 $\Rightarrow\varphi$ 在 $x$ 处连续，$\therefore\varphi$ 在 $(a,b)$ 上连续.}
\handline{72}{349}{810}{由凸性，$\forall a<x_1<x_2<b$，$D^-\varphi(x_1)\le D^+\varphi(x_1)\le D^-\varphi(x_2)\le D^+\varphi(x_2)$.}
\handline{72}{379}{795}{$\therefore D^-\varphi(x),D^+\varphi(x)$ 为单增函数.}
\handline{72}{408}{785}{而单调函数的间断点至多可数.}
\handline{72}{437}{785}{若 $D^+\varphi(x)$ 在 $x_0\in(a,b)$ 处连续.}
\handline{72}{468}{775}{若 $x<x_0$，则 $D^+\varphi(x)\le D^-\varphi(x_0)\le D^+\varphi(x_0)$.}
\handline{72}{500}{755}{令 $x\to x_0-$，则由连续 $\Rightarrow D^+\varphi(x)\to D^+\varphi(x_0)$.}
\handline{72}{532}{765}{$\Rightarrow D^-\varphi(x_0)=D^+\varphi(x_0)$，$\therefore\varphi(x)$ 在 $x_0$ 处可导.}
\handline{72}{563}{765}{$\therefore\varphi'(x)$ 在一个至多可数集以外存在.}
\handline{72}{594}{735}{且 $\varphi'(x)=D^+\varphi(x)$，单增.}
\handline[Teal]{72}{633}{845}{\circled{1} $\varphi$ 是 $(a,b)$ 上的凸函数}
\handline[Teal]{72}{666}{845}{$\Leftrightarrow\forall[x_1,x_2]\subset(a,b)$，$\forall x\in[x_1,x_2]$，$L(x)\ge\varphi(x)$.}
\handline[Teal]{72}{702}{845}{$\Leftrightarrow\forall[x_1,x_2]\subset(a,b)$，$\forall x\in(x_1,x_2)$，$\frac{\varphi(x_2)-\varphi(x_1)}{x_2-x_1}\ge\frac{\varphi(x)-\varphi(x_1)}{x-x_1}$.}
\handline[Teal]{72}{760}{845}{$\Leftrightarrow\forall[x_1,x_2]\subset(a,b)$，$P_1,P_2\ge0$，$P_1+P_2>0$.}
\handline[Teal]{72}{798}{750}{$\varphi\bigl(\frac{P_1x_1+P_2x_2}{P_1+P_2}\bigr)\le\frac{P_1\varphi(x_1)+P_2\varphi(x_2)}{P_1+P_2}$.}
\handline[Teal]{72}{841}{850}{$\Leftrightarrow\forall x\in(a,b)$，$\widetilde\varphi(y)=\frac{\varphi(y)-\varphi(x)}{y-x}$，$\widetilde\varphi$ 在 $(a,x)$ 和 $(x,b)$ 上单调增加.}
\end{handpage}

% SOURCE PAGE 078
\begin{handpage}{78}
\handline[Teal]{72}{130}{523}{\circled{2}\ 几何视角，先固定$3$个点：$(x,\varphi(x))$，$(y_1,\varphi(y_1))$，$(y_2,\varphi(y_2))$，$a<y_1<x<y_2<b$.}
\handline[Teal]{72}{155}{520}{由凸性$\Rightarrow$若$y\in(x,y_2)$，则$(y,\varphi(y))$在割线$L_2$之下}
\handline[Teal]{72}{182}{515}{事实上还要在割线$L_1$之上}
\handline[Teal]{72}{209}{515}{若$(y,\varphi(y))$在$L_1$之下，再作割线$L_3$.}
\handline[Teal]{72}{236}{515}{由凸性$x\in(y_1,y)$，则$(x,\varphi(x))$在$L_3$之下，矛盾!}
\handline[Teal]{72}{263}{510}{于是$(y,\varphi(y))$必须在$L_1$和$L_2$之间.}
\handline[Teal]{72}{291}{500}{$\therefore\quad L_1(y)\le\varphi(y)\le L_2(y)$.}
\handdiagram{613}{135}{315}{163}{
 \draw[Teal] (.1,.2)--(5.8,.2);
 \draw[Teal,thick] plot[smooth] coordinates{(.8,2.7)(1.3,2.0)(2,1.7)(2.6,1.8)(3.1,2.25)(3.6,3.2)(4.2,4.6)(4.8,6.1)};
 \draw[NoteRed] (.1,2.7)--(6.5,.9); \node[NoteRed,above] at(.55,3.1){$L_1$}; \draw[NoteRed,->](.55,2.75)--(.55,3.05);
 \draw[AnnotationBlue] (2,-.3)--(4.9,6.1); \node[AnnotationBlue,right]at(4.9,5.65){$L_2$};
 \draw[Teal] (.55,2.75)--(6.4,-.4); \node[Teal,right]at(6.0,-.5){$L_3$};
 \draw[Teal,dashed] (1.1,.2)--(1.1,2.4); \draw[Teal,dashed](2.9,.2)--(2.9,1.95); \draw[Teal,dashed](3.65,.2)--(3.65,.9); \draw[Teal,dashed](4.4,.2)--(4.4,4.85);
 \node[Teal,below]at(1.1,.2){$y_1$}; \node[Teal,below]at(2.9,.2){$x$}; \node[Teal,below]at(3.65,.2){$y$}; \node[Teal,below]at(4.4,.2){$y_2$};
}
\handline[Teal]{72}{322}{770}{而$L_1(y),L_2(y)$相交于$(x,\varphi(x))$，$\therefore y\to x+$时$L_1(y),L_2(y)\to\varphi(x)$.}
\handline[Teal]{72}{348}{735}{由夹逼定理$\Rightarrow\varphi(y)\to\varphi(x)$，从而$\varphi$必须在$x$处连续.}
\handline{72}{378}{825}{\textbf{Thm.}\quad 若$\varphi$在$(a,b)$上为凸函数，则$\varphi$在$(a,b)$上是紧Lipschitz的.}
\handline{72}{406}{745}{即$\forall[x_1,x_2]\subset(a,b)$，$\varphi$在$[x_1,x_2]$上Lipschitz连续.}
\handline{72}{435}{755}{特别地，$\varphi$在$[x_1,x_2]$上AC. $\Rightarrow\varphi(x_2)-\varphi(x_1)=\int_{x_1}^{x_2}\varphi'(t)\,dt$.}
\handline{72}{468}{830}{\textbf{Pf.}\quad $\forall[x_1,x_2]\subset(a,b)$，$x<y\in[x_1,x_2]$}
\handline{72}{500}{720}{$D^+\varphi(x_1)\le D^+\varphi(x)\le\frac{\varphi(y)-\varphi(x)}{y-x}\le D^-\varphi(y)\le D^-\varphi(x_2)$.}
\handline{72}{546}{730}{$\Rightarrow\quad\left|\frac{\varphi(y)-\varphi(x)}{y-x}\right|\le\max\{|D^+\varphi(x_1)|,|D^-\varphi(x_2)|\}=L>0$.}
\handline{72}{587}{695}{$\Rightarrow |\varphi(y)-\varphi(x)|\le L|y-x|$，$\therefore\varphi$在$[x_1,x_2]$上Lipschitz连续.}
\handline{72}{619}{835}{\textbf{Def.}\quad 若$\varphi$在$(a,b)$上为凸函数，则$\forall x_0\in(a,b)$，}
\handline{72}{648}{740}{存在一条经过$(x_0,\varphi(x_0))$的直线$L(x)$，s.t. $\varphi(x)\ge L(x)\quad\forall x\in(a,b)$.}
\handline{72}{682}{730}{且斜率$m$满足\quad$D^-\varphi(x_0)\le m\le D^+\varphi(x_0)$}
\handline{72}{711}{730}{这条线称为支撑线\ \ann{(Supporting line)}}
\handline[Teal]{72}{745}{755}{若$\varphi'(x_0)$存在，此直线唯一，即为切线$L(x)=\varphi'(x_0)(x-x_0)+\varphi(x_0)$.}
\handline[Teal]{72}{777}{735}{$\forall x\in(a,b)\quad\varphi(x)\ge\varphi'(x_0)(x-x_0)+\varphi(x_0)$.}
\handline[Teal]{72}{811}{755}{若$\varphi'(x_0)$不存在，则直线可能不唯一. 比如：$\varphi(x)=|x|$，$x_0=0$.}
\handline[Teal]{72}{842}{470}{在$0$处有无数条支撑线}
\handdiagram{688}{825}{240}{78}{
 \draw[Teal,thick] (-2,1.9)--(0,0)--(2,1.9);
 \draw[NoteRed] (-3,.4)--(2.4,-.4); \node[NoteRed,right]at(2.4,-.4){$L_1$};
 \draw[NoteRed] (-3,-.6)--(2.4,.6); \node[NoteRed,right]at(2.4,.6){$L_2$};
 \node[NoteRed,below]at(0,-.4){$\cdots$};
}
\end{handpage}

% SOURCE PAGE 079
\begin{handpage}{79}
\handline{72}{142}{856}{\textbf{Thm.}\ \ann{(Jensen Integral Inequality)}}
\handline{72}{170}{820}{$f,p$在$E\subseteq\R^n$上可测，$E\in\mathcal M$，且$f,p$几乎处处有限，$fp,p\in L^1(E)$.}
\handline{72}{196}{815}{$p\ge0$，$\int_Ep>0$，$\varphi$在包含$f$值域的开区间上为凸函数.}
\handline{72}{231}{810}{则\quad$\varphi\left(\frac{\int_Efp}{\int_Ep}\right)\le\frac{\int_E\varphi(f)p}{\int_Ep}$.}
\handline{72}{279}{856}{\textbf{Pf.}\quad $f$几乎处处有限，$\therefore\exists a,b$ s.t. $-\infty\le a<b\le+\infty$.}
\handline{72}{309}{800}{且$\varphi$在$(a,b)$上为凸函数，$a<f(x)<b$ a.e.}
\handline{72}{343}{810}{$\Rightarrow\quad a<\frac{\int_Efp}{\int_Ep}<b$，设$ x_0=\frac{\int_Efp}{\int_Ep}$.}
\handline{72}{390}{800}{则在$x_0$处有一条支撑线，$L(x)=m(x-x_0)+\varphi(x_0)$.}
\handline{72}{423}{790}{$m(x-x_0)+\varphi(x_0)\le\varphi(x)\quad\forall x\in(a,b)$.}
\handline{72}{454}{795}{令$x=f(t)$，则$m(f(t)-x_0)+\varphi(x_0)\le\varphi(f(t))$.}
\handline{72}{486}{770}{$\therefore\quad m(f(t)-x_0)p(t)+\varphi(x_0)p(t)\le\varphi(f(t))p(t)$}
\handline{72}{518}{775}{$\therefore\quad\int_E m(f(t)-x_0)p(t)\,dt+\int_E\varphi(x_0)p(t)\,dt\le\int_E\varphi(f(t))p(t)\,dt$.}
\handline{72}{550}{780}{而$\int_E m(f(t)-x_0)p(t)\,dt=m\left(\int_E f(t)p(t)\,dt-\int_E x_0p(t)\,dt\right)$}
\handline{72}{585}{710}{$=m\left(\int_E f(t)p(t)\,dt-\frac{\int_E f(t)p(t)\,dt}{\int_E p(t)\,dt}\int_E p(t)\,dt\right)=0$.}
\handline{72}{644}{780}{$\therefore\quad\int_E\varphi(x_0)p(t)\,dt\le\int_E\varphi(f(t))p(t)\,dt$}
\handline{72}{683}{700}{即\quad$\varphi\left(\frac{\int_Efp}{\int_Ep}\right)\le\frac{\int_E\varphi(f)p}{\int_Ep}$.}
\handline[Teal]{72}{749}{855}{注意，$ dm=\frac{p\,dx}{\int_Ep}$可以视为概率测度，则$\varphi\left(\int_E f\,dm\right)\le\int_E\varphi(f)\,dm$.}
\handline[Teal]{72}{790}{770}{形式更加简洁，且在概率论中也有广泛应用.}
\end{handpage}

\end{document}